% !TeX program = xelatex
% !TeX encoding = UTF-8
\documentclass[UTF8,10pt,a4paper,fontset=fandol]{ctexart}
\usepackage{amsmath,amssymb,graphicx,tabularx,array,multirow,enumitem,titlesec,xcolor,ragged2e,etoolbox}
\usepackage[a4paper,top=1.6cm,bottom=1.6cm,left=1.6cm,right=1.6cm]{geometry}
\setlength{\parindent}{2em}
\setlength{\parskip}{0.15em}
\setlength{\tabcolsep}{4pt}
\renewcommand{\arraystretch}{1.15}
\setlist[enumerate]{leftmargin=2.6em,itemsep=0.1em,topsep=0.2em}
\titleformat{\section}{\large\bfseries}{\thesection}{0.6em}{}
\titleformat{\subsection}{\normalsize\bfseries}{\thesubsection}{0.6em}{}
\titleformat{\subsubsection}{\normalsize\bfseries}{\thesubsubsection}{0.6em}{}
\titlespacing*{\section}{0pt}{0.7em}{0.25em}
\titlespacing*{\subsection}{0pt}{0.55em}{0.2em}
\titlespacing*{\subsubsection}{0pt}{0.4em}{0.15em}
\raggedbottom
\BeforeBeginEnvironment{tabularx}{\par\noindent}
\newcommand{\fieldvalue}[1]{#1}
\newcommand{\handwritten}[1]{#1}
\begin{document}

\noindent
\begin{minipage}{0.28\textwidth}
\small
\textbf{铂姆卓越生物科技（北京）有限公司}\\[-0.2em]
\scriptsize PLATINUM LIFE EXCELLENCE BIOTECH
\end{minipage}
\hfill
\begin{minipage}{0.68\textwidth}
\raggedleft
文件编码：REC-YZ-SMP-QC-01-001-d\qquad 版本号：1.4
\end{minipage}

\vspace{0.3em}

\begin{center}
\textbf{\large 检验申请单}
\end{center}

\small
\begin{tabularx}{\textwidth}{|>{\centering\arraybackslash}p{2.2cm}|X|>{\centering\arraybackslash}p{2.2cm}|X|}
\hline
样品名称 &
% #VALUE_ID: LEX-P0001-V0001
% #FIELD_VALUE: 样品名称
% #TODO #HANDWRITTEN: 二氧化硅……；手写内容辨识不清
\fieldvalue{\handwritten{二氧化硅……}} &
样品代码 &
% #VALUE_ID: LEX-P0001-V0002
% #FIELD_VALUE: 样品代码
% #TODO #HANDWRITTEN: 25/0402；手写字符部分不清晰
\fieldvalue{\handwritten{25/0402}}
\\ \hline

样品批号 &
% #VALUE_ID: LEX-P0001-V0003
% #FIELD_VALUE: 样品批号
% #TODO #HANDWRITTEN: A1……；手写内容辨识不清
\fieldvalue{\handwritten{A1……}} &
样品规格/数量 &
% #VALUE_ID: LEX-P0001-V0004
% #FIELD_VALUE: 样品规格/数量
% #TODO #HANDWRITTEN: 约30g/1袋；手写内容部分不清晰
\fieldvalue{\handwritten{约30g/1袋}}
\\ \hline

样品类别 &
% #VALUE_ID: LEX-P0001-V0005
% #FIELD_VALUE: 样品类别（选项：生产过程检验样品）
% #HANDWRITTEN: ✓
\fieldvalue{\handwritten{\checkmark}}生产过程检验样品
\quad $\square$包装产品
\quad $\square$成品
\newline
\hspace*{2.2cm}$\square$其他（\hspace{2cm}） &
\multicolumn{2}{l|}{}
\\ \hline

检验类型 &
% #VALUE_ID: LEX-P0001-V0006
% #FIELD_VALUE: 检验类型（选项：初验）
% #HANDWRITTEN: ✓
\fieldvalue{\handwritten{\checkmark}}初验
\quad $\square$复验
\quad $\square$退货
\quad $\square$稳定性
\quad $\square$留样
&
\multicolumn{2}{l|}{}
\\ \hline

请验目的 &
% #VALUE_ID: LEX-P0001-V0007
% #FIELD_VALUE: 请验目的（选项：全检）
% #HANDWRITTEN: ✓
\fieldvalue{\handwritten{\checkmark}}全检
\quad $\square$留样
\quad $\square$其他（\hspace{2cm}）
&
\multicolumn{2}{l|}{}
\\ \hline

请验部门 &
% #VALUE_ID: LEX-P0001-V0008
% #FIELD_VALUE: 请验部门（选项：生产部）
% #HANDWRITTEN: ✓
\fieldvalue{\handwritten{\checkmark}}生产部
\quad $\square$物料管理部
\quad $\square$其他（\hspace{2cm}）
&
\multicolumn{2}{l|}{}
\\ \hline

样品情况 &
$\square$随请验单一起提供
\quad $\square$未随请验单一起提供
&
\multicolumn{2}{l|}{}
\\ \hline

请验人 &
% #VALUE_ID: LEX-P0001-V0009
% #FIELD_VALUE: 请验人
% #TODO #HANDWRITTEN: 俞学宝；手写姓名辨识不完全确定
\fieldvalue{\handwritten{俞学宝}}
&
请验日期 &
% #VALUE_ID: LEX-P0001-V0010
% #FIELD_VALUE: 请验日期
% #HANDWRITTEN: 2025.08.20
\fieldvalue{\handwritten{2025.08.20}}
\\ \hline

收验人 &
% #VALUE_ID: LEX-P0001-V0011
% #FIELD_VALUE: 收验人
% #TODO #HANDWRITTEN: 俞学宝；手写姓名辨识不完全确定
\fieldvalue{\handwritten{俞学宝}}
&
收验日期 &
% #VALUE_ID: LEX-P0001-V0012
% #FIELD_VALUE: 收验日期
% #HANDWRITTEN: 2025.08.20
\fieldvalue{\handwritten{2025.08.20}}
\\ \hline

备注 &
% #VALUE_ID: LEX-P0001-V0013
% #FIELD_VALUE: 备注
% #TODO #HANDWRITTEN: 24sth\quad 8-?；手写内容难以辨识
\fieldvalue{\handwritten{24sth\quad 8-?}} &
\multicolumn{2}{l|}{}
\\ \hline
\end{tabularx}

% LEXOID_PAGE_COMPLETED: 1/60

\newpage

\begin{center}
\small
\begin{tabularx}{0.92\textwidth}{|p{2.2cm}|X|p{2.2cm}|X|}
\hline
\multicolumn{4}{|c|}{\textbf{取样记录}}\\
\hline
\multicolumn{2}{|l|}{文件编号：REC-YZ-SMP-QC-01-004-b} &
\multicolumn{2}{r|}{版本号：1.4}\\
\hline
名称 &
% #VALUE_ID: LEX-P0002-V0001
% #FIELD_VALUE: 名称
% #TODO #HANDWRITTEN: 二氧化硅凝胶瓶; handwriting partially unclear
\fieldvalue{\handwritten{二氧化硅凝胶瓶}} &
代码 &
% #VALUE_ID: LEX-P0002-V0002
% #FIELD_VALUE: 代码
% #TODO #HANDWRITTEN: Q1402; first character partly unclear
\fieldvalue{\handwritten{Q1402}}\\
\hline
规格 &
% #VALUE_ID: LEX-P0002-V0003
% #FIELD_VALUE: 规格
% #HANDWRITTEN: 约20 ml/袋
\fieldvalue{\handwritten{约20 ml/袋}} &
总数量 &
% #VALUE_ID: LEX-P0002-V0004
% #FIELD_VALUE: 总数量
% #HANDWRITTEN: /
\fieldvalue{\handwritten{/}}\\
\hline
批号 &
% #VALUE_ID: LEX-P0002-V0005
% #FIELD_VALUE: 批号
% #TODO #HANDWRITTEN: A13220m08b09; handwriting unclear
\fieldvalue{\handwritten{A13220m08b09}} &
进厂批号 &
% #VALUE_ID: LEX-P0002-V0006
% #FIELD_VALUE: 进厂批号
% #HANDWRITTEN: /
\fieldvalue{\handwritten{/}}\\
\hline
\multicolumn{1}{|c|}{供货单位/厂家} &
\multicolumn{3}{l|}{
% #VALUE_ID: LEX-P0002-V0007
% #FIELD_VALUE: 供货单位/厂家
% #HANDWRITTEN: /
\fieldvalue{\handwritten{/}}}\\
\hline
\multicolumn{1}{|c|}{取样目的} &
\multicolumn{3}{l|}{
% #VALUE_ID: LEX-P0002-V0008
% #FIELD_VALUE: 取样目的
% #HANDWRITTEN: ✓ 正常检验
\fieldvalue{\handwritten{✓ 正常检验}}
\quad $\square$ 复检
\quad $\square$ 稳定性
\quad $\square$ 项确认检验
\quad $\square$ 留样
\quad $\square$ 其他：\underline{\hspace{2cm}}}\\
\hline
\multicolumn{1}{|c|}{取样地点} &
\multicolumn{3}{l|}{
% #VALUE_ID: LEX-P0002-V0009
% #FIELD_VALUE: 取样地点
% #HANDWRITTEN: ✓ 产品出品
\fieldvalue{\handwritten{✓ 产品出品}}
\quad $\square$ 物料库
\quad $\square$ 细胞库
\quad $\square$ 包装间}\\
\hline
\multicolumn{1}{|c|}{贮存条件} &
\multicolumn{3}{l|}{
% #VALUE_ID: LEX-P0002-V0010
% #FIELD_VALUE: 贮存条件
% #HANDWRITTEN: ✓ 2--8°C
\fieldvalue{\handwritten{✓ 2--8°C}}
\quad $\square$ -20°C
\quad $\square$ -80°C
\quad $\square$ 常温
\quad $\square$ 其他：\underline{\hspace{1.2cm}}}\\
\hline
\multicolumn{1}{|c|}{盛装容器} &
\multicolumn{3}{l|}{
% #VALUE_ID: LEX-P0002-V0011
% #FIELD_VALUE: 盛装容器
% #HANDWRITTEN: ✓ 保温箱
\fieldvalue{\handwritten{✓ 保温箱}}
\quad $\square$ 泡沫箱
\quad $\square$ 手提篮
\quad $\square$ 其他：\underline{\hspace{1.5cm}}}\\
\hline
\multicolumn{1}{|c|}{外观复核} &
\multicolumn{3}{l|}{
% #VALUE_ID: LEX-P0002-V0012
% #FIELD_VALUE: 外观复核
% #HANDWRITTEN: ✓ 合格
\fieldvalue{\handwritten{✓ 合格}}
\quad $\square$ 不合格}\\
\hline
取样量 &
% #VALUE_ID: LEX-P0002-V0013
% #FIELD_VALUE: 取样量
% #HANDWRITTEN: /
\fieldvalue{\handwritten{/}} &
检验量 &
% #VALUE_ID: LEX-P0002-V0014
% #FIELD_VALUE: 检验量
% #HANDWRITTEN: /
\fieldvalue{\handwritten{/}}\\
\hline
留样量 &
% #VALUE_ID: LEX-P0002-V0015
% #FIELD_VALUE: 留样量
% #HANDWRITTEN: /
\fieldvalue{\handwritten{/}} &
\multicolumn{2}{c|}{}\\
\hline
取样人 &
% #VALUE_ID: LEX-P0002-V0016
% #FIELD_VALUE: 取样人
% #HANDWRITTEN: 胡丽娟
\fieldvalue{\handwritten{胡丽娟}} &
取样日期 &
% #VALUE_ID: LEX-P0002-V0017
% #FIELD_VALUE: 取样日期
% #HANDWRITTEN: 2023.02.20
\fieldvalue{\handwritten{2023.02.20}}\\
\hline
复核人 &
% #VALUE_ID: LEX-P0002-V0018
% #FIELD_VALUE: 复核人
% #TODO #HANDWRITTEN: 薛永强; handwriting partly unclear
\fieldvalue{\handwritten{薛永强}} &
复核日期 &
% #VALUE_ID: LEX-P0002-V0019
% #FIELD_VALUE: 复核日期
% #HANDWRITTEN: 2023.08.30
\fieldvalue{\handwritten{2023.08.30}}\\
\hline
备注 &
\multicolumn{3}{l|}{
% #VALUE_ID: LEX-P0002-V0020
% #FIELD_VALUE: 备注
% #TODO #HANDWRITTEN: 24.5L，8.2L; handwriting unclear
\fieldvalue{\handwritten{24.5L，8.2L}}}\\
\hline
\end{tabularx}

\vspace{1cm}
第 1 页 共 1 页
\end{center}
% LEXOID_PAGE_COMPLETED: 2/60

\newpage

\section*{生物反应器细胞液检验操作记录}

\begin{center}
\begin{tabular}{|p{0.18\textwidth}|p{0.32\textwidth}|p{0.18\textwidth}|p{0.22\textwidth}|}
\hline
\multicolumn{4}{|c|}{葡萄糖浓度检验} \\
\hline
样品名称 &
% #VALUE_ID: LEX-P0003-V0001
% #FIELD_VALUE: 样品名称
\fieldvalue{二级生物反应器细胞泥液} &
样品规格 &
% #VALUE_ID: LEX-P0003-V0002
% #FIELD_VALUE: 样品规格
% #HANDWRITTEN: 20
\fieldvalue{\handwritten{20}~ml/袋} \\
\hline
样品批号 &
% #VALUE_ID: LEX-P0003-V0003
% #FIELD_VALUE: 样品批号
% TODO #HANDWRITTEN: A?23?00?; handwriting is partially illegible
\fieldvalue{\handwritten{A?23?00?}} &
取样日期 &
% #VALUE_ID: LEX-P0003-V0004
% #FIELD_VALUE: 取样日期
% #HANDWRITTEN: 202?年08月27日
\fieldvalue{\handwritten{202?年08月27日}} \\
\hline
\end{tabular}
\end{center}

\begin{enumerate}
\item \textbf{仪器设备、试剂耗材：}

\subsection*{1.1 \quad 仪器设备}

\begin{center}
\begin{tabular}{|p{0.17\textwidth}|p{0.25\textwidth}|p{0.17\textwidth}|p{0.20\textwidth}|p{0.15\textwidth}|}
\hline
名称 & 厂家 & 型号 & 设备编号 & 检查确认 \\
\hline
血糖仪 & 三诺生物传感股份有限公司 & 安稳+ & & 
% #VALUE_ID: LEX-P0003-V0005
% #FIELD_VALUE: 血糖仪检查确认
% #HANDWRITTEN: ✓
\fieldvalue{\handwritten{\checkmark}}正常 \quad □异常 \\
\hline
移液器 & Eppendorf Corporate & 0.1--2.5$\mu$l &
% #VALUE_ID: LEX-P0003-V0006
% #FIELD_VALUE: 移液器设备编号
% TODO #HANDWRITTEN: 073408?; final character is unclear
\fieldvalue{\handwritten{073408?}} &
% #VALUE_ID: LEX-P0003-V0007
% #FIELD_VALUE: 移液器检查确认
% #HANDWRITTEN: ✓
\fieldvalue{\handwritten{\checkmark}}正常 \quad □异常 \\
\hline
\end{tabular}
\end{center}

\subsection*{1.2 \quad 试剂耗材}

\begin{center}
\begin{tabular}{|p{0.17\textwidth}|p{0.32\textwidth}|p{0.23\textwidth}|p{0.20\textwidth}|}
\hline
名称 & 厂家 & 批号 & 有效期至 \\
\hline
血糖测试条 & 三诺生物传感股份有限公司 &
% #VALUE_ID: LEX-P0003-V0008
% #FIELD_VALUE: 血糖测试条批号
% #HANDWRITTEN: 44C?44-06
\fieldvalue{\handwritten{44C?44-06}} &
% #VALUE_ID: LEX-P0003-V0009
% #FIELD_VALUE: 血糖测试条有效期至
% #HANDWRITTEN: 2027.01.07
\fieldvalue{\handwritten{2027.01.07}} \\
\hline
\end{tabular}
\end{center}

\item \textbf{操作步骤：}

\begin{center}
\begin{tabular}{|p{0.23\textwidth}|p{0.23\textwidth}|p{0.23\textwidth}|p{0.23\textwidth}|}
\hline
\multicolumn{4}{|c|}{操作过程} \\
\hline
\multicolumn{3}{|p{0.69\textwidth}|}{取出测试条，取样品自然沉降后的上清液0.6$\mu$l加入测试条电极反应室，血糖仪自动发出“嘀”提示音，吸入样品完成，血糖仪自动进入倒计时，开始测试；} &
取上清液：
% #VALUE_ID: LEX-P0003-V0010
% #FIELD_VALUE: 取上清液体积
% #HANDWRITTEN: 0.6
\fieldvalue{\handwritten{0.6}}~$\mu$l \\
\hline
\multicolumn{3}{|p{0.69\textwidth}|}{重复以上步骤2次，共计数3次。} &
每个样品分别共计数
% #VALUE_ID: LEX-P0003-V0011
% #FIELD_VALUE: 样品计数次数
% #HANDWRITTEN: 3
\fieldvalue{\handwritten{3}} 次。 \\
\hline
\multicolumn{4}{|c|}{检测结果} \\
\hline
\multicolumn{4}{|c|}{葡萄糖浓度（mmol/L）} \\
\hline
第一次 & 第二次 & 第三次 & 平均值（保留一位小数） \\
\hline
% #VALUE_ID: LEX-P0003-V0012
% #FIELD_VALUE: 葡萄糖浓度第一次
% #HANDWRITTEN: 21.3
\fieldvalue{\handwritten{21.3}} &
% #VALUE_ID: LEX-P0003-V0013
% #FIELD_VALUE: 葡萄糖浓度第二次
% #HANDWRITTEN: 21.6
\fieldvalue{\handwritten{21.6}} &
% #VALUE_ID: LEX-P0003-V0014
% #FIELD_VALUE: 葡萄糖浓度第三次
% #HANDWRITTEN: 21.1
\fieldvalue{\handwritten{21.1}} &
% #VALUE_ID: LEX-P0003-V0015
% #FIELD_VALUE: 葡萄糖浓度平均值
% #HANDWRITTEN: 21.0
\fieldvalue{\handwritten{21.0}} \\
\hline
\multicolumn{1}{|l|}{备注} &
% #VALUE_ID: LEX-P0003-V0016
% #FIELD_VALUE: 备注
% #HANDWRITTEN: /
\fieldvalue{\handwritten{/}} & \multicolumn{2}{l|}{} \\
\hline
\multicolumn{1}{|l|}{实验人/日期} &
% #VALUE_ID: LEX-P0003-V0017
% #FIELD_VALUE: 实验人
% TODO #HANDWRITTEN: 楠?生; handwriting is unclear
\fieldvalue{\handwritten{楠?生}}\quad
% #VALUE_ID: LEX-P0003-V0018
% #FIELD_VALUE: 实验日期
% #HANDWRITTEN: 2021.08.30
\fieldvalue{\handwritten{2021.08.30}} &
\multicolumn{2}{l|}{} \\
\hline
\multicolumn{1}{|l|}{复核人/日期} &
% #VALUE_ID: LEX-P0003-V0019
% #FIELD_VALUE: 复核人
% TODO #HANDWRITTEN: 董?芳; handwriting is unclear
\fieldvalue{\handwritten{董?芳}}\quad
% #VALUE_ID: LEX-P0003-V0020
% #FIELD_VALUE: 复核日期
% #HANDWRITTEN: 2021.08.30
\fieldvalue{\handwritten{2021.08.30}} &
\multicolumn{2}{l|}{} \\
\hline
\end{tabular}
\end{center}
\end{enumerate}

\begin{center}
第1页 共5页
\end{center}
% LEXOID_PAGE_COMPLETED: 3/60

\newpage

\section*{生物反应器细胞液检验操作记录}

\noindent
\textbf{文件编码：} REC-YZ-SOP-QC-99-006-a

\begin{tabularx}{\textwidth}{|l|X|l|X|}
\hline
\multicolumn{4}{|c|}{盛度与活率检验} \\ \hline
样品名称 &
% #VALUE_ID: LEX-P0004-V0001
% #FIELD_VALUE: 样品名称
\fieldvalue{二级生物反应器细胞液} &
样品规格 &
% #VALUE_ID: LEX-P0004-V0002
% #FIELD_VALUE: 样品规格
% #HANDWRITTEN: 20
\fieldvalue{\handwritten{20}} ml/袋
\\ \hline
样品批号 &
% #VALUE_ID: LEX-P0004-V0003
% #FIELD_VALUE: 样品批号
% #TODO #HANDWRITTEN: A312220?; handwriting partially unclear
\fieldvalue{\handwritten{A312220?}} &
取样日期 &
% #VALUE_ID: LEX-P0004-V0004
% #FIELD_VALUE: 取样日期
% #HANDWRITTEN: 2025年8月30日
\fieldvalue{\handwritten{2025年8月30日}} \\ \hline
\end{tabularx}

\section{仪器设备、试剂耗材}

\subsection{仪器设备}

\begin{tabularx}{\textwidth}{|X|X|X|X|X|X|}
\hline
名称 & 厂家 & 型号 & 设备编号 & 校准有效期至 & 检查确认 \\
\hline
细胞计数仪 & Nexcelom & Cellometer K2 & 1410905 &
% #VALUE_ID: LEX-P0004-V0005
% #FIELD_VALUE: 细胞计数仪校准有效期
% #HANDWRITTEN: 2025.11.10
\fieldvalue{\handwritten{2025.11.10}} &
% #VALUE_ID: LEX-P0004-V0006
% #FIELD_VALUE: 细胞计数仪检查确认
\fieldvalue{正常} \\
\hline
高速冷冻离心机 & Thermo & ST16R & 1410810 &
% #VALUE_ID: LEX-P0004-V0007
% #FIELD_VALUE: 高速冷冻离心机校准有效期
% #HANDWRITTEN: 2025.01.17
\fieldvalue{\handwritten{2025.01.17}} &
% #VALUE_ID: LEX-P0004-V0008
% #FIELD_VALUE: 高速冷冻离心机检查确认
\fieldvalue{正常} \\
\hline
电热恒温水槽 & 上海一恒 & DK-8AXX & 1411001 &
% #VALUE_ID: LEX-P0004-V0009
% #FIELD_VALUE: 电热恒温水槽校准有效期
% #HANDWRITTEN: 2025.05.22
\fieldvalue{\handwritten{2025.05.22}} &
% #VALUE_ID: LEX-P0004-V0010
% #FIELD_VALUE: 电热恒温水槽检查确认
\fieldvalue{正常} \\
\hline
\end{tabularx}

\subsection{试剂耗材}

\begin{tabularx}{\textwidth}{|X|X|X|X|}
\hline
名称 & 厂家 & 批号 & 有效期至 \\ \hline
AO/PI 染液 & Nexcelom &
% #VALUE_ID: LEX-P0004-V0011
% #FIELD_VALUE: AO/PI染液批号
% #HANDWRITTEN: 332309
\fieldvalue{\handwritten{332309}} &
% #VALUE_ID: LEX-P0004-V0012
% #FIELD_VALUE: AO/PI染液有效期
% #HANDWRITTEN: 2027.01.04
\fieldvalue{\handwritten{2027.01.04}} \\
\hline
0.9\%氯化钠注射液 & 石家庄四药 &
% #VALUE_ID: LEX-P0004-V0013
% #FIELD_VALUE: 0.9\%氯化钠注射液批号
% #HANDWRITTEN: 241010322
\fieldvalue{\handwritten{241010322}} &
% #VALUE_ID: LEX-P0004-V0014
% #FIELD_VALUE: 0.9\%氯化钠注射液有效期
% #HANDWRITTEN: 2026.10.06
\fieldvalue{\handwritten{2026.10.06}} \\
\hline
20\%人血白蛋白 & Baxter AG &
% #VALUE_ID: LEX-P0004-V0015
% #FIELD_VALUE: 20\%人血白蛋白批号
% #HANDWRITTEN: AB13103A
\fieldvalue{\handwritten{AB13103A}} &
% #VALUE_ID: LEX-P0004-V0016
% #FIELD_VALUE: 20\%人血白蛋白有效期
% #HANDWRITTEN: 2027.02.28
\fieldvalue{\handwritten{2027.02.28}} \\
\hline
\end{tabularx}

\section{操作步骤}

\subsection*{溶液准备}

\begin{tabularx}{\textwidth}{|X|X|}
\hline
20\%人血白蛋白加入到95 ml &
% #VALUE_ID: LEX-P0004-V0017
% #FIELD_VALUE: 20\%人血白蛋白加入量
% #HANDWRITTEN: 1
\fieldvalue{\handwritten{1}} ml \\ \hline
0.9\%氯化钠注射液，混匀即得 &
% #VALUE_ID: LEX-P0004-V0018
% #FIELD_VALUE: 0.9\%氯化钠注射液加入量
% #HANDWRITTEN: 19
\fieldvalue{\handwritten{19}} ml \\ \hline
\end{tabularx}

\begin{tabularx}{\textwidth}{|X|X|X|}
\hline
名称 & 配制批号 & 有效期至 \\ \hline
5$\times$裂解液 &
% #VALUE_ID: LEX-P0004-V0019
% #FIELD_VALUE: 5倍裂解液配制批号
% #HANDWRITTEN: S20250728-01
\fieldvalue{\handwritten{S20250728-01}} &
% #VALUE_ID: LEX-P0004-V0020
% #FIELD_VALUE: 5倍裂解液有效期
% #HANDWRITTEN: 2025.10.27
\fieldvalue{\handwritten{2025.10.27}} \\
\hline
\end{tabularx}

\subsection*{操作过程}

\noindent
按照50 ml离心管的刻度读取样品体积，往样品中加入对应体积的5$\times$裂解液（每4 ml样品加入1 ml 5$\times$裂解液，按照“只进不舍”原则，如26 ml样品，加入7 ml 5$\times$裂解液），充分混匀后倒入至2.3 ml的离心管中，置于37°C电热恒温水槽中裂解20 min后，取出离心管摇晃混匀，使样品充分裂解；再置于37°C电热恒温水槽中裂解10 min得到细胞悬液。

\begin{tabularx}{\textwidth}{|X|X|}
\hline
\parbox[t]{0.57\textwidth}{\raggedright
取样体积：样品中加入对应体积的5$\times$裂解液，充分混匀后倒入离心管中，置于37°C电热恒温水槽中裂解20 min。}
&
\parbox[t]{0.38\textwidth}{\raggedright
样品体积：
% #VALUE_ID: LEX-P0004-V0021
% #FIELD_VALUE: 样品体积
% #HANDWRITTEN: 37
\fieldvalue{\handwritten{37}} ml

裂解液：
% #VALUE_ID: LEX-P0004-V0022
% #FIELD_VALUE: 裂解液体积
% #HANDWRITTEN: 10
\fieldvalue{\handwritten{10}} ml

电热恒温水槽：
% #VALUE_ID: LEX-P0004-V0023
% #FIELD_VALUE: 电热恒温水槽温度
% #HANDWRITTEN: 37
\fieldvalue{\handwritten{37}}\,°C

裂解：
% #VALUE_ID: LEX-P0004-V0024
% #FIELD_VALUE: 第一次裂解时间
% #TODO #HANDWRITTEN: 09:14--20:44; handwriting partially unclear
\fieldvalue{\handwritten{09:14--20:44}}
}
\\ \hline
\parbox[t]{0.57\textwidth}{\raggedright
取出裂解好的细胞悬液，300$\times g$离心5分钟，按照读取样品体积加入对应体积的1\%人白血球裂解液重悬（每1.2 ml样品体积加入0.1 ml 1\%人白血球裂解液，按照“只进不舍”原则，如26 ml样品加入2.2 ml）。
}
&
\parbox[t]{0.38\textwidth}{\raggedright
% #VALUE_ID: LEX-P0004-V0025
% #FIELD_VALUE: 离心力
% #HANDWRITTEN: 300
\fieldvalue{\handwritten{300}}$\times g$ 离心
% #VALUE_ID: LEX-P0004-V0026
% #FIELD_VALUE: 离心时间
% #HANDWRITTEN: 5
\fieldvalue{\handwritten{5}} 分钟

加入1\%人白细胞裂解液总体积：
% #VALUE_ID: LEX-P0004-V0027
% #FIELD_VALUE: 1\%人白细胞裂解液总体积
% #HANDWRITTEN: 31
\fieldvalue{\handwritten{31}} ml

吸取：
% #VALUE_ID: LEX-P0004-V0028
% #FIELD_VALUE: 吸取体积
% #TODO #HANDWRITTEN: 200; handwriting partially unclear
\fieldvalue{\handwritten{200}}\,\textmu l
}
\\ \hline
\end{tabularx}

\begin{center}
第2页 共5页
\end{center}
% LEXOID_PAGE_COMPLETED: 4/60

\newpage

\begin{center}
\small
\textbf{检测结果记录表}
\end{center}

\noindent
文档编号：REC-YS-ZOP-CQ-99-006-a

\vspace{0.3em}

\noindent
加入 $2.2\mathrm{ml}$ 的氯化钠溶液后，反复吹打至单细胞悬液为。 与 $2\mu\mathrm{l}$ AO/PI 染料混合均匀。

\noindent
撕去计数板上下两层保护膜，将细胞悬液混匀，从中吸取 $20\mu\mathrm{l}$ 与 $20\mu\mathrm{l}$ AO/PI 染料混合均匀。

\noindent
取吸
% #VALUE_ID: LEX-P0005-V0001
% #FIELD_VALUE: 吸取体积
% #HANDWRITTEN: 2
\fieldvalue{\handwritten{2}}
$\mu\mathrm{l}$ 计数。

\noindent
每个样品均计数
% #VALUE_ID: LEX-P0005-V0002
% #FIELD_VALUE: 样品计数次数
% #HANDWRITTEN: 2
\fieldvalue{\handwritten{2}}
次。

\noindent
取吸 $20\mu\mathrm{l}$ 细胞悬液与 AO/PI 染料的混合液，注入到计数板的一 个检测室中。

\noindent
将加入混合液的一端插入进板接口。

\noindent
选择 Assay 类型，使用 AO/PI 染料进行荧光光度计检测；

\noindent
在主界面左侧 ID 栏中输入样本批号。

\noindent
点击主界面左下角“Preview”预览细胞图片。

\noindent
调整仪器焦距直至观察到荧光通道下的细胞中心清亮，轮廓清晰。

\noindent
点击“count”按钮进行计数并自动拍摄图片。

\noindent
重复以上步骤 2 次，共计数 3 次。打印数据结果。

\begin{tabularx}{\textwidth}{|p{0.18\textwidth}|X|}
\hline
数据保存路径 \hfill
% #VALUE_ID: LEX-P0005-V0003
% #FIELD_VALUE: 数据保存路径
% #TODO #HANDWRITTEN: 14/10...；后续路径文字难以辨认
\fieldvalue{\handwritten{14/10\ldots}} \\
\hline
\multicolumn{2}{|c|}{检测结果} \\
\hline
\multicolumn{2}{|l|}{检测结果：□ 一级反应血、□ 二级反应血、三级反应血} \\
\hline
\end{tabularx}

\vspace{0.2em}

\small
\begin{tabularx}{\textwidth}{|p{0.10\textwidth}|X|X|X|X|X|X|}
\hline
检测次数 &
活细胞密度（个细胞/ml） &
样品密度（个细胞/ml） &
细胞活率度 CV 值（\%） &
细胞活率（\%） &
样品活率（\%） &
细胞活率 CV 值（\%） \\
\hline
第一次 &
% #VALUE_ID: LEX-P0005-V0004
% #FIELD_VALUE: 第一次活细胞密度
% #TODO #HANDWRITTEN: 1.13$\times 10^{?}$；指数难以辨认
\fieldvalue{\handwritten{1.13$\times 10^{?}$}} &
&
&
% #VALUE_ID: LEX-P0005-V0005
% #FIELD_VALUE: 第一次细胞活率
% #HANDWRITTEN: 89.4
\fieldvalue{\handwritten{89.4}} &
& \\
\hline
第二次 &
% #VALUE_ID: LEX-P0005-V0006
% #FIELD_VALUE: 第二次活细胞密度
% #TODO #HANDWRITTEN: 1.30$\times 10^{?}$；指数难以辨认
\fieldvalue{\handwritten{1.30$\times 10^{?}$}} &
% #VALUE_ID: LEX-P0005-V0007
% #FIELD_VALUE: 样品密度
% #TODO #HANDWRITTEN: 9.4$\times 10^{?}$；指数难以辨认
\fieldvalue{\handwritten{9.4$\times 10^{?}$}} &
% #VALUE_ID: LEX-P0005-V0008
% #FIELD_VALUE: 细胞活率 CV 值
% #HANDWRITTEN: 13
\fieldvalue{\handwritten{13}} &
% #VALUE_ID: LEX-P0005-V0009
% #FIELD_VALUE: 第二次细胞活率
% #HANDWRITTEN: 84.1
\fieldvalue{\handwritten{84.1}} &
% #VALUE_ID: LEX-P0005-V0010
% #FIELD_VALUE: 样品活率
% #HANDWRITTEN: 89
\fieldvalue{\handwritten{89}} &
% #VALUE_ID: LEX-P0005-V0011
% #FIELD_VALUE: 细胞活率 CV 值
% #HANDWRITTEN: 2
\fieldvalue{\handwritten{2}} \\
\hline
第三次 &
% #VALUE_ID: LEX-P0005-V0012
% #FIELD_VALUE: 第三次活细胞密度
% #TODO #HANDWRITTEN: 0.44$\times 10^{?}$；指数难以辨认
\fieldvalue{\handwritten{0.44$\times 10^{?}$}} &
&
&
% #VALUE_ID: LEX-P0005-V0013
% #FIELD_VALUE: 第三次细胞活率
% #HANDWRITTEN: 81.8
\fieldvalue{\handwritten{81.8}} &
& \\
\hline
\multicolumn{2}{|l|}{发酵罐细胞体积（ml）：
% #VALUE_ID: LEX-P0005-V0014
% #FIELD_VALUE: 发酵罐细胞体积
% #TODO #HANDWRITTEN: 890；首位数字辨认不清
\fieldvalue{\handwritten{890}}} &
\multicolumn{2}{l|}{细胞数量（个细胞/罐）} &
\multicolumn{3}{l|}{
% #VALUE_ID: LEX-P0005-V0015
% #FIELD_VALUE: 细胞数量
% #TODO #HANDWRITTEN: 8.13$\times 10^{?}$；指数难以辨认
\fieldvalue{\handwritten{8.13$\times 10^{?}$}}} \\
\hline
\end{tabularx}

\normalsize

\noindent
计算要求：样品密度为 3 次细胞密度结果取平均；细胞数以 1\% 人白血红细胞浓度为基本代表，再以生物反应器细胞浓度换算。

\noindent
细胞液体积；细胞数量为样品密度乘以发酵罐体积。

\begin{tabularx}{\textwidth}{|p{0.18\textwidth}|X|}
\hline
备注 &
% #VALUE_ID: LEX-P0005-V0016
% #FIELD_VALUE: 备注
% #HANDWRITTEN: /
\fieldvalue{\handwritten{/}} \\
\hline
实验人/日期 &
% #VALUE_ID: LEX-P0005-V0017
% #FIELD_VALUE: 实验人/日期
% #TODO #HANDWRITTEN: 李海生 2025.08.20；姓名及日期部分难以辨认
\fieldvalue{\handwritten{李海生\quad 2025.08.20}} \\
\hline
复核人/日期 &
% #VALUE_ID: LEX-P0005-V0018
% #FIELD_VALUE: 复核人/日期
% #TODO #HANDWRITTEN: 胡鸿 2025.08.30；姓名部分难以辨认
\fieldvalue{\handwritten{胡鸿\quad 2025.08.30}} \\
\hline
\end{tabularx}

% TODO: 页面顶部的公司标识、版本信息及印章为固定表单内容，未作为可编辑值标注。
% LEXOID_PAGE_COMPLETED: 5/60

\newpage

\begin{center}
\textbf{生物反应器细胞液检验操作记录}
\end{center}

\noindent 文件编码：REC-YZ-SOP-QC-99-006-a

\begin{center}
\begin{tabularx}{\textwidth}{|p{2.2cm}|X|p{2.2cm}|X|}
\hline
\multicolumn{4}{|c|}{细胞生长状态检验}\\
\hline
样品名称 &
% #VALUE_ID: LEX-P0006-V0001
% #FIELD_VALUE: 样品名称
\fieldvalue{皂生物反应器细胞液} &
样品规格 &
% #VALUE_ID: LEX-P0006-V0002
% #FIELD_VALUE: 样品规格
% #HANDWRITTEN: 20
\fieldvalue{\handwritten{20}} ml/瓶\\
\hline
样品批号 &
% #VALUE_ID: LEX-P0006-V0003
% #FIELD_VALUE: 样品批号
% #TODO #HANDWRITTEN: A312?20?5?08?; handwriting unclear
\fieldvalue{\handwritten{A312?20?5?08?}} &
取样日期 &
% #VALUE_ID: LEX-P0006-V0004
% #FIELD_VALUE: 取样日期
% #HANDWRITTEN: 2025年08月30日
\fieldvalue{\handwritten{2025年08月30日}}\\
\hline
\end{tabularx}
\end{center}

\noindent\textbf{1.\ 仪器设备及试剂：}

\noindent 1.1．仪器设备：

\begin{center}
\begin{tabularx}{\textwidth}{|X|X|X|X|X|}
\hline
仪器名称 & 厂家 & 型号 & 设备编号 & 检查确认\\
\hline
% #VALUE_ID: LEX-P0006-V0005
% #FIELD_VALUE: 仪器名称
\fieldvalue{荧光显微镜} &
% #VALUE_ID: LEX-P0006-V0006
% #FIELD_VALUE: 厂家
\fieldvalue{徕卡} &
% #VALUE_ID: LEX-P0006-V0007
% #FIELD_VALUE: 型号
\fieldvalue{DmIL LED} &
% #VALUE_ID: LEX-P0006-V0008
% #FIELD_VALUE: 设备编号
\fieldvalue{1431404} &
% #VALUE_ID: LEX-P0006-V0009
% #FIELD_VALUE: 检查确认
\fieldvalue{\(\boxtimes\)正常 \(\square\)异常}\\
\hline
\end{tabularx}
\end{center}

\noindent 1.2．试剂耗材：

\begin{center}
\begin{tabularx}{\textwidth}{|X|X|X|X|X|}
\hline
名称 & 厂家 & 货号 & 批号 & 有效期至\\
\hline
% #VALUE_ID: LEX-P0006-V0010
% #FIELD_VALUE: 名称
\fieldvalue{细胞活力检测试剂盒} &
% #VALUE_ID: LEX-P0006-V0011
% #FIELD_VALUE: 厂家
\fieldvalue{江苏凯基生物技术股份有限公司} &
% #VALUE_ID: LEX-P0006-V0012
% #FIELD_VALUE: 货号
\fieldvalue{KGA9501-1000} &
% #VALUE_ID: LEX-P0006-V0013
% #FIELD_VALUE: 批号
% #HANDWRITTEN: 20250410
\fieldvalue{\handwritten{20250410}} &
% #VALUE_ID: LEX-P0006-V0014
% #FIELD_VALUE: 有效期至
% #HANDWRITTEN: 2025.10.09
\fieldvalue{\handwritten{2025.10.09}}\\
\hline
\end{tabularx}
\end{center}

\noindent\textbf{2．操作步骤：}

\begin{center}
\small
\begin{tabularx}{\textwidth}{|p{0.49\textwidth}|X|}
\hline
\multicolumn{2}{|c|}{溶液配制}\\
\hline
20$\times$的荧光染液配制：取
% #VALUE_ID: LEX-P0006-V0015
% #FIELD_VALUE: 荧光染液配制中组分A体积
% #HANDWRITTEN: 0.2
\fieldvalue{\handwritten{0.2}} $\mu$l 组分A 和
% #VALUE_ID: LEX-P0006-V0016
% #FIELD_VALUE: 荧光染液配制中组分B体积
% #HANDWRITTEN: 0.2
\fieldvalue{\handwritten{0.2}} $\mu$l 组分B 与
% #VALUE_ID: LEX-P0006-V0017
% #FIELD_VALUE: 荧光染液配制中PBS体积
% #HANDWRITTEN: 39.1
\fieldvalue{\handwritten{39.1}} $\mu$l PBS 混合均匀配成20$\times$的荧光染液。&
\\
\hline
\multicolumn{2}{|c|}{操作过程}\\
\hline
取混匀后样品200$\mu$l/孔至96孔板中，取2孔。&
样品
% #VALUE_ID: LEX-P0006-V0018
% #FIELD_VALUE: 样品加样体积
% #HANDWRITTEN: 200
\fieldvalue{\handwritten{200}} $\mu$l/孔\\
&
每个样品取
% #VALUE_ID: LEX-P0006-V0019
% #FIELD_VALUE: 每个样品取孔数
% #HANDWRITTEN: 2
\fieldvalue{\handwritten{2}} 孔\\
\hline
自然沉降后，吸弃上清，注意不要吸到沉淀样品。&
% #VALUE_ID: LEX-P0006-V0020
% #FIELD_VALUE: 上清处理
\fieldvalue{\(\boxtimes\)吸弃上清}\\
\hline
取20$\times$的荧光染液加入到96孔板中，200$\mu$l/孔，避光孵育30--60min。&
20$\times$的荧光染液
% #VALUE_ID: LEX-P0006-V0021
% #FIELD_VALUE: 荧光染液加入体积
% #HANDWRITTEN: 200
\fieldvalue{\handwritten{200}} $\mu$l/孔\\
&
避光孵育
% #VALUE_ID: LEX-P0006-V0022
% #FIELD_VALUE: 避光孵育时间
% #HANDWRITTEN: 01:20～01:40
\fieldvalue{\handwritten{01:20\textasciitilde01:40}}\\
\hline
染色结束后，吸弃上清，取PBS加入到96孔板中，200$\mu$l/孔，吸弃上清，注意不要吸到沉淀样品，清洗3次，之后每孔加入100$\mu$l PBS进行拍照！&
加入PBS
% #VALUE_ID: LEX-P0006-V0023
% #FIELD_VALUE: PBS加入体积
% #HANDWRITTEN: 200
\fieldvalue{\handwritten{200}} $\mu$l/孔\\
&
清洗
% #VALUE_ID: LEX-P0006-V0024
% #FIELD_VALUE: 清洗次数
% #HANDWRITTEN: 2
\fieldvalue{\handwritten{2}} 次\\
\hline
将荧光显微镜调至4$\times$物镜，软件中调节荧光参数，调节到合适的光源和曝光，拍摄4张清晰图像。&
\\
\hline
选一张清晰图像，打印并粘贴在实验记录上。&
\\
\hline
数据保存路径&
% #VALUE_ID: LEX-P0006-V0025
% #FIELD_VALUE: 数据保存路径
% #TODO #HANDWRITTEN: 1431404电脑D:\...过程记录\2024年; path is partially illegible
\fieldvalue{\handwritten{1431404电脑D:\textbackslash ...\textbackslash 过程记录\textbackslash 2024年}}\\
\hline
\multicolumn{2}{|c|}{检测结果}\\
\hline
\end{tabularx}
\end{center}

\begin{center}
第4页 共5页
\end{center}
% LEXOID_PAGE_COMPLETED: 6/60

\newpage

\noindent
\begin{tabularx}{\textwidth}{@{}Xr@{}}
\textbf{维生基因生物科技（上海）有限公司} &
文件编码：REC-YZ-SOP-QC-99-006-a\\
\textit{PENTA LIFE EXCELLENCE BIOTECH} &
版本号：14
\end{tabularx}

\vspace{0.2cm}

\noindent
\begin{center}
\setlength{\fboxsep}{8pt}
\fbox{%
\begin{minipage}[t][17.2cm][t]{0.93\textwidth}
\vspace{3.0cm}
\begin{center}
% #VALUE_ID: LEX-P0007-V0001
% #TODO #HANDWRITTEN: 手写标注不清；图像上方蓝色字迹模糊
\handwritten{\text{手写标注不清}}

\vspace{0.1cm}
% TODO: The embedded specimen photograph is visible on this page, but no source filename is available.
\fbox{\parbox[c][3.0cm][c]{5.0cm}{\centering
\textit{图像占位符}\\[-1mm]
\small（此处为显微图像）
}}
\end{center}
\end{minipage}%
}
\end{center}

\vspace{-0.05cm}

\noindent
\begin{tabularx}{0.93\textwidth}{|p{2.0cm}|X|}
\hline
备注 &
% #VALUE_ID: LEX-P0007-V0002
% #FIELD_VALUE: 备注
% #HANDWRITTEN: ✓
\fieldvalue{\handwritten{\checkmark}}
\\
\hline
实验人/日期 &
% #VALUE_ID: LEX-P0007-V0003
% #FIELD_VALUE: 实验人
% #TODO #HANDWRITTEN: 姓名不清；蓝色手写字迹难以辨认
\fieldvalue{\handwritten{\text{姓名不清}}}
\hspace{1.5cm}
% #VALUE_ID: LEX-P0007-V0004
% #FIELD_VALUE: 实验日期
% #HANDWRITTEN: 2025.08.30
\fieldvalue{\handwritten{2025.08.30}}
\\
\hline
复核人/日期 &
% #VALUE_ID: LEX-P0007-V0005
% #FIELD_VALUE: 复核人
% #TODO #HANDWRITTEN: 姓名不清；蓝色手写字迹难以辨认
\fieldvalue{\handwritten{\text{姓名不清}}}
\hspace{1.5cm}
% #VALUE_ID: LEX-P0007-V0006
% #FIELD_VALUE: 复核日期
% #HANDWRITTEN: 2025.08.30
\fieldvalue{\handwritten{2025.08.30}}
\\
\hline
\end{tabularx}

\vfill

\begin{center}
\small 第 5 页 共 5 页
\end{center}
% LEXOID_PAGE_COMPLETED: 7/60

\newpage

% #VALUE_ID: LEX-P0008-V0001
% #FIELD_VALUE: Assay
Assay: \fieldvalue{MSC-1}

% #VALUE_ID: LEX-P0008-V0002
% #FIELD_VALUE: Cell Type F1
Cell Type F1: \fieldvalue{MSC (AO)}

% #VALUE_ID: LEX-P0008-V0003
% #FIELD_VALUE: Cell Type F2
Cell Type F2: \fieldvalue{MSC (PI)}

% #VALUE_ID: LEX-P0008-V0004
% #FIELD_VALUE: Sample ID
Sample ID: \fieldvalue{QC-JS-QY0402-A31Z2202508008-20250829-24.5H-1}

% #VALUE_ID: LEX-P0008-V0005
% #FIELD_VALUE: Dilution
Dilution: \fieldvalue{2.00}

\vspace{0.5em}
\hrule
\vspace{0.3em}

Results:

\vspace{0.5em}
\renewcommand{\arraystretch}{1.15}
\begin{tabularx}{\textwidth}{@{}X X X@{}}
\textbf{Count} & \textbf{Concentration} & \textbf{Mean Diameter} \\
\underline{\rule{0pt}{1.2ex}\texttt{=============}} &
\underline{\rule{0pt}{1.2ex}\texttt{=============}} &
\underline{\rule{0pt}{1.2ex}\texttt{=============}} \\

Total:
% #VALUE_ID: LEX-P0008-V0006
% #FIELD_VALUE: Total count
\fieldvalue{358}
&
% #VALUE_ID: LEX-P0008-V0007
% #FIELD_VALUE: Total concentration
\fieldvalue{$1.26\times10^{6}$ cells/mL}
&
% #VALUE_ID: LEX-P0008-V0008
% #FIELD_VALUE: Total mean diameter
\fieldvalue{15.1 micron}
\\

Live:
% #VALUE_ID: LEX-P0008-V0009
% #FIELD_VALUE: Live count
\fieldvalue{320}
&
% #VALUE_ID: LEX-P0008-V0010
% #FIELD_VALUE: Live concentration
\fieldvalue{$1.13\times10^{6}$ cells/mL}
&
% #VALUE_ID: LEX-P0008-V0011
% #FIELD_VALUE: Live mean diameter
\fieldvalue{15.1 microns}
\\

Dead:
% #VALUE_ID: LEX-P0008-V0012
% #FIELD_VALUE: Dead count
\fieldvalue{38}
&
% #VALUE_ID: LEX-P0008-V0013
% #FIELD_VALUE: Dead concentration
\fieldvalue{$1.34\times10^{6}$ cells/mL}
&
% #VALUE_ID: LEX-P0008-V0014
% #FIELD_VALUE: Dead mean diameter
\fieldvalue{14.4 microns}
\\

Viability:
% #VALUE_ID: LEX-P0008-V0015
% #FIELD_VALUE: Viability
\fieldvalue{89.4\%}
&
&
\end{tabularx}

\vspace{1em}

\begin{center}
\begin{tabular}{@{}c@{\hspace{1.5em}}c@{}}
\begin{minipage}[t]{0.40\textwidth}
\centering
% TODO: Insert the upper-left microscopy image; no filename is available.
\fbox{\rule{0pt}{0.22\textwidth}\rule{0.9\linewidth}{0pt}}
\end{minipage}
&
\begin{minipage}[t]{0.40\textwidth}
\centering
% TODO: Insert the upper-right microscopy image; no filename is available.
\fbox{\rule{0pt}{0.22\textwidth}\rule{0.9\linewidth}{0pt}}
\end{minipage}
\\[1em]
\begin{minipage}[t]{0.40\textwidth}
\centering
% TODO: Insert the lower-left microscopy image; no filename is available.
\fbox{\rule{0pt}{0.22\textwidth}\rule{0.9\linewidth}{0pt}}
\end{minipage}
&
\begin{minipage}[t]{0.40\textwidth}
\centering
% TODO: Insert the lower-right microscopy image; no filename is available.
\fbox{\rule{0pt}{0.22\textwidth}\rule{0.9\linewidth}{0pt}}
\end{minipage}
\end{tabular}
\end{center}

\vfill

操作人/审核人：
% #VALUE_ID: LEX-P0008-V0016
% #FIELD_VALUE: 操作人/审核人
% #TODO #HANDWRITTEN: 邓鹏生; handwritten Chinese name is somewhat unclear
\fieldvalue{\handwritten{邓鹏生}}
\hfill
日期：
% #VALUE_ID: LEX-P0008-V0017
% #FIELD_VALUE: 日期
% #HANDWRITTEN: 2025.08.30
\fieldvalue{\handwritten{2025.08.30}}
% LEXOID_PAGE_COMPLETED: 8/60

\newpage

Assay: 
% #VALUE_ID: LEX-P0009-V0001
% #FIELD_VALUE: Assay
\fieldvalue{MSC-1}

Cell Type F1: 
% #VALUE_ID: LEX-P0009-V0002
% #FIELD_VALUE: Cell Type F1
\fieldvalue{MSC (AO)}

Cell Type F2: 
% #VALUE_ID: LEX-P0009-V0003
% #FIELD_VALUE: Cell Type F2
\fieldvalue{MSC (PI)}

Sample ID: 
% #VALUE_ID: LEX-P0009-V0004
% #FIELD_VALUE: Sample ID
\fieldvalue{QC-JS-QY0402-A31Z2202508008-20250829-24.5H-2}

Dilution: 
% #VALUE_ID: LEX-P0009-V0005
% #FIELD_VALUE: Dilution
\fieldvalue{2.00}

\noindent\rule{\textwidth}{0.4pt}

Results:

\vspace{0.5em}

\begin{tabularx}{\textwidth}{@{}XXX@{}}
\textbf{Count} & \textbf{Concentration} & \textbf{Mean Diameter} \\
\underline{\hspace{0.9\linewidth}} &
\underline{\hspace{0.9\linewidth}} &
\underline{\hspace{0.9\linewidth}} \\[0.2em]
Total: 
% #VALUE_ID: LEX-P0009-V0006
% #FIELD_VALUE: Total count
\fieldvalue{375}
&
% #VALUE_ID: LEX-P0009-V0007
% #FIELD_VALUE: Total concentration
\fieldvalue{1.32$\times 10^{6}$ cells/mL}
&
% #VALUE_ID: LEX-P0009-V0008
% #FIELD_VALUE: Total mean diameter
\fieldvalue{13.6 micron}
\\
Live: 
% #VALUE_ID: LEX-P0009-V0009
% #FIELD_VALUE: Live count
\fieldvalue{342}
&
% #VALUE_ID: LEX-P0009-V0010
% #FIELD_VALUE: Live concentration
\fieldvalue{1.20$\times 10^{6}$ cells/mL}
&
% #VALUE_ID: LEX-P0009-V0011
% #FIELD_VALUE: Live mean diameter
\fieldvalue{13.5 microns}
\\
Dead: 
% #VALUE_ID: LEX-P0009-V0012
% #FIELD_VALUE: Dead count
\fieldvalue{33}
&
% #VALUE_ID: LEX-P0009-V0013
% #FIELD_VALUE: Dead concentration
\fieldvalue{1.16$\times 10^{5}$ cells/mL}
&
% #VALUE_ID: LEX-P0009-V0014
% #FIELD_VALUE: Dead mean diameter
\fieldvalue{14.2 microns}
\\
Viability: 
% #VALUE_ID: LEX-P0009-V0015
% #FIELD_VALUE: Viability
\fieldvalue{91.2\%}
&
&
\end{tabularx}

\vspace{1em}

% TODO: Four microscopy image panels are visible on this page; source filenames are unavailable.
\begin{center}
\setlength{\fboxsep}{4pt}
\begin{tabular}{cc}
\fbox{\rule{0pt}{0.20\textheight}\rule{0.43\textwidth}{0pt}} &
\fbox{\rule{0pt}{0.20\textheight}\rule{0.43\textwidth}{0pt}} \\[1em]
\fbox{\rule{0pt}{0.20\textheight}\rule{0.43\textwidth}{0pt}} &
\fbox{\rule{0pt}{0.20\textheight}\rule{0.43\textwidth}{0pt}}
\end{tabular}
\end{center}

\vfill

\noindent 操作人/审核人：
% #VALUE_ID: LEX-P0009-V0016
% #FIELD_VALUE: 操作人/审核人
% #HANDWRITTEN: 陈彬生
\fieldvalue{\handwritten{陈彬生}}
\hfill
日期：
% #VALUE_ID: LEX-P0009-V0017
% #FIELD_VALUE: 日期
% #HANDWRITTEN: 2025.08.30
\fieldvalue{\handwritten{2025.08.30}}
% LEXOID_PAGE_COMPLETED: 9/60

\newpage

\noindent\textbf{Assay:} %
% #VALUE_ID: LEX-P0010-V0001
% #FIELD_VALUE: Assay
\fieldvalue{MSC-1}

\medskip
\noindent\textbf{Cell Type F1:} %
% #VALUE_ID: LEX-P0010-V0002
% #FIELD_VALUE: Cell Type F1
\fieldvalue{MSC (AO)}

\noindent\textbf{Cell Type F2:} %
% #VALUE_ID: LEX-P0010-V0003
% #FIELD_VALUE: Cell Type F2
\fieldvalue{MSC (PI)}

\medskip
\noindent\textbf{Sample ID:} %
% #VALUE_ID: LEX-P0010-V0004
% #FIELD_VALUE: Sample ID
\fieldvalue{QC-JS-QY0402-A31Z2202508008-20250829-24.5H-3}

\noindent\textbf{Dilution:} %
% #VALUE_ID: LEX-P0010-V0005
% #FIELD_VALUE: Dilution
\fieldvalue{2.00}

\medskip
\hrule
\medskip

\noindent\textbf{Results:}

\medskip
\renewcommand{\arraystretch}{1.15}
\begin{tabularx}{\textwidth}{@{}>{\raggedright\arraybackslash}X
>{\raggedright\arraybackslash}X
>{\raggedright\arraybackslash}X@{}}
\textbf{Count} & \textbf{Concentration} & \textbf{Mean Diameter} \\
\underline{\hspace{1.25cm}} & \underline{\hspace{1.65cm}} & \underline{\hspace{1.75cm}} \\

\textbf{Total:} %
% #VALUE_ID: LEX-P0010-V0006
% #FIELD_VALUE: Total count
\fieldvalue{303}
&
% #VALUE_ID: LEX-P0010-V0007
% #FIELD_VALUE: Total concentration
\fieldvalue{1.07$\times$10$^{6}$ cells/mL}
&
% #VALUE_ID: LEX-P0010-V0008
% #FIELD_VALUE: Total mean diameter
\fieldvalue{15.9 micron}
\\

\textbf{Live:} %
% #VALUE_ID: LEX-P0010-V0009
% #FIELD_VALUE: Live count
\fieldvalue{266}
&
% #VALUE_ID: LEX-P0010-V0010
% #FIELD_VALUE: Live concentration
\fieldvalue{9.44$\times$10$^{5}$ cells/mL}
&
% #VALUE_ID: LEX-P0010-V0011
% #FIELD_VALUE: Live mean diameter
\fieldvalue{16.0 microns}
\\

\textbf{Dead:} %
% #VALUE_ID: LEX-P0010-V0012
% #FIELD_VALUE: Dead count
\fieldvalue{37}
&
% #VALUE_ID: LEX-P0010-V0013
% #FIELD_VALUE: Dead concentration
\fieldvalue{1.31$\times$10$^{5}$ cells/mL}
&
% #VALUE_ID: LEX-P0010-V0014
% #FIELD_VALUE: Dead mean diameter
\fieldvalue{14.9 microns}
\\

\textbf{Viability:} %
% #VALUE_ID: LEX-P0010-V0015
% #FIELD_VALUE: Viability
\fieldvalue{87.8\%}
&
&
\end{tabularx}

\medskip

\begin{center}
\begin{tabular}{cc}
\fbox{\parbox[c][4.0cm][c]{0.40\textwidth}{\centering\textit{Microscopy image placeholder}}}
&
\fbox{\parbox[c][4.0cm][c]{0.40\textwidth}{\centering\textit{Microscopy image placeholder}}}
\\[0.5cm]
\fbox{\parbox[c][4.0cm][c]{0.40\textwidth}{\centering\textit{Microscopy image placeholder}}}
&
\fbox{\parbox[c][4.0cm][c]{0.40\textwidth}{\centering\textit{Microscopy image placeholder}}}
\end{tabular}
\end{center}

\vfill

\noindent 操作人/审核人：%
% #VALUE_ID: LEX-P0010-V0016
% #FIELD_VALUE: 操作人/审核人
% #TODO #HANDWRITTEN: 郭?生; handwritten Chinese signature is unclear
\fieldvalue{\handwritten{郭?生}}
\hfill
日期：%
% #VALUE_ID: LEX-P0010-V0017
% #FIELD_VALUE: 日期
% #TODO #HANDWRITTEN: 2025.8.30; handwritten date is partially unclear
\fieldvalue{\handwritten{2025.8.30}}
% LEXOID_PAGE_COMPLETED: 10/60

\newpage

\begin{center}
\textbf{铂莱生物科技（北京）有限公司}\\
\textit{PLATINUM LIFE \& EXCELLENCE BIOTECH}\\[2pt]
\textbf{检验单}
\end{center}

\noindent
文件编号：
% #VALUE_ID: LEX-P0011-V0001
% #FIELD_VALUE: 文件编号
\fieldvalue{REC-YZ-SMP-QC-01-001-d}
\hfill
版本号：
% #VALUE_ID: LEX-P0011-V0002
% #FIELD_VALUE: 版本号
\fieldvalue{1.4}

\begin{center}
\renewcommand{\arraystretch}{1.45}
\small
\begin{tabularx}{\textwidth}{|p{2.5cm}|X|p{2.5cm}|X|}
\hline
样品名称 &
% #VALUE_ID: LEX-P0011-V0003
% #FIELD_VALUE: 样品名称
% #TODO #HANDWRITTEN: 三氯生的检验样品; handwriting partially unclear
\fieldvalue{\handwritten{三氯生的检验样品}} &
样品代码 &
% #VALUE_ID: LEX-P0011-V0004
% #FIELD_VALUE: 样品代码
\fieldvalue{QY0402}
\\ \hline

样品批号 &
% #VALUE_ID: LEX-P0011-V0005
% #FIELD_VALUE: 样品批号
% #TODO #HANDWRITTEN: 301221?0809; handwritten digits partially unclear
\fieldvalue{\handwritten{301221?0809}} &
样品规格/数量 &
% #VALUE_ID: LEX-P0011-V0006
% #FIELD_VALUE: 样品规格/数量
% #TODO #HANDWRITTEN: 4/20 ml/瓶; unit and preceding quantity partially unclear
\fieldvalue{\handwritten{4/20 ml/瓶}}
\\ \hline

样品类别 &
% #VALUE_ID: LEX-P0011-V0007
% #FIELD_VALUE: 样品类别
\fieldvalue{\checkmark\ 生产过程中检验样品}
&
\multicolumn{2}{l|}{%
\(\square\) 待包装产品\qquad
\(\square\) 成品\qquad
\(\square\) 其他（\hspace{1.2cm}）}
\\ \hline

检验类型 &
% #VALUE_ID: LEX-P0011-V0008
% #FIELD_VALUE: 检验类型
\fieldvalue{\checkmark\ 初验}
&
\multicolumn{2}{l|}{%
\(\square\) 复验\qquad
\(\square\) 退货\qquad
\(\square\) 稳定性\qquad
\(\square\) 留样}
\\ \hline

请验目的 &
% #VALUE_ID: LEX-P0011-V0009
% #FIELD_VALUE: 请验目的
\fieldvalue{\checkmark\ 全检}
&
\multicolumn{2}{l|}{%
\(\square\) 留样\qquad
\(\square\) 其他（\hspace{2.2cm}）}
\\ \hline

请验部门 &
% #VALUE_ID: LEX-P0011-V0010
% #FIELD_VALUE: 请验部门
\fieldvalue{\checkmark\ 生产部}
&
\multicolumn{2}{l|}{%
\(\square\) 物料管理部\qquad
\(\square\) 其他（\hspace{1.6cm}）}
\\ \hline

样品情况 &
% #VALUE_ID: LEX-P0011-V0011
% #FIELD_VALUE: 样品情况
\fieldvalue{\checkmark\ 随请验单一起提供}
&
\multicolumn{2}{l|}{%
\(\square\) 未随请验单一起提供}
\\ \hline

请验人 &
% #VALUE_ID: LEX-P0011-V0012
% #FIELD_VALUE: 请验人
% #TODO #HANDWRITTEN: 侯?梅; handwriting unclear
\fieldvalue{\handwritten{侯?梅}}
&
请验日期 &
% #VALUE_ID: LEX-P0011-V0013
% #FIELD_VALUE: 请验日期
% #HANDWRITTEN: 2023.08.29
\fieldvalue{\handwritten{2023.08.29}}
\\ \hline

收验人 &
% #VALUE_ID: LEX-P0011-V0014
% #FIELD_VALUE: 收验人
% #TODO #HANDWRITTEN: 何?荣; handwriting unclear
\fieldvalue{\handwritten{何?荣}}
&
收验日期 &
% #VALUE_ID: LEX-P0011-V0015
% #FIELD_VALUE: 收验日期
% #HANDWRITTEN: 2023.8.29
\fieldvalue{\handwritten{2023.8.29}}
\\ \hline

备注 &
% #VALUE_ID: LEX-P0011-V0016
% #FIELD_VALUE: 备注
% #TODO #HANDWRITTEN: 二次 0h 9.0L; handwriting partially unclear
\fieldvalue{\handwritten{二次\quad 0h\quad 9.0L}}
&
\multicolumn{2}{l|}{}
\\ \hline
\end{tabularx}
\end{center}

% TODO: The form continues beyond the visible content on this page.
% LEXOID_PAGE_COMPLETED: 11/60

\newpage

\begin{center}
\textbf{取样记录}
\end{center}

\begin{tabularx}{\textwidth}{|p{2.2cm}|X|p{2.2cm}|X|}
\hline
名称 &
% #VALUE_ID: LEX-P0012-V0001
% #FIELD_VALUE: 名称
\fieldvalue{二级生物反应器细胞培养液} &
代码 &
% #VALUE_ID: LEX-P0012-V0002
% #FIELD_VALUE: 代码
\fieldvalue{QY0402}
\\ \hline
规格 &
% #VALUE_ID: LEX-P0012-V0003
% #FIELD_VALUE: 规格
\fieldvalue{约20ml/袋} &
总数量 &
% #VALUE_ID: LEX-P0012-V0004
% #FIELD_VALUE: 总数量
\fieldvalue{1袋}
\\ \hline
批号 &
% #VALUE_ID: LEX-P0012-V0005
% #FIELD_VALUE: 批号
\fieldvalue{A31Z2202508008} &
进厂批号 &
% #VALUE_ID: LEX-P0012-V0006
% #FIELD_VALUE: 进厂批号
\fieldvalue{/}
\\ \hline
供货单位/厂家 &
\multicolumn{3}{l|}{%
% #VALUE_ID: LEX-P0012-V0007
% #FIELD_VALUE: 供货单位/厂家
\fieldvalue{/}} \\ \hline
取样目的 &
\multicolumn{3}{l|}{%
% #VALUE_ID: LEX-P0012-V0008
% #FIELD_VALUE: 取样目的--正常检验
% #HANDWRITTEN: ✓ 正常检验
\fieldvalue{\handwritten{✓\ 正常检验}}\qquad
$\square$ 复检\qquad $\square$ 稳定性
\newline
$\square$ 退货检验\qquad $\square$ 留样\qquad $\square$ 其他：\underline{\hspace{2cm}}
} \\ \hline
取样地点 &
\multicolumn{3}{l|}{%
% #VALUE_ID: LEX-P0012-V0009
% #FIELD_VALUE: 取样地点--产品出厂
% #HANDWRITTEN: ✓ 产品出厂
\fieldvalue{\handwritten{✓\ 产品出厂}}\qquad
$\square$ 物料库\qquad $\square$ 细胞库\qquad $\square$ 包装间
} \\ \hline
贮存条件 &
\multicolumn{3}{l|}{%
% #VALUE_ID: LEX-P0012-V0010
% #FIELD_VALUE: 贮存条件--2--8°C
% #HANDWRITTEN: ✓ 2--8°C
$\fieldvalue{\handwritten{✓\ 2--8^\circ\mathrm{C}}}$\qquad
$\square$ -20°C\qquad $\square$ -80°C\qquad $\square$ 常温\qquad
$\square$ 其他：\underline{\hspace{1.5cm}}
} \\ \hline
盛装容器 &
\multicolumn{3}{l|}{%
% #VALUE_ID: LEX-P0012-V0011
% #FIELD_VALUE: 盛装容器--保温箱
% #HANDWRITTEN: ✓ 保温箱
\fieldvalue{\handwritten{✓\ 保温箱}}\qquad
$\square$ 泡沫箱\qquad $\square$ 手提盒\qquad
$\square$ 其他：\underline{\hspace{2cm}}
} \\ \hline
外观复核 &
\multicolumn{3}{l|}{%
% #VALUE_ID: LEX-P0012-V0012
% #FIELD_VALUE: 外观复核--合格
% #HANDWRITTEN: ✓ 合格
\fieldvalue{\handwritten{✓\ 合格}}\qquad
$\square$ 不合格
} \\ \hline
取样量 &
% #VALUE_ID: LEX-P0012-V0013
% #FIELD_VALUE: 取样量
\fieldvalue{1袋} &
检验量 &
% #VALUE_ID: LEX-P0012-V0014
% #FIELD_VALUE: 检验量
\fieldvalue{1袋}
\\ \hline
\multicolumn{1}{|p{2.2cm}|}{留样量} &
% #VALUE_ID: LEX-P0012-V0015
% #FIELD_VALUE: 留样量
\fieldvalue{/} &
\multicolumn{2}{l|}{}
\\ \hline
取样人/日期 &
\begin{tabular}{@{}l@{}}
% #VALUE_ID: LEX-P0012-V0016
% #FIELD_VALUE: 取样人
% #HANDWRITTEN: 签名（难以辨认）
\fieldvalue{\handwritten{签名（难以辨认）}}\\[-1mm]
% #VALUE_ID: LEX-P0012-V0017
% #FIELD_VALUE: 取样日期
% #TODO #HANDWRITTEN: 2025.8.??; 日期末位难以辨认
\fieldvalue{\handwritten{2025.8.??}}
\end{tabular}
&
复核人/日期 &
\begin{tabular}{@{}l@{}}
% #VALUE_ID: LEX-P0012-V0018
% #FIELD_VALUE: 复核人
% #HANDWRITTEN: 签名（难以辨认）
\fieldvalue{\handwritten{签名（难以辨认）}}\\[-1mm]
% #VALUE_ID: LEX-P0012-V0019
% #FIELD_VALUE: 复核日期
% #TODO #HANDWRITTEN: 2025.8.??; 日期末位难以辨认
\fieldvalue{\handwritten{2025.8.??}}
\end{tabular}
\\ \hline
备注 &
\multicolumn{3}{l|}{%
% #VALUE_ID: LEX-P0012-V0020
% #FIELD_VALUE: 备注
% #TODO #HANDWRITTEN: OK 9.26; handwriting partially unclear
\fieldvalue{\handwritten{OK\ 9.26}}
} \\ \hline
\end{tabularx}

% TODO: The handwritten signatures and the final remark are partially difficult to read; transcriptions are marked for review.
% LEXOID_PAGE_COMPLETED: 12/60

\newpage

\begin{center}
生物反应器细胞液葡萄糖检测操作记录
\end{center}

\noindent 文件编码：REC-YZ-SOP-QC-99-006-a

\begin{tabularx}{\textwidth}{|l|X|l|X|}
\hline
\multicolumn{4}{|c|}{葡萄糖浓度检测} \\ \hline
样品名称 &
% #VALUE_ID: LEX-P0013-V0001
% #FIELD_VALUE: 样品名称
\fieldvalue{二级生物反应器细胞液} &
样品规格 &
% #VALUE_ID: LEX-P0013-V0002
% #FIELD_VALUE: 样品规格
% #HANDWRITTEN: 12.0
\fieldvalue{\handwritten{12.0}} ml/袋 \\ \hline
样品批号 &
% #VALUE_ID: LEX-P0013-V0003
% #FIELD_VALUE: 样品批号
% #HANDWRITTEN: 44?4?06
\fieldvalue{\handwritten{44?4?06}} &
取样日期 &
% #VALUE_ID: LEX-P0013-V0004
% #FIELD_VALUE: 取样日期
% #HANDWRITTEN: 2024年08月27日
\fieldvalue{\handwritten{2024年08月27日}} \\ \hline
\end{tabularx}

\section{仪器设备、试剂耗材}

\subsection{仪器设备}

\begin{tabularx}{\textwidth}{|l|X|l|X|X|}
\hline
名称 & 厂家 & 型号 & 设备编号 & 检查确认 \\ \hline
血糖仪 & 三诺生物传感股份有限公司 & 安稳+ &
% #VALUE_ID: LEX-P0013-V0005
% #FIELD_VALUE: 血糖仪设备编号
% #TODO #HANDWRITTEN: /; handwritten mark is unclear
\fieldvalue{\handwritten{/}} &
% #VALUE_ID: LEX-P0013-V0006
% #FIELD_VALUE: 血糖仪检查确认
\fieldvalue{\(\checkmark\) 正常}\quad $\square$ 异常
\\ \hline
移液器 & Eppendorf Corporate & 0.1--2.5\,\textmu l &
% #VALUE_ID: LEX-P0013-V0007
% #FIELD_VALUE: 移液器设备编号
% #HANDWRITTEN: 027?4?6
\fieldvalue{\handwritten{027?4?6}} &
% #VALUE_ID: LEX-P0013-V0008
% #FIELD_VALUE: 移液器检查确认
\fieldvalue{\(\checkmark\) 正常}\quad $\square$ 异常
\\ \hline
\end{tabularx}

\subsection{试剂耗材}

\begin{tabularx}{\textwidth}{|l|X|X|X|}
\hline
名称 & 厂家 & 批号 & 有效期至 \\ \hline
血糖测试条 & 三诺生物传感股份有限公司 &
% #VALUE_ID: LEX-P0013-V0009
% #FIELD_VALUE: 血糖测试条批号
% #HANDWRITTEN: 444?4?06
\fieldvalue{\handwritten{444?4?06}} &
% #VALUE_ID: LEX-P0013-V0010
% #FIELD_VALUE: 血糖测试条有效期至
% #HANDWRITTEN: 2027.01.08
\fieldvalue{\handwritten{2027.01.08}}
\\ \hline
\end{tabularx}

\section{操作步骤}

\begin{tabularx}{\textwidth}{|X|X|X|X|}
\hline
\multicolumn{4}{|c|}{操作过程} \\ \hline
\multicolumn{3}{|l|}{取出测试条，取样品自然沉降后的上清液 $0.6\mu\mathrm{l}$ 加入测试条电极反应室，血糖仪自动发出“嘀”提示音，吸入样品完成，血糖仪自动进位倒计时，开始测试；} &
取上清液：
% #VALUE_ID: LEX-P0013-V0011
% #FIELD_VALUE: 取上清液体积
% #HANDWRITTEN: 0.6
\fieldvalue{\handwritten{0.6}}\,\textmu l
\\ \hline
\multicolumn{3}{|l|}{重复以上步骤 2 次，共计数 3 次。} &
每个样品分别共计数
% #VALUE_ID: LEX-P0013-V0012
% #FIELD_VALUE: 每个样品分别共计数
% #HANDWRITTEN: 3
\fieldvalue{\handwritten{3}} 次。
\\ \hline
\multicolumn{4}{|c|}{检测结果} \\ \hline
\multicolumn{4}{|c|}{葡萄糖浓度（mmol/L）} \\ \hline
第一次 & 第二次 & 第三次 & 平均值（保留一位小数） \\ \hline
% #VALUE_ID: LEX-P0013-V0013
% #FIELD_VALUE: 第一次葡萄糖浓度
% #HANDWRITTEN: 24.4
\fieldvalue{\handwritten{24.4}} &
% #VALUE_ID: LEX-P0013-V0014
% #FIELD_VALUE: 第二次葡萄糖浓度
% #HANDWRITTEN: 23.7
\fieldvalue{\handwritten{23.7}} &
% #VALUE_ID: LEX-P0013-V0015
% #FIELD_VALUE: 第三次葡萄糖浓度
% #HANDWRITTEN: 23.8
\fieldvalue{\handwritten{23.8}} &
% #VALUE_ID: LEX-P0013-V0016
% #FIELD_VALUE: 平均葡萄糖浓度
% #HANDWRITTEN: 23.9
\fieldvalue{\handwritten{23.9}}
\\ \hline
备注 &
% #VALUE_ID: LEX-P0013-V0017
% #FIELD_VALUE: 备注
% #TODO #HANDWRITTEN: /; handwritten mark is unclear
\fieldvalue{\handwritten{/}} &
\multicolumn{2}{l|}{} \\ \hline
实验人/日期 &
% #VALUE_ID: LEX-P0013-V0018
% #FIELD_VALUE: 实验人
% #TODO #HANDWRITTEN: 杨?; signature is unclear
\fieldvalue{\handwritten{杨?}}\quad
% #VALUE_ID: LEX-P0013-V0019
% #FIELD_VALUE: 实验日期
% #HANDWRITTEN: 2024.08.27
\fieldvalue{\handwritten{2024.08.27}} &
\multicolumn{2}{l|}{} \\ \hline
复核人/日期 &
% #VALUE_ID: LEX-P0013-V0020
% #FIELD_VALUE: 复核人
% #TODO #HANDWRITTEN: 孙?; signature is unclear
\fieldvalue{\handwritten{孙?}}\quad
% #VALUE_ID: LEX-P0013-V0021
% #FIELD_VALUE: 复核日期
% #HANDWRITTEN: 2024.8.29
\fieldvalue{\handwritten{2024.8.29}} &
\multicolumn{2}{l|}{} \\ \hline
\end{tabularx}

% TODO: The footer indicates “第1页 共5页”; the continuation pages are not included here.
% LEXOID_PAGE_COMPLETED: 13/60

\newpage

\begin{center}
\textbf{生物反应器细胞液检验操作记录}
\end{center}

\noindent 文件编号：REC-YZ-SOP-QC-99-006-a

\begin{center}
\begin{tabularx}{\textwidth}{|p{2.5cm}|X|p{2.5cm}|X|}
\hline
\multicolumn{4}{|c|}{密度与活率检验}\\
\hline
样品名称 &
% #VALUE_ID: LEX-P0014-V0001
% #FIELD_VALUE: 样品名称
\fieldvalue{一级生物反应器细胞悬液} &
样品规格 &
% #VALUE_ID: LEX-P0014-V0002
% #FIELD_VALUE: 样品规格
% #HANDWRITTEN: 10 ml/袋
\fieldvalue{\handwritten{10 ml/袋}}\\
\hline
样品批号 &
% #VALUE_ID: LEX-P0014-V0003
% #FIELD_VALUE: 样品批号
% #TODO #HANDWRITTEN: [批号]; handwriting is unclear
\fieldvalue{\handwritten{[批号]}} &
取样日期 &
% #VALUE_ID: LEX-P0014-V0004
% #FIELD_VALUE: 取样日期
% #HANDWRITTEN: 2024年03月27日
\fieldvalue{\handwritten{2024年03月27日}}\\
\hline
\end{tabularx}
\end{center}

\noindent\textbf{1\quad 仪器设备、试剂耗材：}

\noindent 1.1\quad 仪器设备

\begin{center}
\small
\begin{tabularx}{\textwidth}{|X|X|X|X|X|X|}
\hline
名称 & 厂家 & 型号 & 设备编号 & 校准有效期至 & 检查确认\\
\hline
% #VALUE_ID: LEX-P0014-V0005
% #FIELD_VALUE: 名称
\fieldvalue{细胞计数仪} &
% #VALUE_ID: LEX-P0014-V0006
% #FIELD_VALUE: 厂家
\fieldvalue{Nexcelom} &
% #VALUE_ID: LEX-P0014-V0007
% #FIELD_VALUE: 型号
\fieldvalue{Cellometer K2} &
% #VALUE_ID: LEX-P0014-V0008
% #FIELD_VALUE: 设备编号
\fieldvalue{1410905} &
% #VALUE_ID: LEX-P0014-V0009
% #FIELD_VALUE: 校准有效期至
% #HANDWRITTEN: 2026.03.31
\fieldvalue{\handwritten{2026.03.31}} &
% #VALUE_ID: LEX-P0014-V0010
% #FIELD_VALUE: 检查确认
% #HANDWRITTEN: 正常
\fieldvalue{\handwritten{正常}}\\
\hline
% #VALUE_ID: LEX-P0014-V0011
% #FIELD_VALUE: 名称
\fieldvalue{高速冷冻离心机} &
% #VALUE_ID: LEX-P0014-V0012
% #FIELD_VALUE: 厂家
\fieldvalue{Thermo} &
% #VALUE_ID: LEX-P0014-V0013
% #FIELD_VALUE: 型号
\fieldvalue{ST16R} &
% #VALUE_ID: LEX-P0014-V0014
% #FIELD_VALUE: 设备编号
\fieldvalue{1410810} &
% #VALUE_ID: LEX-P0014-V0015
% #FIELD_VALUE: 校准有效期至
% #HANDWRITTEN: 2026.04.12
\fieldvalue{\handwritten{2026.04.12}} &
% #VALUE_ID: LEX-P0014-V0016
% #FIELD_VALUE: 检查确认
% #HANDWRITTEN: 正常
\fieldvalue{\handwritten{正常}}\\
\hline
% #VALUE_ID: LEX-P0014-V0017
% #FIELD_VALUE: 名称
\fieldvalue{电热恒温水槽} &
% #VALUE_ID: LEX-P0014-V0018
% #FIELD_VALUE: 厂家
\fieldvalue{上海一恒} &
% #VALUE_ID: LEX-P0014-V0019
% #FIELD_VALUE: 型号
\fieldvalue{DK-8AXX} &
% #VALUE_ID: LEX-P0014-V0020
% #FIELD_VALUE: 设备编号
\fieldvalue{1411001} &
% #VALUE_ID: LEX-P0014-V0021
% #FIELD_VALUE: 校准有效期至
% #HANDWRITTEN: 2026.07.22
\fieldvalue{\handwritten{2026.07.22}} &
% #VALUE_ID: LEX-P0014-V0022
% #FIELD_VALUE: 检查确认
% #HANDWRITTEN: 正常
\fieldvalue{\handwritten{正常}}\\
\hline
\end{tabularx}
\end{center}

\noindent 1.2\quad 试剂耗材

\begin{center}
\small
\begin{tabularx}{\textwidth}{|X|X|X|X|}
\hline
名称 & 厂家 & 批号 & 有效期至\\
\hline
% #VALUE_ID: LEX-P0014-V0023
% #FIELD_VALUE: 名称
\fieldvalue{AO/PI 染液} &
% #VALUE_ID: LEX-P0014-V0024
% #FIELD_VALUE: 厂家
\fieldvalue{Nexcelom} &
% #VALUE_ID: LEX-P0014-V0025
% #FIELD_VALUE: 批号
% #HANDWRITTEN: 33?449?
\fieldvalue{\handwritten{33?449?}} &
% #VALUE_ID: LEX-P0014-V0026
% #FIELD_VALUE: 有效期至
% #HANDWRITTEN: 2026.06.11
\fieldvalue{\handwritten{2026.06.11}}\\
\hline
% #VALUE_ID: LEX-P0014-V0027
% #FIELD_VALUE: 名称
\fieldvalue{0.9\%氯化钠注射液} &
% #VALUE_ID: LEX-P0014-V0028
% #FIELD_VALUE: 厂家
\fieldvalue{石家庄四药} &
% #VALUE_ID: LEX-P0014-V0029
% #FIELD_VALUE: 批号
% #HANDWRITTEN: 7H?131?320
\fieldvalue{\handwritten{7H?131?320}} &
% #VALUE_ID: LEX-P0014-V0030
% #FIELD_VALUE: 有效期至
% #HANDWRITTEN: 2026.12.06
\fieldvalue{\handwritten{2026.12.06}}\\
\hline
% #VALUE_ID: LEX-P0014-V0031
% #FIELD_VALUE: 名称
\fieldvalue{20\%人血白蛋白} &
% #VALUE_ID: LEX-P0014-V0032
% #FIELD_VALUE: 厂家
\fieldvalue{Baxter AG} &
% #VALUE_ID: LEX-P0014-V0033
% #FIELD_VALUE: 批号
% #HANDWRITTEN: 01151684
\fieldvalue{\handwritten{01151684}} &
% #VALUE_ID: LEX-P0014-V0034
% #FIELD_VALUE: 有效期至
% #HANDWRITTEN: 2027.03.28
\fieldvalue{\handwritten{2027.03.28}}\\
\hline
\end{tabularx}
\end{center}

\noindent\textbf{2\quad 操作步骤：}

\begin{center}
\begin{tabularx}{\textwidth}{|X|}
\hline
\multicolumn{1}{|c|}{溶液准备}\\
\hline
1\%人白蛋白溶液配制：取5 ml 20\%人血白蛋白加入到95 ml的0.9\%氯化钠注射液中，混匀即得。\\
\hline
20\%人血白蛋白 \underline{\hspace{1.2cm}} ml\qquad
% #VALUE_ID: LEX-P0014-V0035
% #FIELD_VALUE: 0.9\%氯化钠注射液体积
% #HANDWRITTEN: 95
\fieldvalue{\handwritten{95}} ml\\
0.9\%氯化钠注射液 \underline{\hspace{1.2cm}} ml\qquad
% #VALUE_ID: LEX-P0014-V0036
% #FIELD_VALUE: 20\%人血白蛋白体积
% #HANDWRITTEN: 5
\fieldvalue{\handwritten{5}} ml\\
\hline
{\begin{tabularx}{\linewidth}{|X|X|X|}
\hline
名称 & 配制批号 & 有效期至\\
\hline
% #VALUE_ID: LEX-P0014-V0037
% #FIELD_VALUE: 名称
\fieldvalue{5$\times$裂解液} &
% #VALUE_ID: LEX-P0014-V0038
% #FIELD_VALUE: 配制批号
% #HANDWRITTEN: 52?071780-1
\fieldvalue{\handwritten{52?071780-1}} &
% #VALUE_ID: LEX-P0014-V0039
% #FIELD_VALUE: 有效期至
% #HANDWRITTEN: 2025.10.27
\fieldvalue{\handwritten{2025.10.27}}\\
\hline
\end{tabularx}}\\
\hline
\multicolumn{1}{|c|}{操作过程}\\
\hline
按照片50 ml离心管的刻度读取样品体积，往样品管中加入对应体积的5$\times$裂解液（每4 ml样品加入1 ml 5$\times$裂解液，按照“只进不合”原则，如26 ml样品，加入7 ml 5$\times$裂解液），充分混匀后分装入2.3的样品中，置于37℃电热恒温水槽中裂解20 min后，取出离心管摇匀，使样品充分裂解；再置于37℃电热恒温水槽中裂解10 min得到细胞悬液。\\
\hline
取出裂解好的细胞悬液300 $\times$ g离心5分钟，按照读取样品体积加入对应体积的1\%人白蛋白溶液重悬（每1.2 ml样品体积加入0.1 ml 1\%人白蛋白溶液，按照“只进不合”原则，如26 ml\\
\end{tabularx}
\end{center}

% TODO：本页底部操作步骤及表格内容在页面边界处继续，以上仅保留本页可见部分。
% LEXOID_PAGE_COMPLETED: 14/60

\newpage

\noindent
\textbf{文件编号：} REC-YS-ZOP-QC-99-006-a \hfill
\textbf{版本：} 1.4

\vspace{0.5em}

\noindent
加入 $2.2\mathrm{ml}$ 的血氧化荧光染液，反复吹打至单细胞悬液。\\
撕去计数板上下两层保护膜，将细胞悬液混匀，从中吸取 $20\mu l$ 与 $20\mu l$ AO/PI 染料混合均匀。\\
吸取
% #VALUE_ID: LEX-P0015-V0001
% #FIELD_VALUE: 吸取量
% #HANDWRITTEN: 20
\fieldvalue{\handwritten{20}}
$\mu l$ 计数。\\
取 $20\mu l$ 细胞悬液与 AO/PI 染料的混合液，注入到计数板的一个检测室中。\\
每个样品均计数
% #VALUE_ID: LEX-P0015-V0002
% #FIELD_VALUE: 计数次数
% #HANDWRITTEN: 2
\fieldvalue{\handwritten{2}}
次。

\vspace{0.5em}

将加入混合液的一端插入进样口。\\
选择 Assay 类型，使用 AO/PI 染料进行双荧光存活率检测；\\
在主界面左侧 ID 栏中输入样本批号。\\
点击主界面左下角“Preview”预览细胞图片。\\
调整仪器焦距直至观察到荧光通道下的细胞中心透亮，轮廓清晰。\\
点击“count”按钮进行计数并自动拍摄图片。\\
重复以上步骤 2 次，共计数 3 次。打印数据结果。

\vspace{0.5em}

\noindent
数据保存路径：
% #TODO #HANDWRITTEN: 路径内容难以辨认；蓝色手写字迹模糊
% #VALUE_ID: LEX-P0015-V0003
% #HANDWRITTEN: 路径内容难以辨认
\handwritten{路径内容难以辨认}

\begin{center}
\textbf{检测结果}
\end{center}

\noindent
检测结果：
% #VALUE_ID: LEX-P0015-V0004
% #FIELD_VALUE: 检测结果（二级反应峰）
% #HANDWRITTEN: ✓
\fieldvalue{\handwritten{\checkmark}}
 一级反应峰、二级反应峰、三级反应峰

\small
\renewcommand{\arraystretch}{1.35}
\begin{tabular}{|p{1.2cm}|p{2.65cm}|p{2.35cm}|p{1.7cm}|p{1.45cm}|p{1.45cm}|p{1.65cm}|}
\hline
检测次数 &
活细胞密度（个细胞/ml） &
样品密度（个细胞/ml） &
活细胞密度\\CV值（\%） &
细胞活率\\（\%） &
样品活率\\（\%） &
细胞活率\\CV值（\%）\\
\hline
第一次 &
% #VALUE_ID: LEX-P0015-V0005
% #FIELD_VALUE: 第一次活细胞密度
% #HANDWRITTEN: 7.15×10^?
$\fieldvalue{\handwritten{7.15\times10^{?}}}$ &
&
&
% #VALUE_ID: LEX-P0015-V0006
% #FIELD_VALUE: 第一次细胞活率
% #HANDWRITTEN: 93.6
\fieldvalue{\handwritten{93.6}} &
&\\
\hline
第二次 &
% #VALUE_ID: LEX-P0015-V0007
% #FIELD_VALUE: 第二次活细胞密度
% #HANDWRITTEN: 7.74×10^?
$\fieldvalue{\handwritten{7.74\times10^{?}}}$ &
% #TODO #HANDWRITTEN: 3.74×10^?；蓝色手写内容不清晰
% #VALUE_ID: LEX-P0015-V0008
% #FIELD_VALUE: 第二次样品密度
% #HANDWRITTEN: 3.74×10^?
$\fieldvalue{\handwritten{3.74\times10^{?}}}$ &
% #VALUE_ID: LEX-P0015-V0009
% #FIELD_VALUE: 第二次活细胞密度CV值
% #HANDWRITTEN: 7
\fieldvalue{\handwritten{7}} &
&
% #VALUE_ID: LEX-P0015-V0010
% #FIELD_VALUE: 第二次样品活率
% #HANDWRITTEN: 80
\fieldvalue{\handwritten{80}} &
% #VALUE_ID: LEX-P0015-V0011
% #FIELD_VALUE: 第二次细胞活率CV值
% #HANDWRITTEN: 1
\fieldvalue{\handwritten{1}}\\
\hline
第三次 &
% #VALUE_ID: LEX-P0015-V0012
% #FIELD_VALUE: 第三次活细胞密度
% #HANDWRITTEN: 6.66×10^?
$\fieldvalue{\handwritten{6.66\times10^{?}}}$ &
&
&
% #VALUE_ID: LEX-P0015-V0013
% #FIELD_VALUE: 第三次细胞活率
% #HANDWRITTEN: 96.0
\fieldvalue{\handwritten{96.0}} &
&\\
\hline
\end{tabular}
\normalsize

\vspace{0.5em}

\noindent
发病毒细胞体积（ml）：
% #TODO #HANDWRITTEN: 900?；蓝色手写数字难以辨认
% #VALUE_ID: LEX-P0015-V0014
% #FIELD_VALUE: 发病毒细胞体积
% #HANDWRITTEN: 900?
\fieldvalue{\handwritten{900?}}
\hfill
细胞数量（个细胞/缸）：
% #TODO #HANDWRITTEN: 3.??×10^8；蓝色手写内容难以辨认
% #VALUE_ID: LEX-P0015-V0015
% #FIELD_VALUE: 细胞数量
% #HANDWRITTEN: 3.??×10^8
$\fieldvalue{\handwritten{3.??\times10^{8}}}$

\vspace{0.5em}

\noindent
计数要求：样品密度为 3 倍活细胞密度结果平均取与血球计数 1\% 人白细胞稀释液质量总体，再除以生物反应器细胞液体积；细胞数量为样品密度乘以发酵罐体积。

\vspace{0.5em}

\begin{tabular}{|p{2cm}|p{13cm}|}
\hline
备注 &
\\[1.2cm]
\hline
实验人/日期 &
% #TODO #HANDWRITTEN: 签名难以辨认；蓝色手写签名模糊
% #VALUE_ID: LEX-P0015-V0016
% #FIELD_VALUE: 实验人签名
% #HANDWRITTEN: 签名难以辨认
\fieldvalue{\handwritten{签名难以辨认}}
\hfill
% #VALUE_ID: LEX-P0015-V0017
% #FIELD_VALUE: 实验日期
% #HANDWRITTEN: 2025.06.19
\fieldvalue{\handwritten{2025.06.19}}\\[0.45cm]
\hline
复核人/日期 &
% #TODO #HANDWRITTEN: 签名疑似“李秀玲”；字迹不清晰
% #VALUE_ID: LEX-P0015-V0018
% #FIELD_VALUE: 复核人签名
% #HANDWRITTEN: 李秀玲
\fieldvalue{\handwritten{李秀玲}}
\hfill
% #TODO #HANDWRITTEN: 2025.07.??；日期末位难以辨认
% #VALUE_ID: LEX-P0015-V0019
% #FIELD_VALUE: 复核日期
% #HANDWRITTEN: 2025.07.??
\fieldvalue{\handwritten{2025.07.??}}\\[0.45cm]
\hline
\end{tabular}

\noindent
\begin{center}
第 3 页 共 5 页
\end{center}
% LEXOID_PAGE_COMPLETED: 15/60

\newpage

\begin{center}
\textbf{生物反应器细胞液检测操作记录}
\end{center}

\noindent
\begin{tabularx}{\textwidth}{|l|X|l|X|}
\hline
\textbf{细胞生长状态检验} & & & \\ \hline
样品名称 &
% #VALUE_ID: LEX-P0016-V0003
% #FIELD_VALUE: 样品名称
\fieldvalue{\,\rule{1.5cm}{0.4pt}\,级生物反应器细胞液}
& 样品规格 & ml/袋 \\ \hline
样品批号 & & 取样日期 & 年\quad 月\quad 日 \\ \hline
\end{tabularx}

\section*{1.\quad 仪器设备及试剂}

\subsection*{1.1\quad 仪器设备}

\begin{tabularx}{\textwidth}{|X|X|X|X|X|}
\hline
仪器名称 & 厂家 & 型号 & 设备编号 & 检查确认 \\ \hline
% #VALUE_ID: LEX-P0016-V0004
% #FIELD_VALUE: 仪器名称
\fieldvalue{荧光显微镜}
&
% #VALUE_ID: LEX-P0016-V0005
% #FIELD_VALUE: 厂家
\fieldvalue{徕卡}
&
% #VALUE_ID: LEX-P0016-V0006
% #FIELD_VALUE: 型号
\fieldvalue{DmiI LED}
&
% #VALUE_ID: LEX-P0016-V0007
% #FIELD_VALUE: 设备编号
\fieldvalue{1431404}
& $\square$ 正常\quad $\square$ 异常 \\ \hline
\end{tabularx}

\subsection*{1.2\quad 试剂耗材}

\begin{tabularx}{\textwidth}{|X|X|X|X|X|}
\hline
名称 & 厂家 & 货号 & 批号 & 有效期至 \\ \hline
细胞活力检测试剂 &
% #VALUE_ID: LEX-P0016-V0008
% #FIELD_VALUE: 厂家
\fieldvalue{江苏凯基生物技术股份有限公司}
&
% #VALUE_ID: LEX-P0016-V0009
% #FIELD_VALUE: 货号
\fieldvalue{KGA9501-1000}
& & \\ \hline
\end{tabularx}

% #VALUE_ID: LEX-P0016-V0010
% #TODO #HANDWRITTEN: 蓝色斜线及批注，内容难以辨认；疑似手写标记
\handwritten{[难以辨认的蓝色手写批注]}

\section*{2.\quad 操作步骤}

\begin{tabularx}{\textwidth}{|X|X|}
\hline
\multicolumn{2}{|c|}{溶液配制} \\ \hline
\multicolumn{2}{|l|}{20$\times$的荧光染液配制：取\_\_\_ $\mu$l 组分A和\_\_\_ $\mu$l 组分B，混匀。} \\ \hline
\multicolumn{2}{|l|}{\hspace{2cm}20$\times$ PBS 混合均匀后配成20$\times$的荧光染液。} \\ \hline
\multicolumn{2}{|c|}{操作过程} \\ \hline
取混匀后样品 200$\mu$l 至96孔板中，取2孔。 &
样品\_\_\_\_\_\_\_\_\_\_\_\_\_\_ $\mu$l/孔，\\
& 每个样品取\_\_\_\_\_\_\_\_\_\_\_\_\_\_孔 \\ \hline
自然沉降后，吸弃上清，注意不要吸到沉淀样品； &
$\square$ 吸弃上清 \\ \hline
取20$\times$的荧光液加入到96孔板中，200$\mu$l/孔，避光静置30--60min。 &
20$\times$的荧光染液\_\_\_\_\_\_\_\_\_\_ $\mu$l/孔，\\
& 避光静置\_\_\_\_\_\_\_\_\_\_ \\ \hline
染色结束后，吸弃上清，取 PBS 加入到以上96孔板中，200$\mu$l/孔，吸弃上清，注意不要吸到沉淀样品，清洗3次，之后每孔加入100$\mu$l PBS进行拍照。 &
加入 PBS\_\_\_\_\_\_\_\_\_\_ $\mu$l/孔，\\
& 清洗\_\_\_\_\_\_\_\_\_\_次 \\ \hline
将荧光显微镜调至4$\times$物镜，软件中调节荧光参数，调节到适合的光源和像距，拍摄4张清晰图像。 &
\\ \hline
选一张清晰图像，打印并粘贴在实验记录上。 &
\\ \hline
数据保存路径 & \\ \hline
\multicolumn{2}{|c|}{检测结果} \\ \hline
\end{tabularx}
% LEXOID_PAGE_COMPLETED: 16/60

\newpage

\noindent
\begin{center}
\small
\begin{tabularx}{\textwidth}{@{}X c X@{}}
\textbf{肽皇生物科技（北京）有限公司} & \fbox{\begin{tabular}{c}原件\end{tabular}} &
\begin{tabular}{r}
文件编码：REC-YZ-SOP-QC-99-006-3\\
\textcolor{red}{版本号：1.4}
\end{tabular}
\end{tabularx}
\end{center}

\vspace{2pt}

\noindent
\begin{picture}(\textwidth,690)
  \put(0,0){\framebox(\textwidth,690){}}
  \put(0,0){\line(1,1){\textwidth}}
  \put(0,0){\line(1,0){\textwidth}}
  \put(0,0){\line(0,1){690}}
  \put(0,0){\line(1,0){\textwidth}}
  \put(0,0){\line(0,1){690}}
  % The diagonal blue line and handwritten annotation are visible in the blank form area.
  \put(270,315){\rotatebox{48}{%
    % #VALUE_ID: LEX-P0017-V0001
    % #HANDWRITTEN: illegible Chinese handwritten annotation
    \handwritten{\textcolor{blue}{\scriptsize illegible Chinese handwritten annotation}}%
  }}
\end{picture}

\vspace{-1pt}

\noindent
\begin{tabularx}{\textwidth}{|p{2.2cm}|X|}
\hline
备注 & \\ \hline
审核人/日期 & \\ \hline
复核人/日期 & \\ \hline
\end{tabularx}

\vfill

\begin{center}
\scriptsize 第 5 页 共 5 页
\end{center}
% LEXOID_PAGE_COMPLETED: 17/60

\newpage

Assay: %
% #VALUE_ID: LEX-P0018-V0001
% #FIELD_VALUE: Assay
\fieldvalue{MSC-1}

Cell Type F1: %
% #VALUE_ID: LEX-P0018-V0002
% #FIELD_VALUE: Cell Type F1
\fieldvalue{MSC (AO)}

Cell Type F2: %
% #VALUE_ID: LEX-P0018-V0003
% #FIELD_VALUE: Cell Type F2
\fieldvalue{MSC (PI)}

Sample ID: %
% #VALUE_ID: LEX-P0018-V0004
% #FIELD_VALUE: Sample ID
\fieldvalue{QC-JS-QY0402-A31Z2202508008-20250829-0H-1}

Dilution: %
% #VALUE_ID: LEX-P0018-V0005
% #FIELD_VALUE: Dilution
\fieldvalue{2.00}

\noindent\rule{\textwidth}{0.4pt}

Results:

\begin{tabularx}{\textwidth}{@{}X X X@{}}
\textbf{Count} & \textbf{Concentration} & \textbf{Mean Diameter} \\
\noalign{\smallskip}
\underline{\hspace{0.9\linewidth}} &
\underline{\hspace{0.9\linewidth}} &
\underline{\hspace{0.9\linewidth}} \\[-2pt]
Total: %
% #VALUE_ID: LEX-P0018-V0006
% #FIELD_VALUE: Total count
\fieldvalue{216}
&
%
% #VALUE_ID: LEX-P0018-V0007
% #FIELD_VALUE: Total concentration
\fieldvalue{$7.64\times 10^{5}$ cells/mL}
&
%
% #VALUE_ID: LEX-P0018-V0008
% #FIELD_VALUE: Total mean diameter
\fieldvalue{15.0 micron}
\\
Live: %
% #VALUE_ID: LEX-P0018-V0009
% #FIELD_VALUE: Live count
\fieldvalue{202}
&
%
% #VALUE_ID: LEX-P0018-V0010
% #FIELD_VALUE: Live concentration
\fieldvalue{$7.15\times 10^{5}$ cells/mL}
&
%
% #VALUE_ID: LEX-P0018-V0011
% #FIELD_VALUE: Live mean diameter
\fieldvalue{15.0 microns}
\\
Dead: %
% #VALUE_ID: LEX-P0018-V0012
% #FIELD_VALUE: Dead count
\fieldvalue{14}
&
%
% #VALUE_ID: LEX-P0018-V0013
% #FIELD_VALUE: Dead concentration
\fieldvalue{$4.90\times 10^{4}$ cells/mL}
&
%
% #VALUE_ID: LEX-P0018-V0014
% #FIELD_VALUE: Dead mean diameter
\fieldvalue{14.2 microns}
\\
Viability: %
% #VALUE_ID: LEX-P0018-V0015
% #FIELD_VALUE: Viability
\fieldvalue{93.6\%}
&
&
\end{tabularx}

\vspace{1em}

% TODO: Four microscopy images are visible on this page; no source filenames are available.
\noindent
\begin{minipage}[t]{0.47\textwidth}
\centering
\fbox{\parbox[c][0.23\textheight][c]{0.95\linewidth}{\centering TODO: microscopy image}}
\end{minipage}
\hfill
\begin{minipage}[t]{0.47\textwidth}
\centering
\fbox{\parbox[c][0.23\textheight][c]{0.95\linewidth}{\centering TODO: microscopy image}}
\end{minipage}

\vspace{1em}

\noindent
\begin{minipage}[t]{0.47\textwidth}
\centering
\fbox{\parbox[c][0.23\textheight][c]{0.95\linewidth}{\centering TODO: microscopy image}}
\end{minipage}
\hfill
\begin{minipage}[t]{0.47\textwidth}
\centering
\fbox{\parbox[c][0.23\textheight][c]{0.95\linewidth}{\centering TODO: microscopy image}}
\end{minipage}

\vspace{1em}

操作人/审核人：%
% #VALUE_ID: LEX-P0018-V0016
% #FIELD_VALUE: 操作人/审核人
% TODO #HANDWRITTEN: [illegible signature]; handwritten signature is not clearly legible
\fieldvalue{\handwritten{[illegible signature]}}

\hfill

日期：%
% #VALUE_ID: LEX-P0018-V0017
% #FIELD_VALUE: 日期
% #HANDWRITTEN: 2025.8.29
\fieldvalue{\handwritten{2025.8.29}}
% LEXOID_PAGE_COMPLETED: 18/60

\newpage

\noindent
\textbf{Assay:}
% #VALUE_ID: LEX-P0019-V0001
% #FIELD_VALUE: Assay
\fieldvalue{MSC-1}

\vspace{0.8em}

\noindent
\textbf{Cell Type F1:}
% #VALUE_ID: LEX-P0019-V0002
% #FIELD_VALUE: Cell Type F1
\fieldvalue{MSC (AO)}

\noindent
\textbf{Cell Type F2:}
% #VALUE_ID: LEX-P0019-V0003
% #FIELD_VALUE: Cell Type F2
\fieldvalue{MSC (PI)}

\vspace{0.8em}

\noindent
\textbf{Sample ID:}
% #VALUE_ID: LEX-P0019-V0004
% #FIELD_VALUE: Sample ID
\fieldvalue{QC-JS-QY0402-A31Z2202508008-20250829-0H-2}

\noindent
\textbf{Dilution:}
% #VALUE_ID: LEX-P0019-V0005
% #FIELD_VALUE: Dilution
\fieldvalue{2.00}

\vspace{0.8em}
\hrule
\vspace{0.4em}

\noindent
\textbf{Results:}

\vspace{0.5em}

\renewcommand{\arraystretch}{1.15}
\begin{tabularx}{\textwidth}{@{}>{\RaggedRight\arraybackslash}X
                                >{\RaggedRight\arraybackslash}X
                                >{\RaggedRight\arraybackslash}X@{}}
\textbf{Count} & \textbf{Concentration} & \textbf{Mean Diameter} \\
\texttt{===============} & \texttt{===============} & \texttt{===============} \\

Total:
% #VALUE_ID: LEX-P0019-V0006
% #FIELD_VALUE: Total count
\fieldvalue{220}
&
% #VALUE_ID: LEX-P0019-V0007
% #FIELD_VALUE: Total concentration
\fieldvalue{$7.77\times 10^{5}$ cells/mL}
&
% #VALUE_ID: LEX-P0019-V0008
% #FIELD_VALUE: Total mean diameter
\fieldvalue{14.3 micron}
\\

Live:
% #VALUE_ID: LEX-P0019-V0009
% #FIELD_VALUE: Live count
\fieldvalue{214}
&
% #VALUE_ID: LEX-P0019-V0010
% #FIELD_VALUE: Live concentration
\fieldvalue{$7.56\times 10^{5}$ cells/mL}
&
% #VALUE_ID: LEX-P0019-V0011
% #FIELD_VALUE: Live mean diameter
\fieldvalue{14.4 microns}
\\

Dead:
% #VALUE_ID: LEX-P0019-V0012
% #FIELD_VALUE: Dead count
\fieldvalue{6}
&
% #VALUE_ID: LEX-P0019-V0013
% #FIELD_VALUE: Dead concentration
\fieldvalue{$2.10\times 10^{4}$ cells/mL}
&
% #VALUE_ID: LEX-P0019-V0014
% #FIELD_VALUE: Dead mean diameter
\fieldvalue{10.4 microns}
\\

Viability:
% #VALUE_ID: LEX-P0019-V0015
% #FIELD_VALUE: Viability
\fieldvalue{97.3\%}
&
&
\end{tabularx}

\vspace{2em}

\begin{center}
\begin{tabular}{cc}
% TODO: Four microscopy images are visible on this page; source filenames are unavailable.
\fbox{\rule{0pt}{0.22\textheight}\rule{0.42\textwidth}{0pt}}
&
\fbox{\rule{0pt}{0.22\textheight}\rule{0.42\textwidth}{0pt}}
\\[1em]
\fbox{\rule{0pt}{0.22\textheight}\rule{0.42\textwidth}{0pt}}
&
\fbox{\rule{0pt}{0.22\textheight}\rule{0.42\textwidth}{0pt}}
\end{tabular}
\end{center}

\vfill

\noindent
操作人/审核人：

% #VALUE_ID: LEX-P0019-V0016
% #FIELD_VALUE: 操作人/审核人
% #TODO #HANDWRITTEN: illegible signature; handwritten signature cannot be reliably transcribed
\fieldvalue{\handwritten{签名}}

\hfill
日期：

% #VALUE_ID: LEX-P0019-V0017
% #FIELD_VALUE: 日期
% #HANDWRITTEN: 2024.8.9
\fieldvalue{\handwritten{2024.8.9}}
% LEXOID_PAGE_COMPLETED: 19/60

\newpage

Assay: %
% #VALUE_ID: LEX-P0020-V0001
% #FIELD_VALUE: Assay
\fieldvalue{MSC-1}

\vspace{0.5em}

Cell Type F1: %
% #VALUE_ID: LEX-P0020-V0002
% #FIELD_VALUE: Cell Type F1
\fieldvalue{MSC (AO)}

Cell Type F2: %
% #VALUE_ID: LEX-P0020-V0003
% #FIELD_VALUE: Cell Type F2
\fieldvalue{MSC (PI)}

\vspace{0.5em}

Sample ID: %
% #VALUE_ID: LEX-P0020-V0004
% #FIELD_VALUE: Sample ID
\fieldvalue{QC-JS-QY0402-A31Z2202508008-20250829-0H-3}

Dilution: %
% #VALUE_ID: LEX-P0020-V0005
% #FIELD_VALUE: Dilution
\fieldvalue{2.00}

\vspace{1em}
\hrule
\vspace{0.5em}

Results:

\vspace{0.5em}

\begin{tabularx}{\textwidth}{@{}lXX@{}}
\textbf{Count} & \textbf{Concentration} & \textbf{Mean Diameter} \\
\hline
Total: %
% #VALUE_ID: LEX-P0020-V0006
% #FIELD_VALUE: Total count
\fieldvalue{196}
&
% #VALUE_ID: LEX-P0020-V0007
% #FIELD_VALUE: Total concentration
\fieldvalue{$6.94\times10^{5}$ cells/mL}
&
% #VALUE_ID: LEX-P0020-V0008
% #FIELD_VALUE: Total mean diameter
\fieldvalue{15.3 micron}
\\
Live: %
% #VALUE_ID: LEX-P0020-V0009
% #FIELD_VALUE: Live count
\fieldvalue{188}
&
% #VALUE_ID: LEX-P0020-V0010
% #FIELD_VALUE: Live concentration
\fieldvalue{$6.66\times10^{5}$ cells/mL}
&
% #VALUE_ID: LEX-P0020-V0011
% #FIELD_VALUE: Live mean diameter
\fieldvalue{15.4 microns}
\\
Dead: %
% #VALUE_ID: LEX-P0020-V0012
% #FIELD_VALUE: Dead count
\fieldvalue{8}
&
% #VALUE_ID: LEX-P0020-V0013
% #FIELD_VALUE: Dead concentration
\fieldvalue{$2.80\times10^{4}$ cells/mL}
&
% #VALUE_ID: LEX-P0020-V0014
% #FIELD_VALUE: Dead mean diameter
\fieldvalue{11.4 microns}
\\
Viability: %
% #VALUE_ID: LEX-P0020-V0015
% #FIELD_VALUE: Viability
\fieldvalue{96.0\%}
&
&
\end{tabularx}

\vspace{2em}

\begin{center}
\begin{tabular}{cc}
% TODO: Image visible on page; no source filename is available.
\fbox{\parbox[c][4.4cm][c]{0.39\textwidth}{\centering Image placeholder}}
&
% TODO: Image visible on page; no source filename is available.
\fbox{\parbox[c][4.4cm][c]{0.39\textwidth}{\centering Image placeholder}}
\\[1em]
% TODO: Image visible on page; no source filename is available.
\fbox{\parbox[c][4.4cm][c]{0.39\textwidth}{\centering Image placeholder}}
&
% TODO: Image visible on page; no source filename is available.
\fbox{\parbox[c][4.4cm][c]{0.39\textwidth}{\centering Image placeholder}}
\end{tabular}
\end{center}

\vfill

操作人/审核人：%
% #VALUE_ID: LEX-P0020-V0016
% #FIELD_VALUE: 操作人/审核人
% TODO #HANDWRITTEN: illegible signature; handwritten signature is unclear
\fieldvalue{\handwritten{illegible signature}}
\hfill
日期：%
% #VALUE_ID: LEX-P0020-V0017
% #FIELD_VALUE: 日期
% TODO #HANDWRITTEN: 24/8/??; handwritten date is partially illegible
\fieldvalue{\handwritten{24/8/??}}
% LEXOID_PAGE_COMPLETED: 20/60

\newpage

\begin{center}
\textbf{原件请验单}
\end{center}

\noindent
文件编码：REC-YS-ZMP-QC-01-01-d \hfill 版本号：1.4

\vspace{0.5em}

\renewcommand{\arraystretch}{1.45}
\small
\begin{tabularx}{\textwidth}{|p{2.4cm}|X|p{2.8cm}|X|}
\hline
样品名称 &
% #VALUE_ID: LEX-P0021-V0001
% #FIELD_VALUE: 样品名称
% #TODO #HANDWRITTEN: 低钽钨硅抛光液; handwriting is partially unclear
\fieldvalue{\handwritten{低钽钨硅抛光液}} &
样品代码 &
% #VALUE_ID: LEX-P0021-V0002
% #FIELD_VALUE: 样品代码
% #TODO #HANDWRITTEN: 0863; final digit is unclear
\fieldvalue{\handwritten{0863}}
\\ \hline

样品批号 &
% #VALUE_ID: LEX-P0021-V0003
% #FIELD_VALUE: 样品批号
% #TODO #HANDWRITTEN: A3212250808; handwriting is partially unclear
\fieldvalue{\handwritten{A3212250808}} &
样品规格/数量 &
% #VALUE_ID: LEX-P0021-V0004
% #FIELD_VALUE: 样品规格/数量
% #TODO #HANDWRITTEN: 4kg/袋×1; handwriting is partially unclear
$\fieldvalue{\handwritten{4kg/袋\times 1}}$
\\ \hline

\multirow{2}{*}{样品类别} &
% #VALUE_ID: LEX-P0021-V0005
% #FIELD_VALUE: 样品类别
\fieldvalue{\checkmark\ 生产过程中的检验样品}
\quad $\square$ 待包装产品
\quad $\square$ 成品
&
\multicolumn{1}{l|}{}
\\[-0.2em]
&
$\square$ 其他（\hspace{2em}）
&
\multicolumn{1}{l|}{}
\\ \hline

检验类型 &
% #VALUE_ID: LEX-P0021-V0006
% #FIELD_VALUE: 检验类型
\fieldvalue{\checkmark\ 初验}
\quad $\square$ 复验
\quad $\square$ 退货
\quad $\square$ 稳定性
\quad $\square$ 留样
&
\multicolumn{1}{l|}{}
\\ \hline

请验目的 &
% #VALUE_ID: LEX-P0021-V0007
% #FIELD_VALUE: 请验目的
\fieldvalue{\checkmark\ 全检}
\quad $\square$ 留样
\quad $\square$ 其他（\hspace{2em}）
&
\multicolumn{1}{l|}{}
\\ \hline

请验部门 &
% #VALUE_ID: LEX-P0021-V0008
% #FIELD_VALUE: 请验部门
\fieldvalue{\checkmark\ 生产部}
\quad $\square$ 物料管理部
\quad $\square$ 其他（\hspace{2em}）
&
\multicolumn{1}{l|}{}
\\ \hline

样品情况 &
% #VALUE_ID: LEX-P0021-V0009
% #FIELD_VALUE: 样品情况
\fieldvalue{\checkmark\ 随请验单一起提供}
\quad $\square$ 未随请验单一起提供
&
\multicolumn{1}{l|}{}
\\ \hline

请验人 &
% #VALUE_ID: LEX-P0021-V0010
% #FIELD_VALUE: 请验人
% #HANDWRITTEN: 张宁
\fieldvalue{\handwritten{张宁}} &
请验日期 &
% #VALUE_ID: LEX-P0021-V0011
% #FIELD_VALUE: 请验日期
% #HANDWRITTEN: 2025-3-24
\fieldvalue{\handwritten{2025-3-24}}
\\ \hline

收验人 &
% #VALUE_ID: LEX-P0021-V0012
% #FIELD_VALUE: 收验人
% #TODO #HANDWRITTEN: 陈祎; handwriting is unclear
\fieldvalue{\handwritten{陈祎}} &
收验日期 &
% #VALUE_ID: LEX-P0021-V0013
% #FIELD_VALUE: 收验日期
% #TODO #HANDWRITTEN: 2025-3-20; final digit is unclear
\fieldvalue{\handwritten{2025-3-20}}
\\ \hline

备注 &
% #VALUE_ID: LEX-P0021-V0014
% #FIELD_VALUE: 备注
% #TODO #HANDWRITTEN: 12 7h; handwriting is illegible
\fieldvalue{\handwritten{12\ 7h}} &
\multicolumn{2}{l|}{}
\\ \hline
\end{tabularx}
% LEXOID_PAGE_COMPLETED: 21/60

\newpage

\begin{center}
文件编码：REC-YZ-SMP-QC-01-004-b\qquad 版本号：1.4

\vspace{0.5em}
{\large\textbf{取样记录}}
\end{center}

\renewcommand{\arraystretch}{1.7}
\small
\begin{tabularx}{\textwidth}{|p{2.0cm}|X|p{2.0cm}|X|}
\hline
名称 &
% #VALUE_ID: LEX-P0022-V0001
% #FIELD_VALUE: 名称
\fieldvalue{一级生物反应器细胞液} &
代码 &
% #VALUE_ID: LEX-P0022-V0002
% #FIELD_VALUE: 代码
\fieldvalue{QY0401}
\\ \hline

规格 &
% #VALUE_ID: LEX-P0022-V0003
% #FIELD_VALUE: 规格
\fieldvalue{约20ml/袋} &
总数量 &
% #VALUE_ID: LEX-P0022-V0004
% #FIELD_VALUE: 总数量
\fieldvalue{1袋}
\\ \hline

批号 &
% #VALUE_ID: LEX-P0022-V0005
% #FIELD_VALUE: 批号
\fieldvalue{A31Z2202508008} &
进厂批号 &
% #VALUE_ID: LEX-P0022-V0006
% #FIELD_VALUE: 进厂批号
\fieldvalue{/}
\\ \hline

\multicolumn{1}{|c|}{供货单位/厂家} &
\multicolumn{3}{l|}{
% #VALUE_ID: LEX-P0022-V0007
% #FIELD_VALUE: 供货单位/厂家
\fieldvalue{/}
}
\\ \hline

取样目的 &
\multicolumn{3}{l|}{
% #VALUE_ID: LEX-P0022-V0008
% #FIELD_VALUE: 取样目的
\fieldvalue{[X]正常检验}\quad
$\square$复检\quad
$\square$稳定性
\newline
\hspace*{5.1em}$\square$退货检验\quad
$\square$留样\quad
$\square$其他：\underline{\hspace{3em}}
}
\\ \hline

取样地点 &
\multicolumn{3}{l|}{
% #VALUE_ID: LEX-P0022-V0009
% #FIELD_VALUE: 取样地点
\fieldvalue{[X]产品出口}\quad
$\square$物料库\quad
$\square$细胞库\quad
$\square$包装间
}
\\ \hline

贮存条件 &
\multicolumn{3}{l|}{
% #VALUE_ID: LEX-P0022-V0010
% #FIELD_VALUE: 贮存条件
\fieldvalue{[X]2--8℃}\quad
$\square$-20℃\quad
$\square$-80℃\quad
$\square$常温\quad
$\square$其他：\underline{\hspace{3em}}
}
\\ \hline

盛装容器 &
\multicolumn{3}{l|}{
% #VALUE_ID: LEX-P0022-V0011
% #FIELD_VALUE: 盛装容器
\fieldvalue{[X]保温箱}\quad
$\square$泡沫箱\quad
$\square$手提篮\quad
$\square$其他：\underline{\hspace{3em}}
}
\\ \hline

外观复核 &
\multicolumn{3}{l|}{
% #VALUE_ID: LEX-P0022-V0012
% #FIELD_VALUE: 外观复核
\fieldvalue{[X]合格}\quad
$\square$不合格
}
\\ \hline

取样量 &
% #VALUE_ID: LEX-P0022-V0013
% #FIELD_VALUE: 取样量
\fieldvalue{1袋} &
检验量 &
% #VALUE_ID: LEX-P0022-V0014
% #FIELD_VALUE: 检验量
\fieldvalue{1袋}
\\ \hline

留样量 &
% #VALUE_ID: LEX-P0022-V0015
% #FIELD_VALUE: 留样量
\fieldvalue{/} &
\multicolumn{2}{l|}{}
\\ \hline

取样人/日期 &
\multicolumn{1}{l|}{
% #VALUE_ID: LEX-P0022-V0016
% #FIELD_VALUE: 取样人
% #TODO #HANDWRITTEN: 杨敏; 姓名笔迹不清
\fieldvalue{\handwritten{杨敏}}
\quad
% #VALUE_ID: LEX-P0022-V0017
% #FIELD_VALUE: 取样日期
% #HANDWRITTEN: 2025.8.7
\fieldvalue{\handwritten{2025.8.7}}
} &
复核人/日期 &
\multicolumn{1}{l|}{
% #VALUE_ID: LEX-P0022-V0018
% #FIELD_VALUE: 复核人
% #TODO #HANDWRITTEN: 李春霞; 姓名笔迹不清
\fieldvalue{\handwritten{李春霞}}
\quad
% #VALUE_ID: LEX-P0022-V0019
% #FIELD_VALUE: 复核日期
% #HANDWRITTEN: 2025.8.9
\fieldvalue{\handwritten{2025.8.9}}
}
\\ \hline

备注 &
\multicolumn{3}{l|}{
% #VALUE_ID: LEX-P0022-V0020
% #FIELD_VALUE: 备注
% #HANDWRITTEN: 1.9L
\fieldvalue{\handwritten{1.9L}}
\qquad
% #VALUE_ID: LEX-P0022-V0021
% #FIELD_VALUE: 备注
% #HANDWRITTEN: 7L
\fieldvalue{\handwritten{7L}}
}
\\ \hline
\end{tabularx}

\vfill
\begin{center}
第 1 页 共 1 页
\end{center}
% LEXOID_PAGE_COMPLETED: 22/60

\newpage

\begin{center}
\textbf{生物反应器细胞液检验操作记录}
\end{center}

\noindent
文件编号：RBC-YZ-SOP-QC-99-006-a

\begin{tabularx}{\textwidth}{|l|X|l|X|}
\hline
样品名称 &
% #VALUE_ID: LEX-P0023-V0001
% #FIELD_VALUE: 样品名称
% #HANDWRITTEN: 二级生物反应器细胞液
\fieldvalue{\handwritten{二级生物反应器细胞液}} &
样品规格 &
% #VALUE_ID: LEX-P0023-V0002
% #FIELD_VALUE: 样品规格
% #TODO #HANDWRITTEN: 162.0; 字迹部分不清
\fieldvalue{\handwritten{162.0}} ml/袋
\\
\hline
样品批号 &
% #VALUE_ID: LEX-P0023-V0003
% #FIELD_VALUE: 样品批号
% #TODO #HANDWRITTEN: AB\ldots; 手写批号难以辨认
\fieldvalue{\handwritten{AB\ldots}} &
取样日期 &
% #VALUE_ID: LEX-P0023-V0004
% #FIELD_VALUE: 取样日期
% #TODO #HANDWRITTEN: 2024年08月9日; 日期末位字迹不清
\fieldvalue{\handwritten{2024年08月9日}}
\\
\hline
\end{tabularx}

\section*{1\quad 仪器设备、试剂耗材}

\subsection*{1.1\quad 仪器设备}

\begin{tabularx}{\textwidth}{|l|X|l|l|l|}
\hline
名称 & 厂家 & 型号 & 设备编号 & 检查确认 \\
\hline
% #VALUE_ID: LEX-P0023-V0005
% #FIELD_VALUE: 仪器名称
\fieldvalue{血球仪} &
% #VALUE_ID: LEX-P0023-V0006
% #FIELD_VALUE: 仪器厂家
\fieldvalue{三诺生物传感股份有限公司} &
% #VALUE_ID: LEX-P0023-V0007
% #FIELD_VALUE: 仪器型号
\fieldvalue{安捷？} &
% #VALUE_ID: LEX-P0023-V0008
% #FIELD_VALUE: 设备编号
% #TODO #HANDWRITTEN: 设备编号约为 173049?; 手写字符不清
\fieldvalue{\handwritten{173049?}} &
% #VALUE_ID: LEX-P0023-V0009
% #FIELD_VALUE: 检查确认
% #HANDWRITTEN: 勾选“正常”
\fieldvalue{\handwritten{\checkmark\ 正常}} \quad 异常
\\
\hline
% #VALUE_ID: LEX-P0023-V0010
% #FIELD_VALUE: 仪器名称
\fieldvalue{移液器} &
% #VALUE_ID: LEX-P0023-V0011
% #FIELD_VALUE: 仪器厂家
\fieldvalue{Eppendorf Corporate} &
% #VALUE_ID: LEX-P0023-V0012
% #FIELD_VALUE: 仪器型号
\fieldvalue{0.1$\sim$2.5$\mu$l} &
% #VALUE_ID: LEX-P0023-V0013
% #FIELD_VALUE: 设备编号
% #TODO #HANDWRITTEN: 约为 173049?; 手写字符不清
\fieldvalue{\handwritten{173049?}} &
% #VALUE_ID: LEX-P0023-V0014
% #FIELD_VALUE: 检查确认
% #HANDWRITTEN: 勾选“正常”
\fieldvalue{\handwritten{\checkmark\ 正常}} \quad 异常
\\
\hline
\end{tabularx}

\subsection*{1.2\quad 试剂耗材}

\begin{tabularx}{\textwidth}{|l|X|l|l|}
\hline
名称 & 厂家 & 批号 & 有效期至 \\
\hline
% #VALUE_ID: LEX-P0023-V0015
% #FIELD_VALUE: 试剂耗材名称
\fieldvalue{血细胞测试条} &
% #VALUE_ID: LEX-P0023-V0016
% #FIELD_VALUE: 试剂耗材厂家
\fieldvalue{三诺生物传感股份有限公司} &
% #VALUE_ID: LEX-P0023-V0017
% #FIELD_VALUE: 批号
% #TODO #HANDWRITTEN: 约为 444?1706; 手写字符不清
\fieldvalue{\handwritten{444?1706}} &
% #VALUE_ID: LEX-P0023-V0018
% #FIELD_VALUE: 有效期至
% #HANDWRITTEN: 2027.1.30
\fieldvalue{\handwritten{2027.1.30}}
\\
\hline
\end{tabularx}

\section*{2\quad 操作步骤}

\begin{tabularx}{\textwidth}{|X|X|X|X|}
\hline
\multicolumn{4}{|c|}{操作过程} \\
\hline
\multicolumn{3}{|l|}{取出测试条，取样品自然沉降后的上清液 $0.6\,\mu l$ 加入测试条电极反应室，血糖仪自动发出“嘀”提示音，吸入样品完成，血糖仪自动进入倒计时，开始测试；} &
取上清液：
% #VALUE_ID: LEX-P0023-V0019
% #FIELD_VALUE: 取上清液体积
% #HANDWRITTEN: 16
\fieldvalue{\handwritten{16}}\,$\mu$l
\\
\hline
\multicolumn{3}{|l|}{重复以上步骤 $2$ 次，共计数 $3$ 次。} &
每个样品分别共计数
% #VALUE_ID: LEX-P0023-V0020
% #FIELD_VALUE: 每个样品计数次数
% #HANDWRITTEN: 3
\fieldvalue{\handwritten{3}} 次。
\\
\hline
\multicolumn{4}{|c|}{检测结果} \\
\hline
\multicolumn{4}{|c|}{葡萄糖浓度（mmol/L）} \\
\hline
第一次 & 第二次 & 第三次 & 平均值（保留一位小数） \\
\hline
% #VALUE_ID: LEX-P0023-V0021
% #FIELD_VALUE: 第一次检测结果
% #HANDWRITTEN: 17.9
\fieldvalue{\handwritten{17.9}} &
% #VALUE_ID: LEX-P0023-V0022
% #FIELD_VALUE: 第二次检测结果
% #HANDWRITTEN: 18.1
\fieldvalue{\handwritten{18.1}} &
% #VALUE_ID: LEX-P0023-V0023
% #FIELD_VALUE: 第三次检测结果
% #HANDWRITTEN: 18.0
\fieldvalue{\handwritten{18.0}} &
% #VALUE_ID: LEX-P0023-V0024
% #FIELD_VALUE: 平均值
% #HANDWRITTEN: 18.0
\fieldvalue{\handwritten{18.0}}
\\
\hline
备注 &
% #VALUE_ID: LEX-P0023-V0025
% #FIELD_VALUE: 备注
% #HANDWRITTEN: 勾
\fieldvalue{\handwritten{\checkmark}}
\\
\hline
实验人/日期 &
% #VALUE_ID: LEX-P0023-V0026
% #FIELD_VALUE: 实验人
% #TODO #HANDWRITTEN: 姓名难以辨认
\fieldvalue{\handwritten{姓名不清}} \qquad
% #VALUE_ID: LEX-P0023-V0027
% #FIELD_VALUE: 实验日期
% #TODO #HANDWRITTEN: 2024.08.21; 前部姓名与日期相连，日期基本可辨
\fieldvalue{\handwritten{2024.08.21}}
\\
\hline
复核人/日期 &
% #VALUE_ID: LEX-P0023-V0028
% #FIELD_VALUE: 复核人
% #TODO #HANDWRITTEN: 姓名难以辨认
\fieldvalue{\handwritten{姓名不清}} \qquad
% #VALUE_ID: LEX-P0023-V0029
% #FIELD_VALUE: 复核日期
% #HANDWRITTEN: 2024.08.29
\fieldvalue{\handwritten{2024.08.29}}
\\
\hline
\end{tabularx}

\begin{center}
第 1 页\quad 共 5 页
\end{center}
% LEXOID_PAGE_COMPLETED: 23/60

\newpage

\section{生物反应器细胞检验操作记录}

\noindent
\textbf{文件编号：} REC-YZ-SOP-QC-99-006-a

\begin{tabularx}{\textwidth}{|l|X|l|X|}
\hline
\multicolumn{4}{|c|}{\textbf{细胞生物}}\\
\hline
原件 & & & \\
\hline
样品名称 &
% #VALUE_ID: LEX-P0024-V0001
% #FIELD_VALUE: 样品名称
\fieldvalue{一级生物反应器细胞液}
& 样品规格 &
% #VALUE_ID: LEX-P0024-V0002
% #FIELD_VALUE: 样品规格
% #HANDWRITTEN: 15?0
\fieldvalue{\handwritten{15?0}}\ \text{ml/袋}\\
\hline
样品批号 &
% #VALUE_ID: LEX-P0024-V0003
% #FIELD_VALUE: 样品批号
% #HANDWRITTEN: 202?年?月?日
\fieldvalue{\handwritten{202?年?月?日}}
& 取样日期 &
% #VALUE_ID: LEX-P0024-V0004
% #FIELD_VALUE: 取样日期
% #HANDWRITTEN: 202?年?月?日
\fieldvalue{\handwritten{202?年?月?日}}\\
\hline
\end{tabularx}

\subsection{1\quad 仪器设备、试剂耗材}

\subsubsection{1.1\quad 仪器设备}

\begin{tabularx}{\textwidth}{|X|X|X|X|X|X|}
\hline
名称 & 厂家 & 型号 & 设备编号 & 校准有效期至 & 检查确认\\
\hline
% #VALUE_ID: LEX-P0024-V0005
% #FIELD_VALUE: 名称
\fieldvalue{细胞计数仪}
&
% #VALUE_ID: LEX-P0024-V0006
% #FIELD_VALUE: 厂家
\fieldvalue{Nexcelom}
&
% #VALUE_ID: LEX-P0024-V0007
% #FIELD_VALUE: 型号
\fieldvalue{Cellometer K2}
&
% #VALUE_ID: LEX-P0024-V0008
% #FIELD_VALUE: 设备编号
\fieldvalue{1410905}
&
% #VALUE_ID: LEX-P0024-V0009
% #FIELD_VALUE: 校准有效期至
% #TODO #HANDWRITTEN: 2026.08.?; 末位不清晰
\fieldvalue{\handwritten{2026.08.?}}
&
% #VALUE_ID: LEX-P0024-V0010
% #FIELD_VALUE: 检查确认
\fieldvalue{正常（已勾选）}\\
\hline
% #VALUE_ID: LEX-P0024-V0011
% #FIELD_VALUE: 名称
\fieldvalue{高速冷冻离心机}
&
% #VALUE_ID: LEX-P0024-V0012
% #FIELD_VALUE: 厂家
\fieldvalue{Thermo}
&
% #VALUE_ID: LEX-P0024-V0013
% #FIELD_VALUE: 型号
\fieldvalue{ST16R}
&
% #VALUE_ID: LEX-P0024-V0014
% #FIELD_VALUE: 设备编号
\fieldvalue{1410810}
&
% #VALUE_ID: LEX-P0024-V0015
% #FIELD_VALUE: 校准有效期至
% #TODO #HANDWRITTEN: 2026.10.?; 末位不清晰
\fieldvalue{\handwritten{2026.10.?}}
&
% #VALUE_ID: LEX-P0024-V0016
% #FIELD_VALUE: 检查确认
\fieldvalue{正常（已勾选）}\\
\hline
% #VALUE_ID: LEX-P0024-V0017
% #FIELD_VALUE: 名称
\fieldvalue{电热恒温水浴箱}
&
% #VALUE_ID: LEX-P0024-V0018
% #FIELD_VALUE: 厂家
\fieldvalue{上海一恒}
&
% #VALUE_ID: LEX-P0024-V0019
% #FIELD_VALUE: 型号
\fieldvalue{DK-8AXX}
&
% #VALUE_ID: LEX-P0024-V0020
% #FIELD_VALUE: 设备编号
\fieldvalue{1411001}
&
% #VALUE_ID: LEX-P0024-V0021
% #FIELD_VALUE: 校准有效期至
% #TODO #HANDWRITTEN: 2026.07.12; handwriting partly unclear
\fieldvalue{\handwritten{2026.07.12}}
&
% #VALUE_ID: LEX-P0024-V0022
% #FIELD_VALUE: 检查确认
\fieldvalue{正常（已勾选）}\\
\hline
\end{tabularx}

\subsubsection{1.2\quad 试剂耗材}

\begin{tabularx}{\textwidth}{|X|X|X|X|}
\hline
名称 & 厂家 & 批号 & 有效期至\\
\hline
% #VALUE_ID: LEX-P0024-V0023
% #FIELD_VALUE: 名称
\fieldvalue{AO/PI 染液}
&
% #VALUE_ID: LEX-P0024-V0024
% #FIELD_VALUE: 厂家
\fieldvalue{Nexcelom}
&
% #VALUE_ID: LEX-P0024-V0025
% #FIELD_VALUE: 批号
% #HANDWRITTEN: 33?4?9?
\fieldvalue{\handwritten{33?4?9?}}
&
% #VALUE_ID: LEX-P0024-V0026
% #FIELD_VALUE: 有效期至
% #HANDWRITTEN: 2026.06.22
\fieldvalue{\handwritten{2026.06.22}}\\
\hline
% #VALUE_ID: LEX-P0024-V0027
% #FIELD_VALUE: 名称
\fieldvalue{0.9\%氯化钠注射液}
&
% #VALUE_ID: LEX-P0024-V0028
% #FIELD_VALUE: 厂家
\fieldvalue{石家庄四药}
&
% #VALUE_ID: LEX-P0024-V0029
% #FIELD_VALUE: 批号
% #HANDWRITTEN: 241?073?30
\fieldvalue{\handwritten{241?073?30}}
&
% #VALUE_ID: LEX-P0024-V0030
% #FIELD_VALUE: 有效期至
% #HANDWRITTEN: 2026.10.08
\fieldvalue{\handwritten{2026.10.08}}\\
\hline
% #VALUE_ID: LEX-P0024-V0031
% #FIELD_VALUE: 名称
\fieldvalue{20\%人血白蛋白}
&
% #VALUE_ID: LEX-P0024-V0032
% #FIELD_VALUE: 厂家
\fieldvalue{Baxter AG}
&
% #VALUE_ID: LEX-P0024-V0033
% #FIELD_VALUE: 批号
% #HANDWRITTEN: 4113/0?4
\fieldvalue{\handwritten{4113/0?4}}
&
% #VALUE_ID: LEX-P0024-V0034
% #FIELD_VALUE: 有效期至
% #HANDWRITTEN: 2027.02.28
\fieldvalue{\handwritten{2027.02.28}}\\
\hline
\end{tabularx}

\section{2\quad 操作步骤}

\subsection{溶液准备}

\noindent
1\%人白蛋白将溶液配制：取5 ml 20\%人血白蛋白加入到95 ml 的0.9\%氯化钠注射液中，混匀即可。

\noindent
20\%人血白蛋白
% #VALUE_ID: LEX-P0024-V0035
% #FIELD_VALUE: 20\%人血白蛋白用量
% #HANDWRITTEN: 1
\fieldvalue{\handwritten{1}}\ \text{ml}

\noindent
0.9\%氯化钠注射液
% #VALUE_ID: LEX-P0024-V0036
% #FIELD_VALUE: 0.9\%氯化钠注射液用量
% #HANDWRITTEN: 19
\fieldvalue{\handwritten{19}}\ \text{ml}

\begin{tabularx}{\textwidth}{|X|X|X|}
\hline
名称 & 配制批号 & 有效期至\\
\hline
5$\times$裂解液 &
% #VALUE_ID: LEX-P0024-V0037
% #FIELD_VALUE: 配制批号
% #HANDWRITTEN: 5?0?7?0?1
\fieldvalue{\handwritten{5?0?7?0?1}}
&
% #VALUE_ID: LEX-P0024-V0038
% #FIELD_VALUE: 有效期至
% #HANDWRITTEN: 2026.12.27
\fieldvalue{\handwritten{2026.12.27}}\\
\hline
\end{tabularx}

\subsection{操作过程}

\begin{tabularx}{\textwidth}{|X|X|}
\hline
按照50 ml离心管的刻度读取样品体积，在取样袋中加入对应体积的5$\times$裂解液（每4 ml样品加入1 ml 5$\times$裂解液，按照“只进不舍”原则，如26 ml样品，加入7 ml 5$\times$裂解液），充分混匀后倒入至2.3的离心管中，置于37°C电热恒温水浴箱中裂解20 min后，取出离心管摇晃混匀，使样品充分裂解；再置于37°C电热恒温水槽中裂解10min得到细胞悬液。 &
样品体积：
% #VALUE_ID: LEX-P0024-V0039
% #FIELD_VALUE: 样品体积
% #HANDWRITTEN: 20
\fieldvalue{\handwritten{20}}\ ml

\\[-2pt]
&
裂解液：
% #VALUE_ID: LEX-P0024-V0040
% #FIELD_VALUE: 裂解液体积
% #HANDWRITTEN: 7
\fieldvalue{\handwritten{7}}\ ml\\
\cline{2-2}
&
电热恒温水箱：
% #VALUE_ID: LEX-P0024-V0041
% #FIELD_VALUE: 电热恒温水箱温度
% #HANDWRITTEN: 37
\fieldvalue{\handwritten{37}}\,$^{\circ}$C

\\[-2pt]
&
裂解：
% #VALUE_ID: LEX-P0024-V0042
% #FIELD_VALUE: 裂解时间
% #HANDWRITTEN: 04:42
\fieldvalue{\handwritten{04:42}}\\
\hline
取出裂解好的细胞悬液300$\times g$离心5分钟，按照读取样品体积加入对应体积的1\%人白蛋白稀释液重悬（每1.2 ml样品体积加入0.1 ml 1\%人白蛋白稀释液，按照“只进不舍”原则，如26 ml &
300$\times g$离心
% #VALUE_ID: LEX-P0024-V0043
% #FIELD_VALUE: 离心力
% #HANDWRITTEN: 300
\fieldvalue{\handwritten{300}}$\times g$，离心
% #VALUE_ID: LEX-P0024-V0044
% #FIELD_VALUE: 离心时间
% #HANDWRITTEN: 5
\fieldvalue{\handwritten{5}}\ 分钟\\
\cline{2-2}
&
加入1\%人白蛋白稀释液体积：
% #VALUE_ID: LEX-P0024-V0045
% #FIELD_VALUE: 1\%人白蛋白稀释液体积
% #HANDWRITTEN: 1.7
\fieldvalue{\handwritten{1.7}}\ ml\\
\cline{2-2}
&
吸取
% #VALUE_ID: LEX-P0024-V0046
% #FIELD_VALUE: 吸取体积
% #HANDWRITTEN: 200
\fieldvalue{\handwritten{200}}\ \text{$\mu$l}\\
\hline
\end{tabularx}

% TODO：本页底部操作步骤文字及表格在下一页继续。
% LEXOID_PAGE_COMPLETED: 24/60

\newpage

\noindent
\textbf{文件编号：} REC-YZ-SOP-PC-99-006-a \hfill
\textbf{版本号：} 1.4

\medskip

\noindent
\textbf{操作步骤：}

\begin{enumerate}
    \item 加入 $2.2\mathrm{ml}$ 人 B 淋巴细胞悬液，反复吹打至细胞悬液均匀。
    \item 摇去计数板上下两层保护膜，将细胞悬液混匀，从中吸取 $20\,\mu\mathrm{l}$ 与 $20\,\mu\mathrm{l}$ AO/PI 染料混合均匀。
    \item 吸取 $20\,\mu\mathrm{l}$ 细胞悬液与 AO/PI 染料的混合液，注入到计数板的一
    个检测室中。
    \item 将加入混合液的一端插入进样口。
    \item 选择 Assay 类型，使用 AO/PI 染料进行双荧光活率检测。
    \item 在主界面左侧“细胞输入”框本栏。
    \item 点击主界面左下角“Preview”预览细胞图片。
    \item 调整各荧光距直至观察到荧光通过下的细胞中心透亮，轮廓清晰。
    \item 点击“count”按钮进行计数并自动拍摄图片。
    \item 重复以上步骤 $2$ 次，共计数 $3$ 次，打印最终结果。
\end{enumerate}

\noindent
与
% #VALUE_ID: LEX-P0025-V0001
% #FIELD_VALUE: AO/PI 染料混合体积
% #TODO #HANDWRITTEN: 2\,\mu\mathrm{l}; 字迹较模糊
$\fieldvalue{\handwritten{2\,\mu\mathrm{l}}}$
AO/PI 染料混合均匀。

\noindent
吸取
% #VALUE_ID: LEX-P0025-V0002
% #FIELD_VALUE: 吸取体积
% #TODO #HANDWRITTEN: 20\,\mu\mathrm{l}; 蓝色手写覆盖印刷数字
$\fieldvalue{\handwritten{20\,\mu\mathrm{l}}}$
计数。

\noindent
每个样品均计数
% #VALUE_ID: LEX-P0025-V0003
% #FIELD_VALUE: 样品计数次数
% #HANDWRITTEN: 2
\fieldvalue{\handwritten{2}}
次。

\medskip

\noindent
\textbf{数据保存路径：}
% #VALUE_ID: LEX-P0025-V0004
% #FIELD_VALUE: 数据保存路径
% #TODO #HANDWRITTEN: /4/10加样单-班-17014-1?; 手写路径难以完全辨认
\fieldvalue{\handwritten{/4/10加样单-班-17014-1?}}

\begin{center}
\begin{tabularx}{\textwidth}{|c|X|X|c|c|c|c|}
\hline
\multicolumn{7}{|c|}{检测结果} \\
\hline
\multicolumn{7}{|l|}{检测结果：
% #VALUE_ID: LEX-P0025-V0005
% #FIELD_VALUE: 检测结果类型
% #HANDWRITTEN: \checkmark
\fieldvalue{\handwritten{\checkmark}}\ 一般反应细胞、二级反应细胞、三级反应细胞} \\
\hline
检测次数 &
活细胞密度（个细胞/ml） &
样品密度（个细胞/ml） &
活细胞密度 CV 值（\%） &
细胞活率（\%） &
样品活率（\%） &
细胞活率 CV 值（\%） \\
\hline
第一次 &
% #VALUE_ID: LEX-P0025-V0006
% #FIELD_VALUE: 第一次活细胞密度
% #TODO #HANDWRITTEN: 3.78\times10^{6}; 末位数字不清晰
\fieldvalue{\handwritten{$3.78\times10^{6}$}} &
&
&
% #VALUE_ID: LEX-P0025-V0007
% #FIELD_VALUE: 第一次细胞活率
% #HANDWRITTEN: 78.1
\fieldvalue{\handwritten{78.1}} &
& \\
\hline
第二次 &
% #VALUE_ID: LEX-P0025-V0008
% #FIELD_VALUE: 第二次活细胞密度
% #TODO #HANDWRITTEN: 3.44\times10^{6}; 手写数字略模糊
\fieldvalue{\handwritten{$3.44\times10^{6}$}} &
% #VALUE_ID: LEX-P0025-V0009
% #FIELD_VALUE: 样品密度
% #TODO #HANDWRITTEN: 3.20\times10^{5}; 指数上标不清晰
\fieldvalue{\handwritten{$3.20\times10^{5}$}} &
% #VALUE_ID: LEX-P0025-V0010
% #FIELD_VALUE: 活细胞密度 CV 值
% #HANDWRITTEN: 7
\fieldvalue{\handwritten{7}} &
% #VALUE_ID: LEX-P0025-V0011
% #FIELD_VALUE: 第二次细胞活率
% #HANDWRITTEN: 90.7
\fieldvalue{\handwritten{90.7}} &
% #VALUE_ID: LEX-P0025-V0012
% #FIELD_VALUE: 样品活率
% #HANDWRITTEN: 92
\fieldvalue{\handwritten{92}} &
% #VALUE_ID: LEX-P0025-V0013
% #FIELD_VALUE: 细胞活率 CV 值
% #HANDWRITTEN: 2
\fieldvalue{\handwritten{2}} \\
\hline
第三次 &
% #VALUE_ID: LEX-P0025-V0014
% #FIELD_VALUE: 第三次活细胞密度
% #TODO #HANDWRITTEN: 3.78\times10^{6}; 手写数字不清晰
\fieldvalue{\handwritten{$3.78\times10^{6}$}} &
&
&
% #VALUE_ID: LEX-P0025-V0015
% #FIELD_VALUE: 第三次细胞活率
% #HANDWRITTEN: 92.8
\fieldvalue{\handwritten{92.8}} &
& \\
\hline
\multicolumn{2}{|l|}{发酵罐细胞体积（ml）：
% #VALUE_ID: LEX-P0025-V0016
% #FIELD_VALUE: 发酵罐细胞体积
% #TODO #HANDWRITTEN: 180; 字迹较模糊
\fieldvalue{\handwritten{180}}}
&
\multicolumn{2}{l|}{细胞数量（个细胞/罐）：
% #VALUE_ID: LEX-P0025-V0017
% #FIELD_VALUE: 细胞数量
% #TODO #HANDWRITTEN: 1.??\times10^{?}; 数字难以辨认
$\fieldvalue{\handwritten{1.??\times10^{?}}}$}
& \multicolumn{3}{l|}{} \\
\hline
\end{tabularx}
\end{center}

\noindent
计算要求：样品密度为 $3$ 次活细胞密度的结果取平均值，除以 $1\%$ 人白血球计数液浓度补体，再除以生物反应器细胞液体积；细胞数量为样品密度乘以发酵罐体积。

\begin{tabularx}{\textwidth}{|p{0.16\textwidth}|X|}
\hline
备注 &
% #VALUE_ID: LEX-P0025-V0018
% #FIELD_VALUE: 备注
% #HANDWRITTEN: /
\fieldvalue{\handwritten{/}} \\
\hline
实验人/日期 &
% #VALUE_ID: LEX-P0025-V0019
% #FIELD_VALUE: 实验人
% #TODO #HANDWRITTEN: 韩??; 姓名字迹不清晰
\fieldvalue{\handwritten{韩??}}\quad
% #VALUE_ID: LEX-P0025-V0020
% #FIELD_VALUE: 实验日期
% #TODO #HANDWRITTEN: 2024.4.?; 日期末位不清晰
\fieldvalue{\handwritten{2024.4.?}} \\
\hline
复核人/日期 &
% #VALUE_ID: LEX-P0025-V0021
% #FIELD_VALUE: 复核人
% #TODO #HANDWRITTEN: 李春霞; 字迹略模糊
\fieldvalue{\handwritten{李春霞}}\quad
% #VALUE_ID: LEX-P0025-V0022
% #FIELD_VALUE: 复核日期
% #HANDWRITTEN: 2025.8.7
\fieldvalue{\handwritten{2025.8.7}} \\
\hline
\end{tabularx}

\begin{center}
第 3 页 共 5 页
\end{center}
% LEXOID_PAGE_COMPLETED: 25/60

\newpage

\begin{center}
\textbf{生物反应器细胞液检验操作记录}
\end{center}

\noindent
文件编码：REC-YZ-SOP-QC-99-006-a

\begin{center}
\begin{tabularx}{\textwidth}{|p{2.5cm}|X|p{2.5cm}|X|}
\hline
\multicolumn{4}{|c|}{细胞生长状态检验}\\
\hline
样品名称 &
% #VALUE_ID: LEX-P0026-V0001
% #FIELD_VALUE: 样品名称
\fieldvalue{一级生物反应器细胞液} &
样品规格 &
% #VALUE_ID: LEX-P0026-V0002
% #FIELD_VALUE: 样品规格
% #HANDWRITTEN: 12? ml/袋
\fieldvalue{\handwritten{12? ml/袋}}\\
\hline
样品批号 &
% #VALUE_ID: LEX-P0026-V0003
% #FIELD_VALUE: 样品批号
% #HANDWRITTEN: [批号，字迹不清]
\fieldvalue{\handwritten{[批号，字迹不清]}} &
取样日期 &
% #VALUE_ID: LEX-P0026-V0004
% #FIELD_VALUE: 取样日期
% #HANDWRITTEN: 2024年03月17日
\fieldvalue{\handwritten{2024年03月17日}}\\
\hline
\end{tabularx}
\end{center}

\section{仪器设备及试剂}

\subsection{仪器设备}

\begin{center}
\begin{tabularx}{\textwidth}{|X|X|X|X|X|}
\hline
仪器名称 & 厂家 & 型号 & 设备编号 & 检查确认\\
\hline
% #VALUE_ID: LEX-P0026-V0005
% #FIELD_VALUE: 仪器名称
\fieldvalue{荧光显微镜} &
% #VALUE_ID: LEX-P0026-V0006
% #FIELD_VALUE: 厂家
\fieldvalue{徕卡} &
% #VALUE_ID: LEX-P0026-V0007
% #FIELD_VALUE: 型号
\fieldvalue{Dmil LED} &
% #VALUE_ID: LEX-P0026-V0008
% #FIELD_VALUE: 设备编号
\fieldvalue{1431404} &
% #VALUE_ID: LEX-P0026-V0009
% #FIELD_VALUE: 检查确认
\fieldvalue{\ensuremath{\boxtimes}正常\quad $\square$异常}\\
\hline
\end{tabularx}
\end{center}

\subsection{试剂耗材}

\begin{center}
\begin{tabularx}{\textwidth}{|X|X|X|X|X|}
\hline
名称 & 厂家 & 货号 & 批号 & 有效期至\\
\hline
% #VALUE_ID: LEX-P0026-V0010
% #FIELD_VALUE: 名称
\fieldvalue{细胞活力检测试剂盒} &
% #VALUE_ID: LEX-P0026-V0011
% #FIELD_VALUE: 厂家
\fieldvalue{江苏凯基生物技术股份有限公司} &
% #VALUE_ID: LEX-P0026-V0012
% #FIELD_VALUE: 货号
\fieldvalue{KGA9501-1000} &
% #VALUE_ID: LEX-P0026-V0013
% #FIELD_VALUE: 批号
% #HANDWRITTEN: 202?0400
\fieldvalue{\handwritten{202?0400}} &
% #VALUE_ID: LEX-P0026-V0014
% #FIELD_VALUE: 有效期至
% #HANDWRITTEN: 2024.10.?
\fieldvalue{\handwritten{2024.10.?}}\\
\hline
\end{tabularx}
\end{center}

\section{操作步骤}

\begin{center}
\begin{tabularx}{\textwidth}{|p{0.56\textwidth}|X|}
\hline
\multicolumn{2}{|c|}{溶液配制}\\
\hline
20$\times$的荧光染液配制：取
% #VALUE_ID: LEX-P0026-V0015
% #FIELD_VALUE: 组分A用量
% #HANDWRITTEN: 0.2 $\mu$l
\fieldvalue{\handwritten{0.2 $\mu$l}}组分A和
% #VALUE_ID: LEX-P0026-V0016
% #FIELD_VALUE: 组分B用量
% #HANDWRITTEN: 0.2 $\mu$l
\fieldvalue{\handwritten{0.2 $\mu$l}}组分B与
% #VALUE_ID: LEX-P0026-V0017
% #FIELD_VALUE: PBS用量
% #HANDWRITTEN: 7.8 $\mu$l
\fieldvalue{\handwritten{7.8 $\mu$l}} PBS混合均匀配成20$\times$的荧光染液。&
\\
\hline
\multicolumn{2}{|c|}{操作过程}\\
\hline
取混合后光样品200$\mu$l至96孔板中，取2孔；&
样品
% #VALUE_ID: LEX-P0026-V0018
% #FIELD_VALUE: 样品加样量
% #HANDWRITTEN: 200 $\mu$l/孔
\fieldvalue{\handwritten{200 $\mu$l/孔}}\\
&
每个样品取
% #VALUE_ID: LEX-P0026-V0019
% #FIELD_VALUE: 每个样品孔数
% #HANDWRITTEN: 2孔
\fieldvalue{\handwritten{2孔}}\\
\hline
自然沉降后，吸弃上清，注意不要吸到沉淀样品；&
% #VALUE_ID: LEX-P0026-V0020
% #FIELD_VALUE: 吸弃上清
\fieldvalue{\ensuremath{\boxtimes}吸弃上清}\\
\hline
取20$\times$的荧光染液加入到96孔板中，200$\mu$l/孔，避光孵育30--60min；&
20$\times$的荧光染液
% #VALUE_ID: LEX-P0026-V0021
% #FIELD_VALUE: 荧光染液加样量
% #HANDWRITTEN: 200 $\mu$l/孔
\fieldvalue{\handwritten{200 $\mu$l/孔}}\\
&
避光孵育
% #VALUE_ID: LEX-P0026-V0022
% #FIELD_VALUE: 避光孵育时间
% #TODO #HANDWRITTEN: [孵育时间，字迹不清]; 蓝色手写内容难以辨认
\fieldvalue{\handwritten{[孵育时间，字迹不清]}}\\
\hline
染色结束后，吸弃上清，取PBS加入到96孔板中，200$\mu$l/孔，吸弃上清，注意不要吸到沉淀样品，清洗3次，之后每孔加入100$\mu$l PBS进行拍照；&
加入PBS
% #VALUE_ID: LEX-P0026-V0023
% #FIELD_VALUE: PBS加样量
% #HANDWRITTEN: 200 $\mu$l/孔
\fieldvalue{\handwritten{200 $\mu$l/孔}}\\
&
清洗
% #VALUE_ID: LEX-P0026-V0024
% #FIELD_VALUE: 清洗次数
% #HANDWRITTEN: 2次
\fieldvalue{\handwritten{2次}}\\
\hline
将荧光显微镜调至4$\times$物镜，软件中调节荧光物镜参数，调节到适合的光源和像距，拍摄4张清晰图像。进一步清晰图像，打印并粘贴在实验记录上。&
\\
\hline
数据保存路径&
% #VALUE_ID: LEX-P0026-V0025
% #FIELD_VALUE: 数据保存路径
% #TODO #HANDWRITTEN: [数据保存路径，字迹不清]; 手写路径部分无法可靠辨认
\fieldvalue{\handwritten{[数据保存路径，字迹不清]}}\\
\hline
\multicolumn{2}{|c|}{检测结果}\\
\hline
\end{tabularx}
\end{center}

% TODO：本页表格及操作记录内容在下一页可能继续。
% LEXOID_PAGE_COMPLETED: 26/60

\newpage

\noindent
\begin{minipage}{\textwidth}
\small
\begin{tabular}{@{}p{0.28\textwidth}p{0.38\textwidth}p{0.28\textwidth}@{}}
\textbf{铂业生产检验科技（北京）有限公司} &
文件编码：REC-YZ-SOP-QC-99-006-a &
\\[-2pt]
\textit{PLATINUM LIME EXCELLENCE PITECH} &
&
\\
\end{tabular}

\vspace{-0.2cm}
\hfill
\begin{tabular}{|c|}
\hline
文件生效章\\
% #VALUE_ID: LEX-P0027-V0001
% #FIELD_VALUE: 文件状态
\fieldvalue{原件}\\
\hline
\end{tabular}
\hspace{0.3cm}
\begin{tabular}{|c|}
\hline
版本号\\
% #VALUE_ID: LEX-P0027-V0002
% #FIELD_VALUE: 版本号
\fieldvalue{1.4}\\
\hline
\end{tabular}

\vspace{0.1cm}
\noindent\rule{\textwidth}{0.5pt}

\vspace{15.5cm}
\begin{center}
% TODO: 插入本页中央可见的显微图像；原 PDF 未提供独立图像文件。
\fbox{\rule{0pt}{3.2cm}\rule{0.36\textwidth}{0pt}}\\[-2pt]
% #VALUE_ID: LEX-P0027-V0003
% #TODO #HANDWRITTEN: [难以辨认]; 图像下方蓝色手写标注不清晰
\handwritten{[难以辨认]}
\end{center}

\vfill

\begin{tabularx}{\textwidth}{|p{0.16\textwidth}|X|}
\hline
备注 &
% #VALUE_ID: LEX-P0027-V0004
% #FIELD_VALUE: 备注
% #HANDWRITTEN: /
\fieldvalue{\handwritten{/}}\\
\hline
实验人/日期 &
% #VALUE_ID: LEX-P0027-V0005
% #FIELD_VALUE: 实验人
% #TODO #HANDWRITTEN: [难以辨认]; 手写姓名字迹不清
\fieldvalue{\handwritten{[难以辨认]}}\quad
% #VALUE_ID: LEX-P0027-V0006
% #FIELD_VALUE: 实验日期
% #HANDWRITTEN: 2025.8.7
\fieldvalue{\handwritten{2025.8.7}}\\
\hline
复核人/日期 &
% #VALUE_ID: LEX-P0027-V0007
% #FIELD_VALUE: 复核人
% #TODO #HANDWRITTEN: [难以辨认]; 手写姓名字迹不清
\fieldvalue{\handwritten{[难以辨认]}}\quad
% #VALUE_ID: LEX-P0027-V0008
% #FIELD_VALUE: 复核日期
% #HANDWRITTEN: 2025.8.7
\fieldvalue{\handwritten{2025.8.7}}\\
\hline
\end{tabularx}

\begin{center}
\scriptsize 第 5 页 共 5 页
\end{center}
\end{minipage}
% LEXOID_PAGE_COMPLETED: 27/60

\newpage

\noindent Assay:
% #VALUE_ID: LEX-P0028-V0001
% #FIELD_VALUE: Assay
\fieldvalue{MSC-1}

\medskip
\noindent Cell Type F1:
% #VALUE_ID: LEX-P0028-V0002
% #FIELD_VALUE: Cell Type F1
\fieldvalue{MSC (AO)}

\noindent Cell Type F2:
% #VALUE_ID: LEX-P0028-V0003
% #FIELD_VALUE: Cell Type F2
\fieldvalue{MSC (PI)}

\medskip
\noindent Sample ID:
% #VALUE_ID: LEX-P0028-V0004
% #FIELD_VALUE: Sample ID
\fieldvalue{QC-JS-QY0401-A31Z2202508008-20250829-70H-1}

\noindent Dilution:
% #VALUE_ID: LEX-P0028-V0005
% #FIELD_VALUE: Dilution
\fieldvalue{2.00}

\medskip
\noindent\rule{\textwidth}{0.4pt}

\noindent Results:

\medskip
\noindent
\begin{tabular}{@{}lll@{}}
\textbf{Count} & \textbf{Concentration} & \textbf{Mean Diameter} \\
\multicolumn{1}{@{}l}{\texttt{=============}} &
\multicolumn{1}{l}{\texttt{=============}} &
\multicolumn{1}{l@{}}{\texttt{=============}} \\[2pt]
Total:
% #VALUE_ID: LEX-P0028-V0006
% #FIELD_VALUE: Count -- Total
\fieldvalue{1201}
&
Total:
% #VALUE_ID: LEX-P0028-V0010
% #FIELD_VALUE: Concentration -- Total
\fieldvalue{$4.22\times10^{6}$ cells/mL}
&
Total:
% #VALUE_ID: LEX-P0028-V0013
% #FIELD_VALUE: Mean Diameter -- Total
\fieldvalue{13.2 micron}
\\
Live:
% #VALUE_ID: LEX-P0028-V0007
% #FIELD_VALUE: Count -- Live
\fieldvalue{1118}
&
Live:
% #VALUE_ID: LEX-P0028-V0011
% #FIELD_VALUE: Concentration -- Live
\fieldvalue{$3.93\times10^{6}$ cells/mL}
&
Live:
% #VALUE_ID: LEX-P0028-V0014
% #FIELD_VALUE: Mean Diameter -- Live
\fieldvalue{13.2 microns}
\\
Dead:
% #VALUE_ID: LEX-P0028-V0008
% #FIELD_VALUE: Count -- Dead
\fieldvalue{83}
&
Dead:
% #VALUE_ID: LEX-P0028-V0012
% #FIELD_VALUE: Concentration -- Dead
\fieldvalue{$2.92\times10^{5}$ cells/mL}
&
Dead:
% #VALUE_ID: LEX-P0028-V0015
% #FIELD_VALUE: Mean Diameter -- Dead
\fieldvalue{13.5 microns}
\\
Viability:
% #VALUE_ID: LEX-P0028-V0009
% #FIELD_VALUE: Viability
\fieldvalue{93.1\%}
&
&
\end{tabular}

\medskip
\begin{center}
\begin{tabular}{cc}
% TODO: Image visible on page; source filename unavailable.
\fbox{\parbox[c][4.2cm][c]{0.38\textwidth}{\centering Image placeholder}}
&
% TODO: Image visible on page; source filename unavailable.
\fbox{\parbox[c][4.2cm][c]{0.38\textwidth}{\centering Image placeholder}}
\\[0.5cm]
% TODO: Image visible on page; source filename unavailable.
\fbox{\parbox[c][4.2cm][c]{0.38\textwidth}{\centering Image placeholder}}
&
% TODO: Image visible on page; source filename unavailable.
\fbox{\parbox[c][4.2cm][c]{0.38\textwidth}{\centering Image placeholder}}
\end{tabular}
\end{center}

\medskip
\noindent 操作人/审核人:
% #VALUE_ID: LEX-P0028-V0016
% #FIELD_VALUE: 操作人/审核人
% #TODO #HANDWRITTEN: illegible signature; handwritten signature is not clearly readable
\fieldvalue{\handwritten{[illegible signature]}}

\hfill 日期:
% #VALUE_ID: LEX-P0028-V0017
% #FIELD_VALUE: 日期
% #HANDWRITTEN: 2025.08.29
\fieldvalue{\handwritten{2025.08.29}}
% LEXOID_PAGE_COMPLETED: 28/60

\newpage

\noindent Assay:
% #VALUE_ID: LEX-P0029-V0001
% #FIELD_VALUE: Assay
\fieldvalue{MSC-1}

\noindent Cell Type F1:
% #VALUE_ID: LEX-P0029-V0002
% #FIELD_VALUE: Cell Type F1
\fieldvalue{MSC (AO)}

\noindent Cell Type F2:
% #VALUE_ID: LEX-P0029-V0003
% #FIELD_VALUE: Cell Type F2
\fieldvalue{MSC (PI)}

\medskip

\noindent Sample ID:
% #VALUE_ID: LEX-P0029-V0004
% #FIELD_VALUE: Sample ID
\fieldvalue{QC-JS-QY0401-A31Z220250808-20250829-70H-2}

\noindent Dilution:
% #VALUE_ID: LEX-P0029-V0005
% #FIELD_VALUE: Dilution
\fieldvalue{2.00}

\medskip
\noindent\rule{\textwidth}{0.4pt}

\noindent Results:

\medskip

\renewcommand{\arraystretch}{1.15}
\noindent
\begin{tabularx}{\textwidth}{@{}>{\RaggedRight\arraybackslash}X
                              >{\RaggedRight\arraybackslash}X
                              >{\RaggedRight\arraybackslash}X@{}}
\textbf{Count} & \textbf{Concentration} & \textbf{Mean Diameter} \\
\underline{\texttt{==============}} &
\underline{\texttt{==============}} &
\underline{\texttt{==============}} \\

Total:
% #VALUE_ID: LEX-P0029-V0006
% #FIELD_VALUE: Total count
\fieldvalue{1092}
&
% #VALUE_ID: LEX-P0029-V0007
% #FIELD_VALUE: Total concentration
\fieldvalue{$3.85\times10^{6}$ cells/mL}
&
% #VALUE_ID: LEX-P0029-V0008
% #FIELD_VALUE: Total mean diameter
\fieldvalue{14.0 micron}
\\

Live:
% #VALUE_ID: LEX-P0029-V0009
% #FIELD_VALUE: Live count
\fieldvalue{990}
&
% #VALUE_ID: LEX-P0029-V0010
% #FIELD_VALUE: Live concentration
\fieldvalue{$3.49\times10^{6}$ cells/mL}
&
% #VALUE_ID: LEX-P0029-V0011
% #FIELD_VALUE: Live mean diameter
\fieldvalue{14.1 microns}
\\

Dead:
% #VALUE_ID: LEX-P0029-V0012
% #FIELD_VALUE: Dead count
\fieldvalue{102}
&
% #VALUE_ID: LEX-P0029-V0013
% #FIELD_VALUE: Dead concentration
\fieldvalue{$3.58\times10^{5}$ cells/mL}
&
% #VALUE_ID: LEX-P0029-V0014
% #FIELD_VALUE: Dead mean diameter
\fieldvalue{12.7 micron}
\\

Viability:
% #VALUE_ID: LEX-P0029-V0015
% #FIELD_VALUE: Viability
\fieldvalue{90.7\%}
&
&
\end{tabularx}

\bigskip

\noindent
\begin{minipage}[t]{0.47\textwidth}
\centering
% TODO: Insert the upper-left microscopy image; no filename is available.
\fbox{\rule{0pt}{0.22\textheight}\rule{0.95\linewidth}{0pt}}
\end{minipage}
\hfill
\begin{minipage}[t]{0.47\textwidth}
\centering
% TODO: Insert the upper-right microscopy image; no filename is available.
\fbox{\rule{0pt}{0.22\textheight}\rule{0.95\linewidth}{0pt}}
\end{minipage}

\medskip

\noindent
\begin{minipage}[t]{0.47\textwidth}
\centering
% TODO: Insert the lower-left microscopy image; no filename is available.
\fbox{\rule{0pt}{0.22\textheight}\rule{0.95\linewidth}{0pt}}
\end{minipage}
\hfill
\begin{minipage}[t]{0.47\textwidth}
\centering
% TODO: Insert the lower-right microscopy image; no filename is available.
\fbox{\rule{0pt}{0.22\textheight}\rule{0.95\linewidth}{0pt}}
\end{minipage}

\bigskip

\noindent 操作人/审核人：
% #VALUE_ID: LEX-P0029-V0016
% #FIELD_VALUE: 操作人/审核人
% #TODO #HANDWRITTEN: illegible signature; handwritten Chinese signature cannot be read confidently
\fieldvalue{\handwritten{[illegible signature]}}

\hfill
日期：
% #VALUE_ID: LEX-P0029-V0017
% #FIELD_VALUE: 日期
% #HANDWRITTEN: 2025.8.29
\fieldvalue{\handwritten{2025.8.29}}
% LEXOID_PAGE_COMPLETED: 29/60

\newpage

\noindent Assay: \\
% #VALUE_ID: LEX-P0030-V0001
% #FIELD_VALUE: Assay
\fieldvalue{MSC-1}

\medskip
\noindent Cell Type F1: \\
% #VALUE_ID: LEX-P0030-V0002
% #FIELD_VALUE: Cell Type F1
\fieldvalue{MSC (AO)}

\noindent Cell Type F2: \\
% #VALUE_ID: LEX-P0030-V0003
% #FIELD_VALUE: Cell Type F2
\fieldvalue{MSC (PI)}

\medskip
\noindent Sample ID: \\
% #VALUE_ID: LEX-P0030-V0004
% #FIELD_VALUE: Sample ID
\fieldvalue{QC-JS-QY0401-A31Z2202508008-20250829-70H-3}

\noindent Dilution: \\
% #VALUE_ID: LEX-P0030-V0005
% #FIELD_VALUE: Dilution
\fieldvalue{2.00}

\medskip
\noindent\rule{\textwidth}{0.4pt}

\noindent Results:

\medskip
\begin{tabularx}{\textwidth}{@{}XXX@{}}
\textbf{Count} & \textbf{Concentration} & \textbf{Mean Diameter} \\
\underline{\hspace{0.9\linewidth}} &
\underline{\hspace{0.9\linewidth}} &
\underline{\hspace{0.9\linewidth}} \\[2pt]
Total:
% #VALUE_ID: LEX-P0030-V0006
% #FIELD_VALUE: Total count
\fieldvalue{1185}
&
% #VALUE_ID: LEX-P0030-V0007
% #FIELD_VALUE: Total concentration
\fieldvalue{$4.18 \times 10^{6}$ cells/mL}
&
% #VALUE_ID: LEX-P0030-V0008
% #FIELD_VALUE: Total mean diameter
\fieldvalue{14.0 micron}
\\
Live:
% #VALUE_ID: LEX-P0030-V0009
% #FIELD_VALUE: Live count
\fieldvalue{1099}
&
% #VALUE_ID: LEX-P0030-V0010
% #FIELD_VALUE: Live concentration
\fieldvalue{$3.88 \times 10^{6}$ cells/mL}
&
% #VALUE_ID: LEX-P0030-V0011
% #FIELD_VALUE: Live mean diameter
\fieldvalue{14.2 microns}
\\
Dead:
% #VALUE_ID: LEX-P0030-V0012
% #FIELD_VALUE: Dead count
\fieldvalue{86}
&
% #VALUE_ID: LEX-P0030-V0013
% #FIELD_VALUE: Dead concentration
\fieldvalue{$3.02 \times 10^{5}$ cells/mL}
&
% #VALUE_ID: LEX-P0030-V0014
% #FIELD_VALUE: Dead mean diameter
\fieldvalue{12.4 microns}
\\
Viability:
% #VALUE_ID: LEX-P0030-V0015
% #FIELD_VALUE: Viability
\fieldvalue{92.8\%}
&
&
\end{tabularx}

\medskip

% TODO: Four microscopy images are visible on this page; no source filenames were provided.
\begin{center}
\begin{minipage}{0.43\textwidth}
\centering
\fbox{\parbox[c][4.2cm][c]{0.95\linewidth}{\centering TODO: microscopy image}}
\end{minipage}
\hfill
\begin{minipage}{0.43\textwidth}
\centering
\fbox{\parbox[c][4.2cm][c]{0.95\linewidth}{\centering TODO: microscopy image}}
\end{minipage}

\vspace{0.4cm}

\begin{minipage}{0.43\textwidth}
\centering
\fbox{\parbox[c][4.2cm][c]{0.95\linewidth}{\centering TODO: microscopy image}}
\end{minipage}
\hfill
\begin{minipage}{0.43\textwidth}
\centering
\fbox{\parbox[c][4.2cm][c]{0.95\linewidth}{\centering TODO: microscopy image}}
\end{minipage}
\end{center}

\vfill

\noindent 操作人/审核人：%
% #VALUE_ID: LEX-P0030-V0016
% #FIELD_VALUE: 操作人/审核人
% #TODO #HANDWRITTEN: illegible signature; cursive blue-ink signature
\fieldvalue{\handwritten{[illegible signature]}}

\hfill 日期：%
% #VALUE_ID: LEX-P0030-V0017
% #FIELD_VALUE: 日期
% #HANDWRITTEN: 2025.3.29
\fieldvalue{\handwritten{2025.3.29}}
% LEXOID_PAGE_COMPLETED: 30/60

\newpage

\begin{center}
\textbf{铂生物科技（北京）有限公司}\\
\textit{PLATINUM LIFE EXCELLENCE BIOTECH}\\[2mm]
原料请验单
\end{center}

\noindent
文件编码：REC-YZ-SMP-QC-01-001-d
\hfill
版本号：1.4

\begin{center}
\small
\begin{tabular}{|p{2.5cm}|p{5.0cm}|p{2.5cm}|p{5.0cm}|}
\hline
样品名称 &
% #VALUE_ID: LEX-P0031-V0001
% #FIELD_VALUE: 样品名称
% #TODO #HANDWRITTEN: [难辨]; 蓝色手写内容辨识度不足
\fieldvalue{\handwritten{[难辨]}} &
样品代码 &
% #VALUE_ID: LEX-P0031-V0002
% #FIELD_VALUE: 样品代码
% #TODO #HANDWRITTEN: 2-Y?4?; 蓝色手写代码部分难辨
\fieldvalue{\handwritten{2-Y?4?}} \\
\hline
样品批号 &
% #VALUE_ID: LEX-P0031-V0003
% #FIELD_VALUE: 样品批号
% #TODO #HANDWRITTEN: A?12 24? 05; 蓝色手写批号部分难辨
\fieldvalue{\handwritten{A?12 24? 05}} &
样品规格/数量 &
% #VALUE_ID: LEX-P0031-V0004
% #FIELD_VALUE: 样品规格/数量
% #TODO #HANDWRITTEN: 2?/2?/2?; 蓝色手写内容部分难辨
\fieldvalue{\handwritten{2?/2?/2?}} \\
\hline
\multirow{2}{*}{样品类别} &
% #VALUE_ID: LEX-P0031-V0005
% #FIELD_VALUE: 样品类别：生产过程检验品
\fieldvalue{\checkmark}\ 生产过程检验品
& \multicolumn{2}{l|}{\(\square\) 待包装产品\qquad \(\square\) 成品} \\
& \(\square\) 其他（\hspace{2cm}） & \multicolumn{2}{l|}{} \\
\hline
检验类型 &
% #VALUE_ID: LEX-P0031-V0006
% #FIELD_VALUE: 检验类型：初验
\fieldvalue{\checkmark}\ 初验
\qquad \(\square\) 复验
\qquad \(\square\) 退货
\qquad \(\square\) 稳定性
\qquad \(\square\) 留样
& \multicolumn{2}{l|}{} \\
\hline
请验目的 &
% #VALUE_ID: LEX-P0031-V0007
% #FIELD_VALUE: 请验目的：全检
\fieldvalue{\checkmark}\ 全检
\qquad \(\square\) 留样
\qquad \(\square\) 其他（\hspace{2cm}）
& \multicolumn{2}{l|}{} \\
\hline
请验部门 &
% #VALUE_ID: LEX-P0031-V0008
% #FIELD_VALUE: 请验部门：生产部
\fieldvalue{\checkmark}\ 生产部
\qquad \(\square\) 物料管理部
\qquad \(\square\) 其他（\hspace{2cm}）
& \multicolumn{2}{l|}{} \\
\hline
样品情况 &
% #VALUE_ID: LEX-P0031-V0009
% #FIELD_VALUE: 样品情况：随请验单一起提供
\fieldvalue{\checkmark}\ 随请验单一起提供
\qquad \(\square\) 未随请验单一起提供
& \multicolumn{2}{l|}{} \\
\hline
请验人 &
% #VALUE_ID: LEX-P0031-V0010
% #FIELD_VALUE: 请验人
% #TODO #HANDWRITTEN: 林??; 蓝色签字部分难辨
\fieldvalue{\handwritten{林??}} &
请验日期 &
% #VALUE_ID: LEX-P0031-V0011
% #FIELD_VALUE: 请验日期
% #HANDWRITTEN: 2024.08.28
\fieldvalue{\handwritten{2024.08.28}} \\
\hline
收验人 &
% #VALUE_ID: LEX-P0031-V0012
% #FIELD_VALUE: 收验人
% #TODO #HANDWRITTEN: 郭??; 蓝色签字部分难辨
\fieldvalue{\handwritten{郭??}} &
收验日期 &
% #VALUE_ID: LEX-P0031-V0013
% #FIELD_VALUE: 收验日期
% #HANDWRITTEN: 2024.8.28
\fieldvalue{\handwritten{2024.8.28}} \\
\hline
备注 &
% #VALUE_ID: LEX-P0031-V0014
% #FIELD_VALUE: 备注
% #TODO #HANDWRITTEN: 5h - 2.0h - 2.0L; 蓝色手写备注部分难辨
\fieldvalue{\handwritten{5h - 2.0h - 2.0L}} &
\multicolumn{2}{l|}{} \\
\hline
\end{tabular}
\end{center}
% LEXOID_PAGE_COMPLETED: 31/60

\newpage

% TODO: Page header includes a company logo; no image filename is available.
\noindent\textbf{铂金生命科技（北京）有限公司}\\[-2pt]
\textit{PLATINUM LIFE EXCELLENCE BIOTECH}\hfill
文件编码：REC-YZ-SMP-QC-01-004-b \quad 版本号：1.4

\begin{center}
    \textbf{\large 取样记录}
\end{center}

\renewcommand{\arraystretch}{1.65}
\small
\begin{tabularx}{\textwidth}{|p{0.14\textwidth}|p{0.18\textwidth}|p{0.12\textwidth}|p{0.14\textwidth}|p{0.12\textwidth}|X|}
\hline
名称 &
\multicolumn{2}{c|}{
% #VALUE_ID: LEX-P0032-V0001
% #FIELD_VALUE: 名称
\fieldvalue{一缓生物反应器细胞液}
} &
代码 &
\multicolumn{2}{c|}{
% #VALUE_ID: LEX-P0032-V0002
% #FIELD_VALUE: 代码
\fieldvalue{QY0401}
}
\\ \hline

规格 &
\multicolumn{2}{c|}{
% #VALUE_ID: LEX-P0032-V0003
% #FIELD_VALUE: 规格
\fieldvalue{约20ml/袋}
} &
总数量 &
\multicolumn{2}{c|}{
% #VALUE_ID: LEX-P0032-V0004
% #FIELD_VALUE: 总数量
\fieldvalue{1袋}
}
\\ \hline

批号 &
\multicolumn{2}{c|}{
% #VALUE_ID: LEX-P0032-V0005
% #FIELD_VALUE: 批号
\fieldvalue{A3122202508008}
} &
进厂批号 &
\multicolumn{2}{c|}{
% #VALUE_ID: LEX-P0032-V0006
% #FIELD_VALUE: 进厂批号
\fieldvalue{/}
}
\\ \hline

供货单位/厂家 &
\multicolumn{5}{c|}{
% #VALUE_ID: LEX-P0032-V0007
% #FIELD_VALUE: 供货单位/厂家
\fieldvalue{/}
}
\\ \hline

取样目的 &
\multicolumn{5}{c|}{
% #VALUE_ID: LEX-P0032-V0008
% #FIELD_VALUE: 取样目的
\fieldvalue{$\boxtimes$ 正常检验}\quad
$\square$ 复检\quad
$\square$ 稳定性
\newline
\hspace*{1.8em}$\square$ 退货检验\quad
$\square$ 留样\quad
$\square$ 其他：\underline{\hspace{3em}}}
\\ \hline

取样地点 &
\multicolumn{5}{c|}{
% #VALUE_ID: LEX-P0032-V0009
% #FIELD_VALUE: 取样地点
\fieldvalue{$\boxtimes$ 岛出口}\quad
$\square$ 物料库\quad
$\square$ 细胞库\quad
$\square$ 包装间
}
\\ \hline

贮存条件 &
\multicolumn{5}{c|}{
% #VALUE_ID: LEX-P0032-V0010
% #FIELD_VALUE: 贮存条件
\fieldvalue{$\boxtimes$2--8$^\circ$C}\quad
$\square$-20$^\circ$C\quad
$\square$-80$^\circ$C\quad
$\square$ 常温\quad
$\square$ 其他：\underline{\hspace{3em}}}
\\ \hline

盛装容器 &
\multicolumn{5}{c|}{
% #VALUE_ID: LEX-P0032-V0011
% #FIELD_VALUE: 盛装容器
\fieldvalue{$\boxtimes$ 保温箱}\quad
$\square$ 泡沫箱\quad
$\square$ 手提筒\quad
$\square$ 其他：\underline{\hspace{3em}}}
\\ \hline

外观复核 &
\multicolumn{5}{c|}{
% #VALUE_ID: LEX-P0032-V0012
% #FIELD_VALUE: 外观复核
\fieldvalue{$\boxtimes$ 合格}\quad
$\square$ 不合格
}
\\ \hline

取样量 &
% #VALUE_ID: LEX-P0032-V0013
% #FIELD_VALUE: 取样量
\fieldvalue{1袋}
&
检验量 &
% #VALUE_ID: LEX-P0032-V0014
% #FIELD_VALUE: 检验量
\fieldvalue{1袋}
&
留样量 &
% #VALUE_ID: LEX-P0032-V0015
% #FIELD_VALUE: 留样量
\fieldvalue{/}
\\ \hline

取样人/日期 &
\multicolumn{2}{c|}{
% #VALUE_ID: LEX-P0032-V0016
% #FIELD_VALUE: 取样人/日期
% #HANDWRITTEN: 钱峰 2025.8.28
\fieldvalue{\handwritten{钱峰 2025.8.28}}
} &
复核人/日期 &
\multicolumn{2}{c|}{
% #VALUE_ID: LEX-P0032-V0017
% #FIELD_VALUE: 复核人/日期
% #HANDWRITTEN: 李春梅 2025.8.28
\fieldvalue{\handwritten{李春梅 2025.8.28}}
}
\\ \hline

备注 &
\multicolumn{5}{c|}{
% #VALUE_ID: LEX-P0032-V0018
% #FIELD_VALUE: 备注
% TODO #HANDWRITTEN: 升 2.0L; handwriting is unclear.
\fieldvalue{\handwritten{升 2.0L}}
}
\\ \hline
\end{tabularx}

\normalsize
\begin{center}
第 1 页 共 1 页
\end{center}
% LEXOID_PAGE_COMPLETED: 32/60

\newpage

\begin{center}
\textbf{生物反应器细胞液检测操作记录}
\end{center}

\begin{tabularx}{\textwidth}{|l|X|l|X|}
\hline
\multicolumn{4}{|c|}{葡萄糖浓度检验} \\ \hline
样品名称 &
% #VALUE_ID: LEX-P0033-V0001
% #FIELD_VALUE: 样品名称
\fieldvalue{生物反应器细胞液} &
样品规格 &
% #VALUE_ID: LEX-P0033-V0002
% #FIELD_VALUE: 样品规格
% #HANDWRITTEN: 1.20
\fieldvalue{\handwritten{1.20}\(\,\mathrm{ml/}\text{袋}\)} \\ \hline
样品批号 &
% #VALUE_ID: LEX-P0033-V0003
% #FIELD_VALUE: 样品批号
% #TODO #HANDWRITTEN: 疑似“2024……”; 字迹不清
\fieldvalue{\handwritten{2024……}} &
取样日期 &
% #VALUE_ID: LEX-P0033-V0004
% #FIELD_VALUE: 取样日期
% #HANDWRITTEN: 2024年3月28日
\fieldvalue{\handwritten{2024年3月28日}} \\ \hline
\end{tabularx}

\vspace{0.4em}

\textbf{1\quad 仪器设备、试剂耗材}

\textbf{1.1\quad 仪器设备}

\begin{tabularx}{\textwidth}{|l|X|l|X|X|}
\hline
名称 & 厂家 & 型号 & 设备编号 & 检查确认 \\ \hline
血糖仪 &
三诺生物传感股份有限公司 &
安稳+ &
 &
% #VALUE_ID: LEX-P0033-V0005
% #FIELD_VALUE: 检查确认（血糖仪）
\fieldvalue{\ensuremath{\boxtimes}\ 正常\quad \(\square\)\ 异常}
\\ \hline
移液器 &
Eppendorf Corporate &
0.1$\sim$2.5\,\textmu l &
% #VALUE_ID: LEX-P0033-V0006
% #FIELD_VALUE: 设备编号
% #TODO #HANDWRITTEN: 027?26; 字迹不清
\fieldvalue{\handwritten{027?26}} &
% #VALUE_ID: LEX-P0033-V0007
% #FIELD_VALUE: 检查确认（移液器）
\fieldvalue{\ensuremath{\boxtimes}\ 正常\quad \(\square\)\ 异常}
\\ \hline
\end{tabularx}

\textbf{1.2\quad 试剂耗材}

\begin{tabularx}{\textwidth}{|l|X|X|X|}
\hline
名称 & 厂家 & 批号 & 有效期至 \\ \hline
血糖测试条 &
三诺生物传感股份有限公司 &
% #VALUE_ID: LEX-P0033-V0008
% #FIELD_VALUE: 批号
% #TODO #HANDWRITTEN: 444?440; 字迹不清
\fieldvalue{\handwritten{444?440}} &
% #VALUE_ID: LEX-P0033-V0009
% #FIELD_VALUE: 有效期至
% #HANDWRITTEN: 2021.3.28
\fieldvalue{\handwritten{2021.3.28}}
\\ \hline
\end{tabularx}

\textbf{2\quad 操作步骤：}

\begin{tabularx}{\textwidth}{|p{0.32\textwidth}|X|X|X|}
\hline
\multicolumn{4}{|c|}{操作过程} \\ \hline
取出测试条，取样品自然沉降后的上清液0.6\,\textmu l加入测试条电极反应室，血糖仪自动发出“嘀”提示音，吸入样品完成，血糖仪自动进入倒计时，开始测试； &
\multicolumn{3}{l|}{
取上清液：
% #VALUE_ID: LEX-P0033-V0010
% #FIELD_VALUE: 取上清液
% #HANDWRITTEN: 0.6
\fieldvalue{\handwritten{0.6}}\,\textmu l}
\\ \hline
重复以上步骤2次，共计3次。 &
\multicolumn{3}{l|}{
每个样品分别共计数
% #VALUE_ID: LEX-P0033-V0011
% #FIELD_VALUE: 测试次数
% #HANDWRITTEN: 3
\fieldvalue{\handwritten{3}} 次。}
\\ \hline
\multicolumn{4}{|c|}{检测结果} \\ \hline
\multicolumn{4}{|c|}{葡萄糖浓度（mmol/L）} \\ \hline
\multicolumn{1}{|c|}{第一次} &
\multicolumn{1}{c|}{第二次} &
\multicolumn{1}{c|}{第三次} &
\multicolumn{1}{c|}{平均值（保留一位小数）}
\\ \hline
% #VALUE_ID: LEX-P0033-V0012
% #FIELD_VALUE: 第一次检测结果
% #HANDWRITTEN: 2.17
\fieldvalue{\handwritten{2.17}} &
% #VALUE_ID: LEX-P0033-V0013
% #FIELD_VALUE: 第二次检测结果
% #HANDWRITTEN: 2.19
\fieldvalue{\handwritten{2.19}} &
% #VALUE_ID: LEX-P0033-V0014
% #FIELD_VALUE: 第三次检测结果
% #HANDWRITTEN: 2.17
\fieldvalue{\handwritten{2.17}} &
% #VALUE_ID: LEX-P0033-V0015
% #FIELD_VALUE: 平均值
% #HANDWRITTEN: 2.18
\fieldvalue{\handwritten{2.18}}
\\ \hline
备注 &
% #VALUE_ID: LEX-P0033-V0016
% #HANDWRITTEN: /
\handwritten{/} & \multicolumn{2}{l|}{} \\ \hline
实验人/日期 &
% #VALUE_ID: LEX-P0033-V0017
% #FIELD_VALUE: 实验人
% #TODO #HANDWRITTEN: 姓名不清; 手写签名难以辨认
\fieldvalue{\handwritten{姓名不清}}\quad
% #VALUE_ID: LEX-P0033-V0018
% #FIELD_VALUE: 实验日期
% #HANDWRITTEN: 2024.3.28
\fieldvalue{\handwritten{2024.3.28}} &
\multicolumn{2}{l|}{} \\ \hline
复核人/日期 &
% #VALUE_ID: LEX-P0033-V0019
% #FIELD_VALUE: 复核人
% #TODO #HANDWRITTEN: 姓名不清; 手写签名难以辨认
\fieldvalue{\handwritten{姓名不清}}\quad
% #VALUE_ID: LEX-P0033-V0020
% #FIELD_VALUE: 复核日期
% #HANDWRITTEN: 2024.3.28
\fieldvalue{\handwritten{2024.3.28}} &
\multicolumn{2}{l|}{} \\ \hline
\end{tabularx}
% LEXOID_PAGE_COMPLETED: 33/60

\newpage

\section*{生物反应器细胞液检验操作记录}

\begin{tabularx}{\textwidth}{|p{3cm}|X|p{3cm}|X|}
\hline
样品名称 &
% #VALUE_ID: LEX-P0034-V0001
% #FIELD_VALUE: 样品名称
\fieldvalue{\handwritten{一级生物反应器细胞液}} &
样品规格 &
% #VALUE_ID: LEX-P0034-V0002
% #FIELD_VALUE: 样品规格
% #HANDWRITTEN: 20
\fieldvalue{\handwritten{20}\,ml/袋}
\\ \hline
样品批号 &
% #VALUE_ID: LEX-P0034-V0003
% #FIELD_VALUE: 样品批号
% #TODO #HANDWRITTEN: 231?104?69; 手写内容较模糊
\fieldvalue{\handwritten{231?104?69}} &
取样日期 &
% #VALUE_ID: LEX-P0034-V0004
% #FIELD_VALUE: 取样日期
% #TODO #HANDWRITTEN: 2024年?月8日; 月份字迹不清
\fieldvalue{\handwritten{2024年?月8日}}
\\ \hline
\end{tabularx}

\section{仪器设备、试剂耗材}

\subsection{仪器设备}

\begin{tabularx}{\textwidth}{|p{2.4cm}|p{2.2cm}|p{2.5cm}|p{2.1cm}|p{2.3cm}|X|}
\hline
名称 & 厂家 & 型号 & 设备编号 & 校准有效期至 & 检查确认 \\ \hline
% #VALUE_ID: LEX-P0034-V0005
% #FIELD_VALUE: 名称
\fieldvalue{细胞计数仪} &
% #VALUE_ID: LEX-P0034-V0006
% #FIELD_VALUE: 厂家
\fieldvalue{Nexcelom} &
% #VALUE_ID: LEX-P0034-V0007
% #FIELD_VALUE: 型号
\fieldvalue{Cellometer K2} &
% #VALUE_ID: LEX-P0034-V0008
% #FIELD_VALUE: 设备编号
\fieldvalue{1410905} &
% #VALUE_ID: LEX-P0034-V0009
% #FIELD_VALUE: 校准有效期至
% #HANDWRITTEN: 2024.09.25
\fieldvalue{\handwritten{2024.09.25}} &
% #VALUE_ID: LEX-P0034-V0010
% #FIELD_VALUE: 检查确认
% #HANDWRITTEN: ✓
\fieldvalue{\handwritten{✓}\,正常 \quad □异常}
\\ \hline
% #VALUE_ID: LEX-P0034-V0011
% #FIELD_VALUE: 名称
\fieldvalue{高速冷冻离心机} &
% #VALUE_ID: LEX-P0034-V0012
% #FIELD_VALUE: 厂家
\fieldvalue{Thermo} &
% #VALUE_ID: LEX-P0034-V0013
% #FIELD_VALUE: 型号
\fieldvalue{ST16R} &
% #VALUE_ID: LEX-P0034-V0014
% #FIELD_VALUE: 设备编号
\fieldvalue{1410810} &
% #VALUE_ID: LEX-P0034-V0015
% #FIELD_VALUE: 校准有效期至
% #HANDWRITTEN: 2024.10.18
\fieldvalue{\handwritten{2024.10.18}} &
% #VALUE_ID: LEX-P0034-V0016
% #FIELD_VALUE: 检查确认
% #HANDWRITTEN: ✓
\fieldvalue{\handwritten{✓}\,正常 \quad □异常}
\\ \hline
% #VALUE_ID: LEX-P0034-V0017
% #FIELD_VALUE: 名称
\fieldvalue{电热恒温水槽} &
% #VALUE_ID: LEX-P0034-V0018
% #FIELD_VALUE: 厂家
\fieldvalue{上海一恒} &
% #VALUE_ID: LEX-P0034-V0019
% #FIELD_VALUE: 型号
\fieldvalue{DK-8AXX} &
% #VALUE_ID: LEX-P0034-V0020
% #FIELD_VALUE: 设备编号
\fieldvalue{1411001} &
% #VALUE_ID: LEX-P0034-V0021
% #FIELD_VALUE: 校准有效期至
% #HANDWRITTEN: 2024.10.22
\fieldvalue{\handwritten{2024.10.22}} &
% #VALUE_ID: LEX-P0034-V0022
% #FIELD_VALUE: 检查确认
% #HANDWRITTEN: ✓
\fieldvalue{\handwritten{✓}\,正常 \quad □异常}
\\ \hline
\end{tabularx}

\subsection{试剂耗材}

\begin{tabularx}{\textwidth}{|p{3.3cm}|p{2.8cm}|X|p{2.6cm}|}
\hline
名称 & 厂家 & 批号 & 有效期至 \\ \hline
% #VALUE_ID: LEX-P0034-V0023
% #FIELD_VALUE: 名称
\fieldvalue{AO/PI 染液} &
% #VALUE_ID: LEX-P0034-V0024
% #FIELD_VALUE: 厂家
\fieldvalue{Nexcelom} &
% #VALUE_ID: LEX-P0034-V0025
% #FIELD_VALUE: 批号
% #HANDWRITTEN: 324?49
\fieldvalue{\handwritten{324?49}} &
% #VALUE_ID: LEX-P0034-V0026
% #FIELD_VALUE: 有效期至
% #HANDWRITTEN: 2026.07.22
\fieldvalue{\handwritten{2026.07.22}}
\\ \hline
% #VALUE_ID: LEX-P0034-V0027
% #FIELD_VALUE: 名称
\fieldvalue{0.9\%氯化钠注射液} &
% #VALUE_ID: LEX-P0034-V0028
% #FIELD_VALUE: 厂家
\fieldvalue{石家庄四药} &
% #VALUE_ID: LEX-P0034-V0029
% #FIELD_VALUE: 批号
% #HANDWRITTEN: 24?375?20
\fieldvalue{\handwritten{24?375?20}} &
% #VALUE_ID: LEX-P0034-V0030
% #FIELD_VALUE: 有效期至
% #HANDWRITTEN: 2026.10.06
\fieldvalue{\handwritten{2026.10.06}}
\\ \hline
% #VALUE_ID: LEX-P0034-V0031
% #FIELD_VALUE: 名称
\fieldvalue{20\%人血白蛋白} &
% #VALUE_ID: LEX-P0034-V0032
% #FIELD_VALUE: 厂家
\fieldvalue{Baxter AG} &
% #VALUE_ID: LEX-P0034-V0033
% #FIELD_VALUE: 批号
% #HANDWRITTEN: AB116?9
\fieldvalue{\handwritten{AB116?9}} &
% #VALUE_ID: LEX-P0034-V0034
% #FIELD_VALUE: 有效期至
% #TODO #HANDWRITTEN: 2027.?.?; 日期字迹不清
\fieldvalue{\handwritten{2027.?.?}}
\\ \hline
\end{tabularx}

\section{操作步骤}

\subsection*{溶液准备}

取 5\,ml 20\%人血白蛋白加入到 95\,ml 的 0.9\%氯化钠注射液中，混匀即得。

\begin{tabularx}{\textwidth}{|p{3.2cm}|X|p{3.2cm}|}
\hline
名称 & 配制批号 & 有效期至 \\ \hline
% #VALUE_ID: LEX-P0034-V0035
% #FIELD_VALUE: 名称
\fieldvalue{5$\times$裂解液} &
% #VALUE_ID: LEX-P0034-V0036
% #FIELD_VALUE: 配制批号
% #HANDWRITTEN: 24?10?80
\fieldvalue{\handwritten{24?10?80}} &
% #VALUE_ID: LEX-P0034-V0037
% #FIELD_VALUE: 有效期至
% #HANDWRITTEN: 2024.12.7
\fieldvalue{\handwritten{2024.12.7}}
\\ \hline
\end{tabularx}

\subsection*{操作过程}

\begin{tabularx}{\textwidth}{|X|p{5.5cm}|}
\hline
操作步骤 & 操作记录 \\ \hline
按照 50\,ml 离心管的刻度读取样品体积，在取样管中加入对应体积的 5$\times$裂解液（每 4\,ml 样品加入 1\,ml 5$\times$裂解液，按照“只进不合”原则，如 26\,ml 样品，加入 7\,ml 5$\times$裂解液），充分混匀后倒入至 2.3 的样品中，置于 37\,\textdegree C 电热恒温水槽中裂解 20\,min 后，取出离心管混匀，使样品充分裂解；再置于 37\,\textdegree C 电热恒温水槽中裂解 10\,min 得到细胞悬液。 &
样品体积：
% #VALUE_ID: LEX-P0034-V0038
% #FIELD_VALUE: 样品体积
% #HANDWRITTEN: 24
\fieldvalue{\handwritten{24}}\,ml

裂解液：
% #VALUE_ID: LEX-P0034-V0039
% #FIELD_VALUE: 裂解液
% #HANDWRITTEN: 6
\fieldvalue{\handwritten{6}}\,ml

电热恒温水槽：
% #VALUE_ID: LEX-P0034-V0040
% #FIELD_VALUE: 电热恒温水槽
% #HANDWRITTEN: 37
\fieldvalue{\handwritten{37}}\,\textdegree C

裂解：
% #VALUE_ID: LEX-P0034-V0041
% #FIELD_VALUE: 裂解
% #TODO #HANDWRITTEN: 9?秒; 手写记录不清
\fieldvalue{\handwritten{9?秒}}
\\ \hline
取出裂解好的细胞悬液，加入对应体积的 1\%人白蛋白，混匀。 &
% #VALUE_ID: LEX-P0034-V0042
% #FIELD_VALUE: 离心条件
% #TODO #HANDWRITTEN: 30?×g 离心1分钟; 手写内容部分不清
\fieldvalue{\handwritten{30?}\(\times g\)}\,离心
% #VALUE_ID: LEX-P0034-V0043
% #FIELD_VALUE: 离心时间
% #HANDWRITTEN: 1
\fieldvalue{\handwritten{1}}\,分钟

加 1\% 人白蛋白溶液重悬体积：
% #VALUE_ID: LEX-P0034-V0044
% #FIELD_VALUE: 重悬体积
% #HANDWRITTEN: 2
\fieldvalue{\handwritten{2}}\,ml

吸取：
% #VALUE_ID: LEX-P0034-V0045
% #FIELD_VALUE: 吸取体积
% #HANDWRITTEN: 20
\fieldvalue{\handwritten{20}}\,\textmu l
\\ \hline
\end{tabularx}

% TODO: The operation record continues beyond the bottom boundary of this page.
% LEXOID_PAGE_COMPLETED: 34/60

\newpage

\begin{center}
\textbf{细胞计数及细胞活率检测记录}
\end{center}

\noindent
\begin{tabularx}{\textwidth}{@{}lXr@{}}
 & & 文件编号：%
% #VALUE_ID: LEX-P0035-V0001
% #FIELD_VALUE: 文件编号
% #TODO #HANDWRITTEN: REC-YZ-SOP-CQ-99-006; printed value is partially unclear
\fieldvalue{\handwritten{REC-YZ-SOP-CQ-99-006}}\\
\end{tabularx}

\noindent
加入 22 ml 1\% 人白蛋白稀释液，反复吹打至单细胞悬液；\\
摇去计数板上下两层保护膜，将细胞悬液混匀，从中吸取 20\,$\mu$l\\
与 20\,$\mu$l AO/PI 染料混合均匀。\\
吸取 20\,$\mu$l 细胞悬液与 AO/PI 染料的混合液，注入到计数板的一\\
个检测室中。\\
将加入混合液的一端插入进样口。\\
选择 Assay 类型，使用 AO/PI 染料进行双荧光存活率检测；\\
在主界面点击“阻冲输入样本”。\\
点击主界面左下角“Preview”预览细胞图片。\\
调整仪器焦距直至观察到荧光通道下的细胞中心透亮，轮廓清\\
晰。\\
点击“count”按钮进行计数并自动拍摄图片。\\
重复以上步骤 2 次，共计数 3 次。打印数据结果。

\medskip

\noindent
\begin{tabularx}{\textwidth}{@{}lX@{}}
与 & %
% #VALUE_ID: LEX-P0035-V0002
% #FIELD_VALUE: 与 AO/PI 染料混合的细胞悬液体积
% #TODO #HANDWRITTEN: 20?; handwritten volume is unclear
\fieldvalue{\handwritten{20?}} $\mu$l AO/PI 染料混合\\[2pt]
吸取 & %
% #VALUE_ID: LEX-P0035-V0003
% #FIELD_VALUE: 吸取体积
% #HANDWRITTEN: 22
\fieldvalue{\handwritten{22}} $\mu$l 计数。\\[2pt]
每个样品均计数 & %
% #VALUE_ID: LEX-P0035-V0004
% #FIELD_VALUE: 每个样品计数次数
% #HANDWRITTEN: 2
\fieldvalue{\handwritten{2}} 次。
\end{tabularx}

\medskip

\noindent
\begin{tabularx}{\textwidth}{@{}lX@{}}
数据保存路径：&
% #VALUE_ID: LEX-P0035-V0005
% #FIELD_VALUE: 数据保存路径
% #TODO #HANDWRITTEN: 路径不可辨认; blue handwritten text is illegible
\fieldvalue{\handwritten{路径不可辨认}}
\end{tabularx}

\begin{center}
检测结果
\end{center}

\noindent
检测结果：$\square$ 一级反应　$\square$ 二级反应　%
% #VALUE_ID: LEX-P0035-V0006
% #FIELD_VALUE: 检测结果选择
% #HANDWRITTEN: ✓（三级反应）
\fieldvalue{\handwritten{\checkmark\ （三级反应）}}

\medskip

\small
\renewcommand{\arraystretch}{1.35}
\begin{tabularx}{\textwidth}{|c|X|X|X|X|X|X|}
\hline
检测次 &
活细胞浓度（个细胞/ml） &
样品密度（个细胞/ml） &
细胞密度 CV 值（\%） &
细胞活率（\%） &
样品活率（\%） &
细胞活率 CV 值（\%）\\
\hline
第一次 &
% #VALUE_ID: LEX-P0035-V0007
% #FIELD_VALUE: 第一次检测的活细胞浓度
% #TODO #HANDWRITTEN: 1.3?×10^6; handwritten digits are partially unclear
\fieldvalue{\handwritten{1.3?$\times10^{6}$}} &
% #VALUE_ID: LEX-P0035-V0008
% #FIELD_VALUE: 样品密度
% #TODO #HANDWRITTEN: 1.9?×10^6; handwritten value is partially unclear
\fieldvalue{\handwritten{1.9?$\times10^{6}$}} &
% #VALUE_ID: LEX-P0035-V0009
% #FIELD_VALUE: 细胞密度 CV 值
% #HANDWRITTEN: 8
\fieldvalue{\handwritten{8}} &
% #VALUE_ID: LEX-P0035-V0010
% #FIELD_VALUE: 第一次检测的细胞活率
% #HANDWRITTEN: 87.2
\fieldvalue{\handwritten{87.2}} &
% #VALUE_ID: LEX-P0035-V0011
% #FIELD_VALUE: 样品活率
% #HANDWRITTEN: 87
\fieldvalue{\handwritten{87}} &
% #VALUE_ID: LEX-P0035-V0012
% #FIELD_VALUE: 细胞活率 CV 值
% #HANDWRITTEN: 3
\fieldvalue{\handwritten{3}}\\
\hline
第二次 &
% #VALUE_ID: LEX-P0035-V0013
% #FIELD_VALUE: 第二次检测的活细胞浓度
% #TODO #HANDWRITTEN: 1.3?×10^6; handwritten digits are partially unclear
\fieldvalue{\handwritten{1.3?$\times10^{6}$}} &
&
&
% #VALUE_ID: LEX-P0035-V0014
% #FIELD_VALUE: 第二次检测的细胞活率
% #HANDWRITTEN: 88.2
\fieldvalue{\handwritten{88.2}} &
&\\
\hline
第三次 &
% #VALUE_ID: LEX-P0035-V0015
% #FIELD_VALUE: 第三次检测的活细胞浓度
% #TODO #HANDWRITTEN: 1.5?×10^6; handwritten digits are partially unclear
\fieldvalue{\handwritten{1.5?$\times10^{6}$}} &
&
&
% #VALUE_ID: LEX-P0035-V0016
% #FIELD_VALUE: 第三次检测的细胞活率
% #HANDWRITTEN: 88.8
\fieldvalue{\handwritten{88.8}} &
&\\
\hline
\end{tabularx}
\normalsize

\noindent
\begin{tabularx}{\textwidth}{|l|X|l|X|}
\hline
发酵罐细胞体积（ml） &
% #VALUE_ID: LEX-P0035-V0017
% #FIELD_VALUE: 发酵罐细胞体积
% #TODO #HANDWRITTEN: 2???; handwritten value is unclear
\fieldvalue{\handwritten{2???}} &
细胞数量（个细胞/罐） &
% #VALUE_ID: LEX-P0035-V0018
% #FIELD_VALUE: 细胞数量
% #TODO #HANDWRITTEN: 2.18×10^8; exponent and final digit are somewhat unclear
\fieldvalue{\handwritten{2.18$\times10^{8}$}}\\
\hline
\multicolumn{4}{|l|}{计算要求：样品密度为 3 次活细胞密度结果取平均值，细胞浓度以 1\% 人白蛋白稀释液总体积计，再除以生物反应器}\\
\multicolumn{4}{|l|}{细胞液体积；细胞数量为样品密度乘以发酵罐体积。}\\
\hline
备注 &
% #VALUE_ID: LEX-P0035-V0019
% #FIELD_VALUE: 备注
% #TODO #HANDWRITTEN: /; handwritten mark has no legible text
\fieldvalue{\handwritten{/}} &
&\\
\hline
实验人/日期 &
% #VALUE_ID: LEX-P0035-V0020
% #FIELD_VALUE: 实验人
% #TODO #HANDWRITTEN: 签名不可辨认; handwritten signature is illegible
\fieldvalue{\handwritten{签名不可辨认}}%
% #VALUE_ID: LEX-P0035-V0021
% #FIELD_VALUE: 实验日期
% #TODO #HANDWRITTEN: 2025.8.?; final digit is unclear
\fieldvalue{\handwritten{2025.8.?}} &
&\\
\hline
复核人/日期 &
% #VALUE_ID: LEX-P0035-V0022
% #FIELD_VALUE: 复核人
% #TODO #HANDWRITTEN: 签名不可辨认; handwritten signature is illegible
\fieldvalue{\handwritten{签名不可辨认}}%
% #VALUE_ID: LEX-P0035-V0023
% #FIELD_VALUE: 复核日期
% #HANDWRITTEN: 2025.8.28
\fieldvalue{\handwritten{2025.8.28}} &
&\\
\hline
\end{tabularx}

% TODO: The page contains a continuation-style laboratory form; only the rows and fields visible on this page are rendered.
% LEXOID_PAGE_COMPLETED: 35/60

\newpage

\begin{center}
\textbf{生物反应器细胞液检测记录}
\end{center}

文件编号：\fieldvalue{REC-YZ-SOP-QC-99-006-a}

\begin{tabularx}{\textwidth}{|l|X|l|X|}
\hline
样品名称 &
% #VALUE_ID: LEX-P0036-V0002
% #FIELD_VALUE: 样品名称
\fieldvalue{绿生物反应器细胞液} &
样品规格 &
% #VALUE_ID: LEX-P0036-V0003
% #FIELD_VALUE: 样品规格
% #HANDWRITTEN: 20
\fieldvalue{\handwritten{20}} ml/袋
\\ \hline
样品批号 &
% #VALUE_ID: LEX-P0036-V0004
% #FIELD_VALUE: 样品批号
% #HANDWRITTEN: (8)23?5?070?
\fieldvalue{\handwritten{(8)23?5?070?}} &
取样日期 &
% #VALUE_ID: LEX-P0036-V0005
% #FIELD_VALUE: 取样日期
% #HANDWRITTEN: 2024年?月18日
\fieldvalue{\handwritten{2024年?月18日}}
\\ \hline
\end{tabularx}

\section{仪器设备及试剂}

\subsection{仪器设备}

\begin{tabularx}{\textwidth}{|X|X|X|X|X|}
\hline
仪器名称 & 厂家 & 型号 & 设备编号 & 检查确认 \\ \hline
% #VALUE_ID: LEX-P0036-V0006
% #FIELD_VALUE: 仪器名称
\fieldvalue{荧光显微镜} &
% #VALUE_ID: LEX-P0036-V0007
% #FIELD_VALUE: 厂家
\fieldvalue{徕卡} &
% #VALUE_ID: LEX-P0036-V0008
% #FIELD_VALUE: 型号
\fieldvalue{DmiL LED} &
% #VALUE_ID: LEX-P0036-V0009
% #FIELD_VALUE: 设备编号
\fieldvalue{1431404} &
% #VALUE_ID: LEX-P0036-V0010
% #FIELD_VALUE: 检查确认
\fieldvalue{\(\boxtimes\)正常 \(\square\)异常}
\\ \hline
\end{tabularx}

\subsection{试剂耗材}

\begin{tabularx}{\textwidth}{|X|X|X|X|X|}
\hline
名称 & 厂家 & 货号 & 批号 & 有效期至 \\ \hline
% #VALUE_ID: LEX-P0036-V0011
% #FIELD_VALUE: 名称
\fieldvalue{细胞活力检测试剂盒} &
% #VALUE_ID: LEX-P0036-V0012
% #FIELD_VALUE: 厂家
\fieldvalue{江苏凯基生物技术股份有限公司} &
% #VALUE_ID: LEX-P0036-V0013
% #FIELD_VALUE: 货号
\fieldvalue{KGA9501-1000} &
% #VALUE_ID: LEX-P0036-V0014
% #FIELD_VALUE: 批号
% #HANDWRITTEN: 204?40?
\fieldvalue{\handwritten{204?40?}} &
% #VALUE_ID: LEX-P0036-V0015
% #FIELD_VALUE: 有效期至
% #HANDWRITTEN: 2027.10.07
\fieldvalue{\handwritten{2027.10.07}}
\\ \hline
\end{tabularx}

\section{操作步骤}

\begin{tabularx}{\textwidth}{|p{0.48\textwidth}|X|}
\hline
\multicolumn{2}{|c|}{\textbf{溶液配制}} \\ \hline
20$\times$的荧光染液配制：取
% #VALUE_ID: LEX-P0036-V0016
% #FIELD_VALUE: 组分A用量
% #HANDWRITTEN: 0.2
\fieldvalue{\handwritten{0.2}} $\mu$l 组分A 和
% #VALUE_ID: LEX-P0036-V0017
% #FIELD_VALUE: 组分B用量
% #HANDWRITTEN: 20
\fieldvalue{\handwritten{20}} $\mu$l 组分B 与
% #VALUE_ID: LEX-P0036-V0018
% #FIELD_VALUE: PBS用量
% #HANDWRITTEN: 396
\fieldvalue{\handwritten{396}} $\mu$l PBS 混合均匀配成20$\times$的荧光染液。 &
\\ \hline
\multicolumn{2}{|c|}{\textbf{操作过程}} \\ \hline
取配好后的样品200$\mu$l至96孔板中，取2孔； &
样品
% #VALUE_ID: LEX-P0036-V0019
% #FIELD_VALUE: 样品加入量
% #HANDWRITTEN: 200
\fieldvalue{\handwritten{200}} $\mu$l/孔\\
每个样品取
% #VALUE_ID: LEX-P0036-V0020
% #FIELD_VALUE: 每个样品孔数
% #HANDWRITTEN: 2
\fieldvalue{\handwritten{2}} 孔
\\ \hline
自然沉降后，吸弃上清，注意不要吸到沉淀样品； &
% #VALUE_ID: LEX-P0036-V0021
% #FIELD_VALUE: 操作确认
\fieldvalue{\(\boxtimes\)吸弃上清}
\\ \hline
取20$\times$的荧光染液加入到96孔板中，200$\mu$l/孔，避光孵育30--60min； &
20$\times$的荧光染液
% #VALUE_ID: LEX-P0036-V0022
% #FIELD_VALUE: 荧光染液加入量
% #HANDWRITTEN: 20
\fieldvalue{\handwritten{20}} $\mu$l/孔\\
避光孵育
% #VALUE_ID: LEX-P0036-V0023
% #FIELD_VALUE: 避光孵育条件
% #HANDWRITTEN: 37℃
\fieldvalue{\handwritten{37℃}}
\\ \hline
染色结束后，吸弃上清，取PBS加入到以上96孔板中，200$\mu$l/孔，吸弃上清，注意不要吸到沉淀样品，清洗3次，之后每孔加入100$\mu$l PBS进行细胞检测。 &
加入PBS
% #VALUE_ID: LEX-P0036-V0024
% #FIELD_VALUE: PBS加入量
% #HANDWRITTEN: 200
\fieldvalue{\handwritten{200}} $\mu$l/孔\\
清洗
% #VALUE_ID: LEX-P0036-V0025
% #FIELD_VALUE: 清洗次数
% #HANDWRITTEN: 2
\fieldvalue{\handwritten{2}} 次
\\ \hline
将荧光显微镜调至4$\times$物镜，软件中调节节至相应物镜参数，调节到适合的光源和修距，拍摄4张清晰图像。\\
选一张清晰图像，打印并粘贴在实验记录上。 &
\\ \hline
数据保存路径 &
% #VALUE_ID: LEX-P0036-V0026
% #FIELD_VALUE: 数据保存路径
% #HANDWRITTEN: I4\textbackslash 细胞检测\textbackslash DMS\textbackslash 设备图
\fieldvalue{\handwritten{I4\textbackslash 细胞检测\textbackslash DMS\textbackslash 设备图}}
\\ \hline
\multicolumn{2}{|c|}{检测结果} \\ \hline
\end{tabularx}

% TODO：页面底部“检测结果”区域在本页仅显示标题，后续内容未见于本页。

\begin{center}
第4页 共5页
\end{center}
% LEXOID_PAGE_COMPLETED: 36/60

\newpage

\noindent
\begin{tabularx}{\textwidth}{@{}Xr@{}}
\textbf{钛生物材料科技（北京）有限公司} &
\begin{tabular}{|c|}
\hline
留存生物\\
% #VALUE_ID: LEX-P0037-V0001
% #FIELD_VALUE: 文档类型
\fieldvalue{原件}\\
\hline
\end{tabular}
\quad
\begin{tabular}{|c|}
\hline
管理性文件\\
版本号：%
% #VALUE_ID: LEX-P0037-V0002
% #FIELD_VALUE: 版本号
\fieldvalue{4}\\
\hline
\end{tabular}
\\
\end{tabularx}

\noindent
文件编码：REC-YZ-SOR-QC-99-006-a

\vspace{0.5em}
\noindent\rule{\textwidth}{0.4pt}

\vfill

\begin{center}
% TODO: 插入本页中部可见的显微图像；原始图像文件名不可用。
\fbox{\parbox[c][3.0cm][c]{0.38\textwidth}{\centering 显微图像占位符}}

% #VALUE_ID: LEX-P0037-V0003
% #HANDWRITTEN: 阴性
\handwritten{阴性}
\quad
% #VALUE_ID: LEX-P0037-V0004
% #HANDWRITTEN: 2016/04/08
\handwritten{2016/04/08}
\end{center}

\vfill

\noindent
\begin{tabularx}{\textwidth}{|p{0.16\textwidth}|X|}
\hline
备注 &
% #VALUE_ID: LEX-P0037-V0005
% #FIELD_VALUE: 备注
% #HANDWRITTEN: ✓
\fieldvalue{\handwritten{✓}}\\
\hline
实验人/日期 &
% #VALUE_ID: LEX-P0037-V0006
% #FIELD_VALUE: 实验人
% #HANDWRITTEN: 姓名不清
\fieldvalue{\handwritten{姓名不清}}
\quad
% #VALUE_ID: LEX-P0037-V0007
% #FIELD_VALUE: 实验日期
% #HANDWRITTEN: 2016.4.28
\fieldvalue{\handwritten{2016.4.28}}\\
\hline
复核人/日期 &
% #VALUE_ID: LEX-P0037-V0008
% #FIELD_VALUE: 复核人
% #HANDWRITTEN: 姓名不清
\fieldvalue{\handwritten{姓名不清}}
\quad
% #VALUE_ID: LEX-P0037-V0009
% #FIELD_VALUE: 复核日期
% #HANDWRITTEN: 2016.08.28
\fieldvalue{\handwritten{2016.08.28}}\\
\hline
\end{tabularx}

\begin{center}
第 5 页 共 5 页
\end{center}
% LEXOID_PAGE_COMPLETED: 37/60

\newpage

\noindent Assay: %
% #VALUE_ID: LEX-P0038-V0001
% #FIELD_VALUE: Assay
\fieldvalue{MSC-1}

\medskip
\noindent Cell Type F1: %
% #VALUE_ID: LEX-P0038-V0002
% #FIELD_VALUE: Cell Type F1
\fieldvalue{MSC (AO)}

\noindent Cell Type F2: %
% #VALUE_ID: LEX-P0038-V0003
% #FIELD_VALUE: Cell Type F2
\fieldvalue{MSC (PI)}

\medskip
\noindent Sample ID: %
% #VALUE_ID: LEX-P0038-V0004
% #FIELD_VALUE: Sample ID
\fieldvalue{QC-JS-QY0401-A31Z2202508008-20250828-51H-1}

\noindent Dilution: %
% #VALUE_ID: LEX-P0038-V0005
% #FIELD_VALUE: Dilution
\fieldvalue{2.00}

\bigskip
\noindent\rule{\textwidth}{0.4pt}

\noindent Results:

\medskip
\begin{tabular}{lll}
Count & Concentration & Mean Diameter \\
\texttt{============} & \texttt{============} & \texttt{=============} \\

Total: %
% #VALUE_ID: LEX-P0038-V0006
% #FIELD_VALUE: Total count
\fieldvalue{451}
&
% #VALUE_ID: LEX-P0038-V0007
% #FIELD_VALUE: Total concentration
\fieldvalue{$1.60{\times}10^{6}$ cells/mL}
&
% #VALUE_ID: LEX-P0038-V0008
% #FIELD_VALUE: Total mean diameter
\fieldvalue{15.0 micron}
\\

Live: %
% #VALUE_ID: LEX-P0038-V0009
% #FIELD_VALUE: Live count
\fieldvalue{384}
&
% #VALUE_ID: LEX-P0038-V0010
% #FIELD_VALUE: Live concentration
\fieldvalue{$1.36{\times}10^{6}$ cells/mL}
&
% #VALUE_ID: LEX-P0038-V0011
% #FIELD_VALUE: Live mean diameter
\fieldvalue{15.0 microns}
\\

Dead: %
% #VALUE_ID: LEX-P0038-V0012
% #FIELD_VALUE: Dead count
\fieldvalue{67}
&
% #VALUE_ID: LEX-P0038-V0013
% #FIELD_VALUE: Dead concentration
\fieldvalue{$2.37{\times}10^{5}$ cells/mL}
&
% #VALUE_ID: LEX-P0038-V0014
% #FIELD_VALUE: Dead mean diameter
\fieldvalue{15.3 micron}
\\

Viability: %
% #VALUE_ID: LEX-P0038-V0015
% #FIELD_VALUE: Viability
\fieldvalue{85.2\%}
&&
\end{tabular}

\bigskip

\begin{center}
\begin{tabular}{cc}
% TODO: Image visible on this page; no source filename is available.
\fbox{\rule{0pt}{0.26\textheight}\rule{0.42\textwidth}{0pt}}
&
% TODO: Image visible on this page; no source filename is available.
\fbox{\rule{0pt}{0.26\textheight}\rule{0.42\textwidth}{0pt}}
\\[1em]
% TODO: Image visible on this page; no source filename is available.
\fbox{\rule{0pt}{0.26\textheight}\rule{0.42\textwidth}{0pt}}
&
% TODO: Image visible on this page; no source filename is available.
\fbox{\rule{0pt}{0.26\textheight}\rule{0.42\textwidth}{0pt}}
\end{tabular}
\end{center}

\vfill

\noindent 操作人/审核人：%
% #VALUE_ID: LEX-P0038-V0016
% #FIELD_VALUE: 操作人/审核人
% #TODO #HANDWRITTEN: illegible signature; handwritten signature is not confidently recognizable
\fieldvalue{\handwritten{[illegible signature]}}
\hfill
日期：%
% #VALUE_ID: LEX-P0038-V0017
% #FIELD_VALUE: 日期
% #HANDWRITTEN: 2024.8.28
\fieldvalue{\handwritten{2024.8.28}}
% LEXOID_PAGE_COMPLETED: 38/60

\newpage

\noindent
% #VALUE_ID: LEX-P0039-V0001
% #FIELD_VALUE: Assay
\textbf{Assay:} \fieldvalue{MSC-1}

\medskip

% #VALUE_ID: LEX-P0039-V0002
% #FIELD_VALUE: Cell Type F1
\textbf{Cell Type F1:} \fieldvalue{MSC (AO)}

% #VALUE_ID: LEX-P0039-V0003
% #FIELD_VALUE: Cell Type F2
\textbf{Cell Type F2:} \fieldvalue{MSC (PI)}

\medskip

% #VALUE_ID: LEX-P0039-V0004
% #FIELD_VALUE: Sample ID
\textbf{Sample ID:} \fieldvalue{QC-JS-QY0401-A31Z2202508008-20250828-51H-2}

% #VALUE_ID: LEX-P0039-V0005
% #FIELD_VALUE: Dilution
\textbf{Dilution:} \fieldvalue{2.00}

\medskip
\hrule
\medskip

\textbf{Results:}

\medskip

\small
\begin{tabularx}{\textwidth}{@{}l X X@{}}
\textbf{Count} & \textbf{Concentration} & \textbf{Mean Diameter} \\
\underline{\hspace{4.5em}} & \underline{\hspace{8em}} & \underline{\hspace{8em}} \\[2pt]

% #VALUE_ID: LEX-P0039-V0006
% #FIELD_VALUE: Total count
Total: \fieldvalue{441}
&
% #VALUE_ID: LEX-P0039-V0007
% #FIELD_VALUE: Total concentration
\fieldvalue{$1.57\times10^{6}$ cells/mL}
&
% #VALUE_ID: LEX-P0039-V0008
% #FIELD_VALUE: Total mean diameter
\fieldvalue{15.8 micron}
\\

% #VALUE_ID: LEX-P0039-V0009
% #FIELD_VALUE: Live count
Live: \fieldvalue{388}
&
% #VALUE_ID: LEX-P0039-V0010
% #FIELD_VALUE: Live concentration
\fieldvalue{$1.38\times10^{6}$ cells/mL}
&
% #VALUE_ID: LEX-P0039-V0011
% #FIELD_VALUE: Live mean diameter
\fieldvalue{15.9 microns}
\\

% #VALUE_ID: LEX-P0039-V0012
% #FIELD_VALUE: Dead count
Dead: \fieldvalue{53}
&
% #VALUE_ID: LEX-P0039-V0013
% #FIELD_VALUE: Dead concentration
\fieldvalue{$1.87\times10^{5}$ cells/mL}
&
% #VALUE_ID: LEX-P0039-V0014
% #FIELD_VALUE: Dead mean diameter
\fieldvalue{14.5 microns}
\\

% #VALUE_ID: LEX-P0039-V0015
% #FIELD_VALUE: Viability
Viability: \fieldvalue{88.1\%}
&
&
\end{tabularx}
\normalsize

\medskip

\begin{center}
\begin{tabular}{cc}
% TODO: Image visible on this page; source filename unavailable.
\fbox{\rule{0pt}{0.28\textheight}\rule{0.43\textwidth}{0pt}}
&
% TODO: Image visible on this page; source filename unavailable.
\fbox{\rule{0pt}{0.28\textheight}\rule{0.43\textwidth}{0pt}}
\\[1em]
% TODO: Image visible on this page; source filename unavailable.
\fbox{\rule{0pt}{0.28\textheight}\rule{0.43\textwidth}{0pt}}
&
% TODO: Image visible on this page; source filename unavailable.
\fbox{\rule{0pt}{0.28\textheight}\rule{0.43\textwidth}{0pt}}
\end{tabular}
\end{center}

\vfill

\noindent
操作人/审核人：

% #VALUE_ID: LEX-P0039-V0016
% #FIELD_VALUE: 操作人/审核人
% #TODO #HANDWRITTEN: [illegible signature]; handwritten signature is not clearly legible.
\fieldvalue{\handwritten{[illegible signature]}}

\hfill
日期：

% #VALUE_ID: LEX-P0039-V0017
% #FIELD_VALUE: 日期
% #HANDWRITTEN: 2025.08.28
\fieldvalue{\handwritten{2025.08.28}}
% LEXOID_PAGE_COMPLETED: 39/60

\newpage

\noindent\textbf{Assay:}
% #VALUE_ID: LEX-P0040-V0001
% #FIELD_VALUE: Assay
\fieldvalue{MSC-1}

\medskip
\noindent\textbf{Cell Type F1:}
% #VALUE_ID: LEX-P0040-V0002
% #FIELD_VALUE: Cell Type F1
\fieldvalue{MSC (AO)}

\noindent\textbf{Cell Type F2:}
% #VALUE_ID: LEX-P0040-V0003
% #FIELD_VALUE: Cell Type F2
\fieldvalue{MSC (PI)}

\medskip
\noindent\textbf{Sample ID:}
% #VALUE_ID: LEX-P0040-V0004
% #FIELD_VALUE: Sample ID
\fieldvalue{QC-JS-QY0401-A3172202508008-20250828-51H-3}

\noindent\textbf{Dilution:}
% #VALUE_ID: LEX-P0040-V0005
% #FIELD_VALUE: Dilution
\fieldvalue{2.00}

\medskip
\noindent\rule{\textwidth}{0.4pt}

\noindent Results:

\medskip
\small
\begin{tabularx}{\textwidth}{@{}X X X@{}}
\textbf{Count} & \textbf{Concentration} & \textbf{Mean Diameter} \\
\underline{\hspace{1.2cm}} & \underline{\hspace{1.8cm}} & \underline{\hspace{2cm}} \\[2pt]

Total:
% #VALUE_ID: LEX-P0040-V0006
% #FIELD_VALUE: Total count
\fieldvalue{381}
&
% #VALUE_ID: LEX-P0040-V0007
% #FIELD_VALUE: Total concentration
\fieldvalue{$1.35\times10^{6}\ \mathrm{cells/mL}$}
&
% #VALUE_ID: LEX-P0040-V0008
% #FIELD_VALUE: Total mean diameter
\fieldvalue{$15.6\ \mathrm{micron}$}
\\

Live:
% #VALUE_ID: LEX-P0040-V0009
% #FIELD_VALUE: Live count
\fieldvalue{338}
&
% #VALUE_ID: LEX-P0040-V0010
% #FIELD_VALUE: Live concentration
\fieldvalue{$1.20\times10^{6}\ \mathrm{cells/mL}$}
&
% #VALUE_ID: LEX-P0040-V0011
% #FIELD_VALUE: Live mean diameter
\fieldvalue{$15.9\ \mathrm{micron}$}
\\

Dead:
% #VALUE_ID: LEX-P0040-V0012
% #FIELD_VALUE: Dead count
\fieldvalue{43}
&
% #VALUE_ID: LEX-P0040-V0013
% #FIELD_VALUE: Dead concentration
\fieldvalue{$1.51\times10^{5}\ \mathrm{cells/mL}$}
&
% #VALUE_ID: LEX-P0040-V0014
% #FIELD_VALUE: Dead mean diameter
\fieldvalue{$13.1\ \mathrm{micron}$}
\\

Viability:
% #VALUE_ID: LEX-P0040-V0015
% #FIELD_VALUE: Viability
\fieldvalue{88.8\%}
&
&
\end{tabularx}
\normalsize

\medskip
\begin{center}
\begin{tabular}{@{}c@{\hspace{1.2cm}}c@{}}
% TODO: Image placeholder for upper-left microscopy image; no filename is available.
\fbox{\rule{0pt}{0.22\textwidth}\rule{0.38\textwidth}{0pt}}
&
% TODO: Image placeholder for upper-right microscopy image; no filename is available.
\fbox{\rule{0pt}{0.22\textwidth}\rule{0.38\textwidth}{0pt}}
\\[8pt]
% TODO: Image placeholder for lower-left microscopy image; no filename is available.
\fbox{\rule{0pt}{0.22\textwidth}\rule{0.38\textwidth}{0pt}}
&
% TODO: Image placeholder for lower-right microscopy image; no filename is available.
\fbox{\rule{0pt}{0.22\textwidth}\rule{0.38\textwidth}{0pt}}
\end{tabular}
\end{center}

\vfill

\noindent 操作人/审核人：
% #VALUE_ID: LEX-P0040-V0016
% #FIELD_VALUE: 操作人/审核人
% #TODO #HANDWRITTEN: illegible signature; handwritten signature is not clearly readable
\fieldvalue{\handwritten{[签名]}}

\hfill
日期：
% #VALUE_ID: LEX-P0040-V0017
% #FIELD_VALUE: 日期
% #HANDWRITTEN: 2025.7.18
\fieldvalue{\handwritten{2025.7.18}}
% LEXOID_PAGE_COMPLETED: 40/60

\newpage

\begin{center}
\textbf{铂金瑞生物科技（北京）有限公司}\\
\textit{PLATINUM LIFE EXCELLENCE BIOTECH}\\[0.5em]
{\Large 文件编号：RBC-YZ-SMP-QC-01-01-d \qquad 版本号：1.4}\\
{\LARGE 原 \quad 请 \quad 验 \quad 单}
\end{center}

\noindent
\begin{tabularx}{\textwidth}{|p{0.13\textwidth}|X|p{0.13\textwidth}|X|}
\hline
样品名称 &
% #VALUE_ID: LEX-P0041-V0001
% #FIELD_VALUE: 样品名称
% #TODO #HANDWRITTEN: 二级纯化水？；字迹较潦草
\fieldvalue{\handwritten{二级纯化水？}} &
样品代码 &
% #VALUE_ID: LEX-P0041-V0002
% #FIELD_VALUE: 样品代码
% #TODO #HANDWRITTEN: 62640？；末位及首位字迹不清
\fieldvalue{\handwritten{62640？}}\\
\hline
样品批号 &
% #VALUE_ID: LEX-P0041-V0003
% #FIELD_VALUE: 样品批号
% #TODO #HANDWRITTEN: A372202578300？；部分字符难以辨认
\fieldvalue{\handwritten{A372202578300？}} &
样品规格/数量 &
% #VALUE_ID: LEX-P0041-V0004
% #FIELD_VALUE: 样品规格/数量
% #TODO #HANDWRITTEN: 约20 ml／瓶×1；字迹不清
\fieldvalue{\handwritten{约20 ml／瓶×1}}\\
\hline
样品类别 &
\multicolumn{3}{l|}{
% #VALUE_ID: LEX-P0041-V0005
% #FIELD_VALUE: 样品类别（已选项）
% #HANDWRITTEN: √ 生产过程中检验样品
\fieldvalue{\handwritten{√ 生产过程中检验样品}}
\quad $\square$ 待包装产品
\quad $\square$ 成品}\\
\cline{1-4}
& \multicolumn{3}{l|}{$\square$ 其他（\hspace{4em}）}\\
\hline
检验类型 &
\multicolumn{3}{l|}{
% #VALUE_ID: LEX-P0041-V0006
% #FIELD_VALUE: 检验类型（已选项）
% #HANDWRITTEN: √ 初验
\fieldvalue{\handwritten{√ 初验}}
\quad $\square$ 复验
\quad $\square$ 退货
\quad $\square$ 稳定性
\quad $\square$ 留样}\\
\hline
请验目的 &
\multicolumn{3}{l|}{
% #VALUE_ID: LEX-P0041-V0007
% #FIELD_VALUE: 请验目的（已选项）
% #HANDWRITTEN: √ 全检
\fieldvalue{\handwritten{√ 全检}}
\quad $\square$ 留样
\quad $\square$ 其他（\hspace{4em}）}\\
\hline
请验部门 &
\multicolumn{3}{l|}{$\square$ 生产部
\quad $\square$ 物料管理部
\quad $\square$ 其他（\hspace{4em}）}\\
\hline
样品情况 &
\multicolumn{3}{l|}{
% #VALUE_ID: LEX-P0041-V0008
% #FIELD_VALUE: 样品情况（已选项）
% #HANDWRITTEN: √ 随请验单一起提供
\fieldvalue{\handwritten{√ 随请验单一起提供}}
\quad $\square$ 未随请验单一起提供}\\
\hline
请验人 &
% #VALUE_ID: LEX-P0041-V0009
% #FIELD_VALUE: 请验人
% #TODO #HANDWRITTEN: 邵金波？；签名字迹难以辨认
\fieldvalue{\handwritten{邵金波？}} &
请验日期 &
% #VALUE_ID: LEX-P0041-V0010
% #FIELD_VALUE: 请验日期
% #HANDWRITTEN: 2025-08-27
\fieldvalue{\handwritten{2025-08-27}}\\
\hline
收验人 &
% #VALUE_ID: LEX-P0041-V0011
% #FIELD_VALUE: 收验人
% #TODO #HANDWRITTEN: 张？；签名字迹难以辨认
\fieldvalue{\handwritten{张？}} &
收验日期 &
% #VALUE_ID: LEX-P0041-V0012
% #FIELD_VALUE: 收验日期
% #HANDWRITTEN: 2025-08-27
\fieldvalue{\handwritten{2025-08-27}}\\
\hline
备注 &
\multicolumn{3}{l|}{
% #VALUE_ID: LEX-P0041-V0013
% #FIELD_VALUE: 备注
% #TODO #HANDWRITTEN: 27.4h \quad 2.0？；第二项单位及末位不清
\fieldvalue{\handwritten{27.4h \quad 2.0？}}}\\
\hline
\end{tabularx}

% TODO: 页面上方的圆形公司标志及右上角“受控文件”印章未单独绘制。
% LEXOID_PAGE_COMPLETED: 41/60

\newpage

\begin{center}
{\large \textbf{取样记录}}
\end{center}

\noindent
文件编码：REC-YZ-SMP-QC-01-004-6 \qquad 版本号：1.4

\vspace{0.5em}

\renewcommand{\arraystretch}{1.8}
\begin{tabular}{|c|p{4.0cm}|c|p{3.5cm}|}
\hline
名称 &
% #VALUE_ID: LEX-P0042-V0001
% #FIELD_VALUE: 名称
\fieldvalue{一般生物反应器细胞液} &
代码 &
% #VALUE_ID: LEX-P0042-V0002
% #FIELD_VALUE: 代码
\fieldvalue{QY0401}
\\ \hline

规格 &
% #VALUE_ID: LEX-P0042-V0003
% #FIELD_VALUE: 规格
\fieldvalue{约20ml/袋} &
总数量 &
% #VALUE_ID: LEX-P0042-V0004
% #FIELD_VALUE: 总数量
\fieldvalue{1袋}
\\ \hline

批号 &
% #VALUE_ID: LEX-P0042-V0005
% #FIELD_VALUE: 批号
\fieldvalue{A31Z2202508008} &
进厂批号 &
% #VALUE_ID: LEX-P0042-V0006
% #FIELD_VALUE: 进厂批号
\fieldvalue{/}
\\ \hline

供货单位/厂家 &
\multicolumn{3}{l|}{
% #VALUE_ID: LEX-P0042-V0007
% #FIELD_VALUE: 供货单位/厂家
\fieldvalue{/}}
\\ \hline

取样目的 &
\multicolumn{3}{l|}{
% #VALUE_ID: LEX-P0042-V0008
% #FIELD_VALUE: 取样目的
\fieldvalue{\checkmark\,正常检验}
\quad $\square$\,复检
\quad $\square$\,稳定性}
\\[-1.2em]
& \multicolumn{3}{r|}{
$\square$\,退回待检
\quad $\square$\,留样
\quad $\square$\,其他：\underline{\hspace{1.5cm}}}
\\ \hline

取样地点 &
\multicolumn{3}{l|}{
% #VALUE_ID: LEX-P0042-V0009
% #FIELD_VALUE: 取样地点
\fieldvalue{\checkmark\,产品出口}
\quad $\square$\,原料库
\quad $\square$\,细胞库
\quad $\square$\,包装间}
\\ \hline

贮存条件 &
\multicolumn{3}{l|}{
% #VALUE_ID: LEX-P0042-V0010
% #FIELD_VALUE: 贮存条件
\fieldvalue{\checkmark\,2--8℃}
\quad $\square$\,{-20℃}
\quad $\square$\,{-80℃}
\quad $\square$\,常温
\quad $\square$\,其他：\underline{\hspace{1.2cm}}}
\\ \hline

盛装容器 &
\multicolumn{3}{l|}{
% #VALUE_ID: LEX-P0042-V0011
% #FIELD_VALUE: 盛装容器
\fieldvalue{\checkmark\,保温箱}
\quad $\square$\,泡沫箱
\quad $\square$\,手提篮
\quad $\square$\,其他：\underline{\hspace{1.2cm}}}
\\ \hline

外观复核 &
\multicolumn{3}{l|}{
% #VALUE_ID: LEX-P0042-V0012
% #FIELD_VALUE: 外观复核
\fieldvalue{\checkmark\,合格}
\quad $\square$\,不合格}
\\ \hline

取样量 &
% #VALUE_ID: LEX-P0042-V0013
% #FIELD_VALUE: 取样量
\fieldvalue{1袋} &
检验量 &
% #VALUE_ID: LEX-P0042-V0014
% #FIELD_VALUE: 检验量
\fieldvalue{1袋}
\\ \hline

留样量 &
% #VALUE_ID: LEX-P0042-V0015
% #FIELD_VALUE: 留样量
\fieldvalue{/} &
\multicolumn{2}{c|}{}
\\ \hline

取样人/日期 &
% #VALUE_ID: LEX-P0042-V0016
% #FIELD_VALUE: 取样人
% #TODO #HANDWRITTEN: 徐峰; handwritten signature is partially unclear
\fieldvalue{\handwritten{徐峰}}
\quad
% #VALUE_ID: LEX-P0042-V0017
% #FIELD_VALUE: 取样日期
% #TODO #HANDWRITTEN: 2024.08.7; handwritten date is partially unclear
\fieldvalue{\handwritten{2024.08.7}} &
复核人/日期 &
% #VALUE_ID: LEX-P0042-V0018
% #FIELD_VALUE: 复核人
% #TODO #HANDWRITTEN: 肖身立; handwritten signature is partially unclear
\fieldvalue{\handwritten{肖身立}}
\quad
% #VALUE_ID: LEX-P0042-V0019
% #FIELD_VALUE: 复核日期
% #TODO #HANDWRITTEN: 2024.08.7; handwritten date is partially unclear
\fieldvalue{\handwritten{2024.08.7}}
\\ \hline

备注 &
\multicolumn{3}{l|}{
% #VALUE_ID: LEX-P0042-V0020
% #FIELD_VALUE: 备注
% #TODO #HANDWRITTEN: -7.4h; handwritten note is partially unclear
\fieldvalue{\handwritten{-7.4h}}
\qquad
% #VALUE_ID: LEX-P0042-V0021
% #FIELD_VALUE: 备注
% #TODO #HANDWRITTEN: 2.0L; handwritten note is partially unclear
\fieldvalue{\handwritten{2.0L}}}
\\ \hline
\end{tabular}

\vfill
\begin{center}
第 1 页 共 1 页
\end{center}
% LEXOID_PAGE_COMPLETED: 42/60

\newpage

\section*{生物反应器细胞液检验操作记录}

\noindent\textbf{葡萄糖浓度检验}

\begin{tabularx}{\textwidth}{|p{2.4cm}|X|p{2.4cm}|X|}
\hline
样品名称 &
% #VALUE_ID: LEX-P0043-V0001
% #FIELD_VALUE: 样品名称
\fieldvalue{一级生物反应器细胞液} &
样品规格 &
% #VALUE_ID: LEX-P0043-V0002
% #FIELD_VALUE: 样品规格
% #HANDWRITTEN: 4×20
\fieldvalue{\handwritten{4×20}} ml/袋
\\ \hline
样品批号 &
% #VALUE_ID: LEX-P0043-V0003
% #FIELD_VALUE: 样品批号
% #TODO #HANDWRITTEN: 20140527; handwriting partially unclear
\fieldvalue{\handwritten{20140527}} &
取样日期 &
% #VALUE_ID: LEX-P0043-V0004
% #FIELD_VALUE: 取样日期
% #HANDWRITTEN: 2014.05.27
\fieldvalue{\handwritten{2014.05.27}} 日
\\ \hline
\end{tabularx}

\subsection*{1\quad 仪器设备、试剂耗材：}

\subsubsection*{1.1\quad 仪器设备}

\begin{tabularx}{\textwidth}{|p{2.1cm}|X|p{2.1cm}|p{2.4cm}|p{2.7cm}|}
\hline
名称 & 厂家 & 型号 & 设备编号 & 检查确认
\\ \hline
血糖仪 & 三诺生物传感股份有限公司 & 安稳+ &
% #VALUE_ID: LEX-P0043-V0005
% #FIELD_VALUE: 设备编号
% #HANDWRITTEN: /
\fieldvalue{\handwritten{/}} &
% #VALUE_ID: LEX-P0043-V0006
% #FIELD_VALUE: 检查确认（正常）
% #HANDWRITTEN: 勾选
\fieldvalue{\handwritten{\checkmark}} 正常\quad $\square$ 异常
\\ \hline
移液器 & Eppendorf Corporate & 0.1--2.5\,$\mu$l &
% #VALUE_ID: LEX-P0043-V0007
% #FIELD_VALUE: 设备编号
% #TODO #HANDWRITTEN: 不清晰; handwritten equipment number is illegible
\fieldvalue{\handwritten{不清晰}} &
% #VALUE_ID: LEX-P0043-V0008
% #FIELD_VALUE: 检查确认（正常）
% #HANDWRITTEN: 勾选
\fieldvalue{\handwritten{\checkmark}} 正常\quad $\square$ 异常
\\ \hline
\end{tabularx}

\subsubsection*{1.2\quad 试剂耗材}

\begin{tabularx}{\textwidth}{|p{2.1cm}|X|p{2.4cm}|p{2.4cm}|}
\hline
名称 & 厂家 & 批号 & 有效期至
\\ \hline
血糖测试条 & 三诺生物传感股份有限公司 &
% #VALUE_ID: LEX-P0043-V0009
% #FIELD_VALUE: 批号
% #TODO #HANDWRITTEN: 4?4?07; handwriting unclear
\fieldvalue{\handwritten{4?4?07}} &
% #VALUE_ID: LEX-P0043-V0010
% #FIELD_VALUE: 有效期至
% #HANDWRITTEN: 2021.10.27
\fieldvalue{\handwritten{2021.10.27}}
\\ \hline
\end{tabularx}

\section*{2\quad 操作步骤：}

\begin{tabularx}{\textwidth}{|p{0.28\textwidth}|X|X|X|}
\hline
\multicolumn{4}{|c|}{操作过程}
\\ \hline
\begin{minipage}[t]{0.27\textwidth}
取出测试条，取样品自然沉降后的上清液0.6\,\textmu l加入测试条电极反应室，血糖仪自动发出“嘀”提示音，吸入样品完成，血糖仪自动进入倒计时，开始测试；
\end{minipage}
&
\multicolumn{3}{l|}{
取上清液：
% #VALUE_ID: LEX-P0043-V0011
% #FIELD_VALUE: 取上清液体积
% #HANDWRITTEN: 0.6
\fieldvalue{\handwritten{0.6}}\,\textmu l}
\\ \hline
重复以上步骤2次，共计数3次。 &
\multicolumn{3}{l|}{
每个样品分别共计数
% #VALUE_ID: LEX-P0043-V0012
% #FIELD_VALUE: 每个样品计数次数
% #HANDWRITTEN: 3
\fieldvalue{\handwritten{3}} 次。}
\\ \hline
\multicolumn{4}{|c|}{检测结果}
\\ \hline
\multicolumn{4}{|c|}{葡萄糖浓度（mmol/L）}
\\ \hline
第一次 & 第二次 & 第三次 & 平均值（保留一位小数）
\\ \hline
% #VALUE_ID: LEX-P0043-V0013
% #FIELD_VALUE: 第一次检测结果
% #HANDWRITTEN: 21.9
\fieldvalue{\handwritten{21.9}} &
% #VALUE_ID: LEX-P0043-V0014
% #FIELD_VALUE: 第二次检测结果
% #HANDWRITTEN: 21.1
\fieldvalue{\handwritten{21.1}} &
% #VALUE_ID: LEX-P0043-V0015
% #FIELD_VALUE: 第三次检测结果
% #HANDWRITTEN: 21.3
\fieldvalue{\handwritten{21.3}} &
% #VALUE_ID: LEX-P0043-V0016
% #FIELD_VALUE: 平均值
% #HANDWRITTEN: 21.1
\fieldvalue{\handwritten{21.1}}
\\ \hline
备注 &
% #VALUE_ID: LEX-P0043-V0017
% #FIELD_VALUE: 备注
% #HANDWRITTEN: /
\fieldvalue{\handwritten{/}} &
\multicolumn{2}{l|}{}
\\ \hline
实验人/日期 &
% #VALUE_ID: LEX-P0043-V0018
% #FIELD_VALUE: 实验人/日期
% #TODO #HANDWRITTEN: 签名，2014.5.27; signature is unclear
\fieldvalue{\handwritten{签名，2014.5.27}} &
\multicolumn{2}{l|}{}
\\ \hline
复核人/日期 & \multicolumn{3}{l|}{}
\\ \hline
\end{tabularx}

% TODO: The source page contains handwritten marks whose exact character forms are partially illegible.
% LEXOID_PAGE_COMPLETED: 43/60

\newpage

\begin{center}
生物反应器细胞液检验记录
\end{center}

\begin{tabularx}{\textwidth}{|c|X|c|X|}
\hline
\multicolumn{4}{|c|}{密度与活率检验}\\
\hline
样品名称 &
% #VALUE_ID: LEX-P0044-V0001
% #FIELD_VALUE: 样品名称
\fieldvalue{一级生物反应器细胞液} &
样品规格 &
% #VALUE_ID: LEX-P0044-V0002
% #FIELD_VALUE: 样品规格
% #HANDWRITTEN: 6250
\fieldvalue{\handwritten{6250}} ml/袋\\
\hline
样品批号 &
% #VALUE_ID: LEX-P0044-V0003
% #FIELD_VALUE: 样品批号
% #TODO #HANDWRITTEN: 例？；字迹不清
\fieldvalue{\handwritten{例？}} &
取样日期 &
% #VALUE_ID: LEX-P0044-V0004
% #FIELD_VALUE: 取样日期
% #TODO #HANDWRITTEN: 24年？月？日；字迹不清
\fieldvalue{\handwritten{24年？月？日}}\\
\hline
\end{tabularx}

\section{仪器设备、试剂耗材}

\subsection{仪器设备}

\begin{tabularx}{\textwidth}{|c|c|c|c|c|c|}
\hline
名称 & 厂家 & 型号 & 设备编号 & 校准有效期至 & 检查确认\\
\hline
细胞计数仪 &
% #VALUE_ID: LEX-P0044-V0005
% #FIELD_VALUE: 厂家
\fieldvalue{Nexcelom} &
% #VALUE_ID: LEX-P0044-V0006
% #FIELD_VALUE: 型号
\fieldvalue{Cellometer K2} &
% #VALUE_ID: LEX-P0044-V0007
% #FIELD_VALUE: 设备编号
\fieldvalue{1410905} &
% #VALUE_ID: LEX-P0044-V0008
% #FIELD_VALUE: 校准有效期至
% #HANDWRITTEN: 2025.06.15
\fieldvalue{\handwritten{2025.06.15}} &
% #VALUE_ID: LEX-P0044-V0009
% #FIELD_VALUE: 检查确认
\fieldvalue{\checkmark 正常\quad 异常}\\
\hline
高速冷冻离心机 &
% #VALUE_ID: LEX-P0044-V0010
% #FIELD_VALUE: 厂家
\fieldvalue{Thermo} &
% #VALUE_ID: LEX-P0044-V0011
% #FIELD_VALUE: 型号
\fieldvalue{ST16R} &
% #VALUE_ID: LEX-P0044-V0012
% #FIELD_VALUE: 设备编号
\fieldvalue{1410810} &
% #VALUE_ID: LEX-P0044-V0013
% #FIELD_VALUE: 校准有效期至
% #HANDWRITTEN: 2025.07.09
\fieldvalue{\handwritten{2025.07.09}} &
% #VALUE_ID: LEX-P0044-V0014
% #FIELD_VALUE: 检查确认
\fieldvalue{\checkmark 正常\quad 异常}\\
\hline
电热恒温水槽 &
% #VALUE_ID: LEX-P0044-V0015
% #FIELD_VALUE: 厂家
\fieldvalue{上海一恒} &
% #VALUE_ID: LEX-P0044-V0016
% #FIELD_VALUE: 型号
\fieldvalue{DK-8AXX} &
% #VALUE_ID: LEX-P0044-V0017
% #FIELD_VALUE: 设备编号
\fieldvalue{1411001} &
% #VALUE_ID: LEX-P0044-V0018
% #FIELD_VALUE: 校准有效期至
% #HANDWRITTEN: 2025.07.12
\fieldvalue{\handwritten{2025.07.12}} &
% #VALUE_ID: LEX-P0044-V0019
% #FIELD_VALUE: 检查确认
\fieldvalue{\checkmark 正常\quad 异常}\\
\hline
\end{tabularx}

\subsection{试剂耗材}

\begin{tabularx}{\textwidth}{|c|c|c|c|}
\hline
名称 & 厂家 & 批号 & 有效期至\\
\hline
AO/PI 染液 &
% #VALUE_ID: LEX-P0044-V0020
% #FIELD_VALUE: 厂家
\fieldvalue{Nexcelom} &
% #VALUE_ID: LEX-P0044-V0021
% #FIELD_VALUE: 批号
% #TODO #HANDWRITTEN: 3？47？；字迹不清
\fieldvalue{\handwritten{3？47？}} &
% #VALUE_ID: LEX-P0044-V0022
% #FIELD_VALUE: 有效期至
% #HANDWRITTEN: 2026.01.12
\fieldvalue{\handwritten{2026.01.12}}\\
\hline
0.9\%氯化钠注射液 &
% #VALUE_ID: LEX-P0044-V0023
% #FIELD_VALUE: 厂家
\fieldvalue{石家庄四药} &
% #VALUE_ID: LEX-P0044-V0024
% #FIELD_VALUE: 批号
% #TODO #HANDWRITTEN: 240？130？；字迹不清
\fieldvalue{\handwritten{240？130？}} &
% #VALUE_ID: LEX-P0044-V0025
% #FIELD_VALUE: 有效期至
% #HANDWRITTEN: 2024.12.06
\fieldvalue{\handwritten{2024.12.06}}\\
\hline
20\%人血白蛋白 &
% #VALUE_ID: LEX-P0044-V0026
% #FIELD_VALUE: 厂家
\fieldvalue{Baxter AG} &
% #VALUE_ID: LEX-P0044-V0027
% #FIELD_VALUE: 批号
% #TODO #HANDWRITTEN: 841？；字迹不清
\fieldvalue{\handwritten{841？}} &
% #VALUE_ID: LEX-P0044-V0028
% #FIELD_VALUE: 有效期至
% #HANDWRITTEN: 2027.12.28
\fieldvalue{\handwritten{2027.12.28}}\\
\hline
\end{tabularx}

\section{操作步骤}

\begin{tabularx}{\textwidth}{|X|X|}
\hline
\multicolumn{2}{|c|}{溶液准备}\\
\hline
1\%人白蛋白溶液配制：取 7.5 ml 20\%人血白蛋白加入到 95 ml 的 0.9\%氯化钠注射液中，混匀即得。 &
20\%人血白蛋白
% #VALUE_ID: LEX-P0044-V0029
% #FIELD_VALUE: 20\%人血白蛋白用量
% #HANDWRITTEN: 1
\fieldvalue{\handwritten{1}} ml\\
&
0.9\%氯化钠注射液
% #VALUE_ID: LEX-P0044-V0030
% #FIELD_VALUE: 0.9\%氯化钠注射液用量
% #HANDWRITTEN: 19
\fieldvalue{\handwritten{19}} ml\\
\hline
\end{tabularx}

\begin{tabularx}{\textwidth}{|c|X|X|}
\hline
名称 & 配制批号 & 有效期至\\
\hline
5$\times$裂解液 &
% #VALUE_ID: LEX-P0044-V0031
% #FIELD_VALUE: 配制批号
% #TODO #HANDWRITTEN: 5？；字迹不清
\fieldvalue{\handwritten{5？}} &
% #VALUE_ID: LEX-P0044-V0032
% #FIELD_VALUE: 有效期至
% #HANDWRITTEN: 2025.10.07
\fieldvalue{\handwritten{2025.10.07}}\\
\hline
\end{tabularx}

\begin{tabularx}{\textwidth}{|X|X|}
\hline
\multicolumn{2}{|c|}{操作过程}\\
\hline
按照 50 ml 离心管的刻度读取样品体积，往样品中加入对应体积的 5$\times$裂解液（每 4 ml 样品加入 1 ml 5$\times$裂解液，按照“只进不舍”原则，如 26 ml 样品，加入 7 ml 5$\times$裂解液），充分混匀后倒入至 2、3 的样品中，置于 37℃电热恒温水槽中裂解 20 min 后，取出离心管摇匀，使样品充分裂解；再置于 37℃电热恒温水槽中裂解 10 min 得到细胞悬液。 &
样品体积：
% #VALUE_ID: LEX-P0044-V0033
% #FIELD_VALUE: 样品体积
% #HANDWRITTEN: 24
\fieldvalue{\handwritten{24}} ml\\
裂解液：
% #VALUE_ID: LEX-P0044-V0034
% #FIELD_VALUE: 裂解液
% #HANDWRITTEN: 6
\fieldvalue{\handwritten{6}} ml\\
电热恒温水槽：
% #VALUE_ID: LEX-P0044-V0035
% #FIELD_VALUE: 电热恒温水槽温度
% #HANDWRITTEN: 37
\fieldvalue{\handwritten{37}}℃\\
裂解：
% #VALUE_ID: LEX-P0044-V0036
% #FIELD_VALUE: 裂解时间
% #TODO #HANDWRITTEN: 0？12？--0？30？；字迹不清
\fieldvalue{\handwritten{0？12？--0？30？}}\\
\hline
取出裂解好的细胞悬液 300$\times$g 离心 5 分钟，按照读取样品体积加入对应体积的 1\%人白蛋白溶液重量（每 1.2 ml 样品体积加入 0.1 ml 1\%人白蛋白溶液，按照“只进不舍”原则，如 26 ml 样品，吸取 2.6 ml 1\%人白蛋白溶液）。 &
% #VALUE_ID: LEX-P0044-V0037
% #FIELD_VALUE: 离心条件
% #HANDWRITTEN: 300×g 离心 5分钟
\fieldvalue{\handwritten{300$\times$g 离心 5分钟}}\\
加入 1\%人白蛋白溶液覆盖体积：
% #VALUE_ID: LEX-P0044-V0038
% #FIELD_VALUE: 1\%人白蛋白溶液覆盖体积
% #HANDWRITTEN: 2
\fieldvalue{\handwritten{2}} ml\\
吸取：
% #VALUE_ID: LEX-P0044-V0039
% #FIELD_VALUE: 吸取体积
% #HANDWRITTEN: 20
\fieldvalue{\handwritten{20}} $\mu$l\\
\hline
\end{tabularx}

\begin{center}
第 2 页 共 5 页
\end{center}
% LEXOID_PAGE_COMPLETED: 44/60

\newpage

\begin{center}
\textbf{实验生物}\\
% #VALUE_ID: LEX-P0045-V0001
% #FIELD_VALUE: 实验生物类型
\fieldvalue{原代细胞}
\hspace{3em}
% #VALUE_ID: LEX-P0045-V0002
% #FIELD_VALUE: 稀释度
\fieldvalue{4}
\end{center}

\noindent
文件编码：REC-YZ-SOP-09-006-a

\begin{minipage}{0.58\textwidth}
加入2.2 ml 1\%人血氰化钠溶液，反复吹打至单细胞悬液。

撕去计数板上下两层保护膜，将细胞悬液混匀，从中吸取20~$\mu$l与20~$\mu$l AO/PI染料混合均匀。

吸取20~$\mu$l细胞悬液与AO/PI染料的混合液，注入到计数板的一个检测室中。

将加入混合液的一端插入进样口。

选择Assay类型，使用AO/PI染料进行双荧光存活率检测；

在主界面左侧ID栏中输入样本批号。

点击主界面左下角“Preview”预览细胞图片。

调整仪器焦距直至观察到荧光通道下的细胞中心透亮，轮廓清晰。

点击“count”按钮进行计数并自动拍摄图片。

重复以上步骤2次，共计数3次，打印数据结果。
\end{minipage}
\hfill
\begin{minipage}{0.36\textwidth}
与
% #VALUE_ID: LEX-P0045-V0003
% #FIELD_VALUE: AO/PI染料用量
% #HANDWRITTEN: 2 $\mu$l
\fieldvalue{\handwritten{2~$\mu$l}}
AO/PI染料混合均匀。

吸取
% #VALUE_ID: LEX-P0045-V0004
% #FIELD_VALUE: 计数吸取量
% #HANDWRITTEN: 20 $\mu$l
\fieldvalue{\handwritten{20~$\mu$l}}
计数。

每个样品共计数
% #VALUE_ID: LEX-P0045-V0005
% #FIELD_VALUE: 计数次数
% #HANDWRITTEN: 3
\fieldvalue{\handwritten{3}}
次。
\end{minipage}

\noindent
\begin{tabularx}{\textwidth}{|l|X|}
\hline
数据保存路径 &
% #VALUE_ID: LEX-P0045-V0006
% #FIELD_VALUE: 数据保存路径
% #HANDWRITTEN: 14/19细胞培养箱\textbackslash 2024/12/14-不良细胞\textbackslash -2代
\fieldvalue{\handwritten{14/19细胞培养箱\textbackslash 2024/12/14-不良细胞\textbackslash -2代}}
\\
\hline
\multicolumn{2}{|c|}{检测结果}\\
\hline
检测结果 &
$\square$
% #VALUE_ID: LEX-P0045-V0007
% #FIELD_VALUE: 检测结果类型
\fieldvalue{一级反应细胞}
、二级反应细胞、三级反应细胞
\\
\hline
\end{tabularx}

\begin{center}
\begin{tabularx}{\textwidth}{|c|X|X|X|X|X|X|}
\hline
检测次数 &
\multicolumn{1}{c|}{活细胞密度}\\
($\times10^6$/ml) &
\multicolumn{1}{c|}{样品密度}\\
($\times10^6$/ml) &
\multicolumn{1}{c|}{细胞密度}\\
CV值（\%） &
\multicolumn{1}{c|}{细胞活率}\\
（\%） &
\multicolumn{1}{c|}{样品活率}\\
（\%） &
\multicolumn{1}{c|}{细胞活率}\\
CV值（\%）\\
\hline
第一次 &
% #VALUE_ID: LEX-P0045-V0008
% #FIELD_VALUE: 第一次活细胞密度
% #TODO #HANDWRITTEN: 7.?? $\times10^6$; 数字笔迹不清
\fieldvalue{\handwritten{7.??$\times10^6$}} &
&
&
% #VALUE_ID: LEX-P0045-V0013
% #FIELD_VALUE: 第一次细胞活率
% #HANDWRITTEN: 84
\fieldvalue{\handwritten{84}} &
&
\\
\hline
第二次 &
% #VALUE_ID: LEX-P0045-V0009
% #FIELD_VALUE: 第二次活细胞密度
% #TODO #HANDWRITTEN: 4.?? $\times10^6$; 数字笔迹不清
\fieldvalue{\handwritten{4.??$\times10^6$}} &
% #VALUE_ID: LEX-P0045-V0011
% #FIELD_VALUE: 样品密度
% #HANDWRITTEN: 4.79 $\times10^6$
\fieldvalue{\handwritten{4.79$\times10^6$}} &
% #VALUE_ID: LEX-P0042-V0012
% #FIELD_VALUE: 细胞密度CV值
% #HANDWRITTEN: 8
\fieldvalue{\handwritten{8}} &
% #VALUE_ID: LEX-P0045-V0014
% #FIELD_VALUE: 第二次细胞活率
% #HANDWRITTEN: 90.6
\fieldvalue{\handwritten{90.6}} &
% #VALUE_ID: LEX-P0045-V0016
% #FIELD_VALUE: 样品活率
% #HANDWRITTEN: 91
\fieldvalue{\handwritten{91}} &
% #VALUE_ID: LEX-P0045-V0017
% #FIELD_VALUE: 细胞活率CV值
% #HANDWRITTEN: 4
\fieldvalue{\handwritten{4}}\\
\hline
第三次 &
% #VALUE_ID: LEX-P0045-V0010
% #FIELD_VALUE: 第三次活细胞密度
% #HANDWRITTEN: 4.21 $\times10^6$
\fieldvalue{\handwritten{4.21$\times10^6$}} &
&
&
% #VALUE_ID: LEX-P0045-V0015
% #FIELD_VALUE: 第三次细胞活率
% #HANDWRITTEN: 88.6
\fieldvalue{\handwritten{88.6}} &
&
\\
\hline
\end{tabularx}
\end{center}

\noindent
\begin{tabularx}{\textwidth}{|l|X|l|X|}
\hline
发解缩细胞体积（ml） &
% #VALUE_ID: LEX-P0045-V0018
% #FIELD_VALUE: 解冻细胞体积
% #TODO #HANDWRITTEN: 20.??; 数字笔迹不清
\fieldvalue{\handwritten{20.??}} &
细胞数量（个细胞/瓶） &
% #VALUE_ID: LEX-P0045-V0019
% #FIELD_VALUE: 细胞数量
% #HANDWRITTEN: 8.18 $\times10^5$
\fieldvalue{\handwritten{8.18$\times10^5$}}
\\
\hline
\multicolumn{4}{|l|}{注：细胞密度为3次活细胞密度结果取平均值乘以仪器乘以1\%人血氰化钠液稀释体积再除以生物反应器}\\
\multicolumn{4}{|l|}{细胞总体积；细胞数量为样品密度乘以发解缩体积。}\\
\hline
备注 &
% #VALUE_ID: LEX-P0045-V0020
% #FIELD_VALUE: 备注
% #HANDWRITTEN: /
\fieldvalue{\handwritten{/}} &
\multicolumn{2}{l|}{}\\
\hline
实验人/日期 &
% #VALUE_ID: LEX-P0045-V0021
% #FIELD_VALUE: 实验人
% #TODO #HANDWRITTEN: [姓名]; 签名难以辨认
\fieldvalue{\handwritten{[姓名]}}
\quad
% #VALUE_ID: LEX-P0045-V0022
% #FIELD_VALUE: 实验日期
% #HANDWRITTEN: 2024.3.7
\fieldvalue{\handwritten{2024.3.7}} &
\multicolumn{2}{l|}{}\\
\hline
复核人/日期 &
% #VALUE_ID: LEX-P0045-V0023
% #FIELD_VALUE: 复核人
% #TODO #HANDWRITTEN: [姓名]; 签名难以辨认
\fieldvalue{\handwritten{[姓名]}}
\quad
% #VALUE_ID: LEX-P0045-V0024
% #FIELD_VALUE: 复核日期
% #TODO #HANDWRITTEN: 2025.8.2; 日期笔迹不清
\fieldvalue{\handwritten{2025.8.2}} &
\multicolumn{2}{l|}{}\\
\hline
\end{tabularx}

\begin{center}
第3页 共5页
\end{center}
% LEXOID_PAGE_COMPLETED: 45/60

\newpage

\begin{center}
\textbf{生物反应器细胞液检测操作记录}
\end{center}

\begin{tabularx}{\textwidth}{|>{\centering\arraybackslash}p{2.2cm}|X|>{\centering\arraybackslash}p{2.2cm}|X|}
\hline
\multicolumn{4}{|c|}{细胞生长状态检验}\\
\hline
样品名称 &
% #VALUE_ID: LEX-P0046-V0001
% #FIELD_VALUE: 样品名称
\fieldvalue{生物反应器细胞液} &
样品规格 &
% #VALUE_ID: LEX-P0046-V0002
% #FIELD_VALUE: 样品规格
% #TODO #HANDWRITTEN: 1? ml/袋; handwriting is unclear
\fieldvalue{\handwritten{1? ml/袋}}\\
\hline
样品批号 &
% #VALUE_ID: LEX-P0046-V0003
% #FIELD_VALUE: 样品批号
% #TODO #HANDWRITTEN: [难以辨认]; handwritten batch number is unclear
\fieldvalue{\handwritten{[难以辨认]}} &
取样日期 &
% #VALUE_ID: LEX-P0046-V0004
% #FIELD_VALUE: 取样日期
% #HANDWRITTEN: 2024年06月27日
\fieldvalue{\handwritten{2024年06月27日}}\\
\hline
\end{tabularx}

\subsection*{1.\quad 仪器设备及试剂}

\subsubsection*{1.1.\quad 仪器设备：}

\begin{tabularx}{\textwidth}{|>{\centering\arraybackslash}X|>{\centering\arraybackslash}X|>{\centering\arraybackslash}X|>{\centering\arraybackslash}X|>{\centering\arraybackslash}X|}
\hline
仪器名称 & 厂家 & 型号 & 设备编号 & 检查确认\\
\hline
% #VALUE_ID: LEX-P0046-V0005
% #FIELD_VALUE: 仪器名称
\fieldvalue{荧光显微镜} &
% #VALUE_ID: LEX-P0046-V0006
% #FIELD_VALUE: 厂家
\fieldvalue{朱伟} &
% #VALUE_ID: LEX-P0046-V0007
% #FIELD_VALUE: 型号
\fieldvalue{Dmi1 LED} &
% #VALUE_ID: LEX-P0046-V0008
% #FIELD_VALUE: 设备编号
\fieldvalue{1431404} &
% #VALUE_ID: LEX-P0046-V0009
% #FIELD_VALUE: 检查确认
% #HANDWRITTEN: √
\fieldvalue{\handwritten{\checkmark}\,正常\quad□异常}\\
\hline
\end{tabularx}

\subsubsection*{1.2.\quad 试剂耗材}

\begin{tabularx}{\textwidth}{|>{\centering\arraybackslash}X|>{\centering\arraybackslash}X|>{\centering\arraybackslash}X|>{\centering\arraybackslash}X|>{\centering\arraybackslash}X|}
\hline
名称 & 厂家 & 货号 & 批号 & 有效期至\\
\hline
% #VALUE_ID: LEX-P0046-V0010
% #FIELD_VALUE: 名称
\fieldvalue{细胞活力检测试剂盒} &
% #VALUE_ID: LEX-P0046-V0011
% #FIELD_VALUE: 厂家
\fieldvalue{江苏凯基生物技术股份有限公司} &
% #VALUE_ID: LEX-P0046-V0012
% #FIELD_VALUE: 货号
\fieldvalue{KGA9501-1000} &
% #VALUE_ID: LEX-P0046-V0013
% #FIELD_VALUE: 批号
% #TODO #HANDWRITTEN: 2?; handwritten batch number is unclear
\fieldvalue{\handwritten{2?}} &
% #VALUE_ID: LEX-P0046-V0014
% #FIELD_VALUE: 有效期至
% #TODO #HANDWRITTEN: 2024.09?; handwritten expiry date is partly unclear
\fieldvalue{\handwritten{2024.09?}}\\
\hline
\end{tabularx}

\subsection*{2.\quad 操作步骤：}

\begin{tabularx}{\textwidth}{|>{\centering\arraybackslash}X|}
\hline
\textbf{溶液配制}\\
\hline
20$\times$的荧光染液配制：取
% #VALUE_ID: LEX-P0046-V0015
% #FIELD_VALUE: 组分A用量
% #HANDWRITTEN: 0.2
\fieldvalue{\handwritten{0.2}}\,\textmu l 组分 A 和
% #VALUE_ID: LEX-P0046-V0016
% #FIELD_VALUE: 组分B用量
% #HANDWRITTEN: 0.2
\fieldvalue{\handwritten{0.2}}\,\textmu l 组分 B 与
% #VALUE_ID: LEX-P0046-V0017
% #FIELD_VALUE: PBS用量
% #HANDWRITTEN: 39.6
\fieldvalue{\handwritten{39.6}}\,\textmu l PBS 混合均匀配成 20$\times$ 的荧光染液。\\
\hline
\textbf{操作过程}\\
\hline
\begin{minipage}{0.96\linewidth}
取混匀后样品 200\textmu l/孔至 96 孔板中，取 2 孔；
\hfill 样品
% #VALUE_ID: LEX-P0046-V0018
% #FIELD_VALUE: 样品用量
% #HANDWRITTEN: 200
\fieldvalue{\handwritten{200}}\,\textmu l/孔

每个样品
% #VALUE_ID: LEX-P0046-V0019
% #FIELD_VALUE: 每个样品孔数
% #TODO #HANDWRITTEN: 2孔; handwriting is unclear
\fieldvalue{\handwritten{2孔}}

自然沉降后，吸弃上清，注意不要吸到沉淀样品；
\hfill $\square$ 吸弃上清

取 20$\times$ 的荧光染液加入到 96 孔板中，200\textmu l/孔，避光孵育 30--60min；
\hfill 20$\times$ 的荧光染液
% #VALUE_ID: LEX-P0046-V0020
% #FIELD_VALUE: 荧光染液用量
% #HANDWRITTEN: 200
\fieldvalue{\handwritten{200}}\,\textmu l/孔

避光孵育
% #VALUE_ID: LEX-P0046-V0021
% #FIELD_VALUE: 避光孵育时间
% #TODO #HANDWRITTEN: 9:??--11:30; handwritten time range is partly illegible
\fieldvalue{\handwritten{9:??--11:30}}

染色结束后，吸弃上清，取 PBS 加入到以上 96 孔板中，200\textmu l/孔，吸弃上清，注意不要吸到沉淀样品，清洗 3 次，之后每孔加入 100\textmu l PBS 进行拍照；
\hfill 加入 PBS
% #VALUE_ID: LEX-P0046-V0022
% #FIELD_VALUE: PBS用量
% #HANDWRITTEN: 200
\fieldvalue{\handwritten{200}}\,\textmu l/孔

清洗
% #VALUE_ID: LEX-P0046-V0023
% #FIELD_VALUE: 清洗次数
% #HANDWRITTEN: 3
\fieldvalue{\handwritten{3}} 次

将荧光显微镜调至 4$\times$镜，软件中调节至相应倍数参数，调节到适合的光源和像距，拍摄 4 张清晰图像。

选一张清晰图像，打印并粘贴在实验记录上。
\end{minipage}\\
\hline
数据保存路径 \hfill
% #VALUE_ID: LEX-P0046-V0024
% #FIELD_VALUE: 数据保存路径
% #TODO #HANDWRITTEN: 147?10?细胞...; handwritten storage path is largely unclear
\fieldvalue{\handwritten{147?10?细胞\ldots}}\\
\hline
\textbf{检测结果}\\
\hline
\end{tabularx}

\begin{center}
第 4 页 共 5 页
\end{center}
% LEXOID_PAGE_COMPLETED: 46/60

\newpage

\noindent
\begin{minipage}{\textwidth}
\small
铂立生物科技（北京）有限公司 \hfill 文件编码：REC-YZ-SOP-QC-99-006-a\\
\hfill \textcolor{red}{受控文件}\\
\hfill \textcolor{red}{版本号：1.4}
\end{minipage}

\vspace{0.4cm}

\noindent
\framebox{\parbox[c][17cm][c]{\dimexpr\textwidth-2\fboxsep-2\fboxrule\relax}{%
\centering
% TODO: Insert the photograph visible on this page.
\fbox{\rule{0pt}{3.2cm}\rule{6.5cm}{0pt}}%

\vspace{0.1cm}

% #VALUE_ID: LEX-P0047-V0001
% #HANDWRITTEN: 梅... 2025/08/27
\handwritten{梅... 2025/08/27}
}}

\vspace{0.1cm}

\noindent
\begin{tabularx}{\textwidth}{|p{2.2cm}|X|}
\hline
备注 &
% #VALUE_ID: LEX-P0047-V0002
% #FIELD_VALUE: 备注
% #HANDWRITTEN: /
\fieldvalue{\handwritten{/}}\\
\hline
实验人/日期 &
% #VALUE_ID: LEX-P0047-V0003
% #FIELD_VALUE: 实验人
% #TODO #HANDWRITTEN: 梅...; handwriting unclear
\fieldvalue{\handwritten{梅...}}\quad
% #VALUE_ID: LEX-P0047-V0004
% #FIELD_VALUE: 日期
% #HANDWRITTEN: 2025.08.27
\fieldvalue{\handwritten{2025.08.27}}\\
\hline
复核人/日期 &
% #VALUE_ID: LEX-P0047-V0005
% #FIELD_VALUE: 复核人
% #TODO #HANDWRITTEN: 邱...; handwriting unclear
\fieldvalue{\handwritten{邱...}}\quad
% #VALUE_ID: LEX-P0047-V0006
% #FIELD_VALUE: 日期
% #HANDWRITTEN: 2025.08.27
\fieldvalue{\handwritten{2025.08.27}}\\
\hline
\end{tabularx}

\vfill

\begin{center}
第 5 页 共 5 页
\end{center}
% LEXOID_PAGE_COMPLETED: 47/60

\newpage

Assay: 
% #VALUE_ID: LEX-P0048-V0001
% #FIELD_VALUE: Assay
\fieldvalue{MSC-1}

Cell Type F1: 
% #VALUE_ID: LEX-P0048-V0002
% #FIELD_VALUE: Cell Type F1
\fieldvalue{MSC (AO)}

Cell Type F2: 
% #VALUE_ID: LEX-P0048-V0003
% #FIELD_VALUE: Cell Type F2
\fieldvalue{MSC (PI)}

Sample ID: 
% #VALUE_ID: LEX-P0048-V0004
% #FIELD_VALUE: Sample ID
\fieldvalue{QC-JS-QY0401-A3122202508008-20250827-27.4H-1}

Dilution: 
% #VALUE_ID: LEX-P0048-V0005
% #FIELD_VALUE: Dilution
\fieldvalue{2.00}

\noindent\rule{\textwidth}{0.4pt}

Results:

\begin{tabularx}{\textwidth}{@{}lXlXlX@{}}
\textbf{Count} & & \textbf{Concentration} & & \textbf{Mean Diameter} & \\
\cline{1-1}\cline{3-3}\cline{5-5}
Total: &
% #VALUE_ID: LEX-P0048-V0006
% #FIELD_VALUE: Total count
\fieldvalue{172}
&
Total: &
% #VALUE_ID: LEX-P0048-V0007
% #FIELD_VALUE: Total concentration
\fieldvalue{$6.06\times10^{5}$ cells/mL}
&
Total: &
% #VALUE_ID: LEX-P0048-V0008
% #FIELD_VALUE: Total mean diameter
\fieldvalue{14.0 micron}
\\
Live: &
% #VALUE_ID: LEX-P0048-V0009
% #FIELD_VALUE: Live count
\fieldvalue{163}
&
Live: &
% #VALUE_ID: LEX-P0048-V0010
% #FIELD_VALUE: Live concentration
\fieldvalue{$5.75\times10^{5}$ cells/mL}
&
Live: &
% #VALUE_ID: LEX-P0048-V0011
% #FIELD_VALUE: Live mean diameter
\fieldvalue{14.1 micron}
\\
Dead: &
% #VALUE_ID: LEX-P0048-V0012
% #FIELD_VALUE: Dead count
\fieldvalue{9}
&
Dead: &
% #VALUE_ID: LEX-P0048-V0013
% #FIELD_VALUE: Dead concentration
\fieldvalue{$3.10\times10^{5}$ cells/mL}
&
Dead: &
% #VALUE_ID: LEX-P0048-V0014
% #FIELD_VALUE: Dead mean diameter
\fieldvalue{11.5 micron}
\\
Viability: &
% #VALUE_ID: LEX-P0048-V0015
% #FIELD_VALUE: Viability
\fieldvalue{94.9\%}
&
&
&
&
\\
\end{tabularx}

\vspace{1em}

\begin{center}
\begin{tabular}{cc}
\fbox{\parbox[c][4.2cm][c]{0.38\textwidth}{\centering
% TODO: Insert the upper-left microscopy image visible on this page.
}} &
\fbox{\parbox[c][4.2cm][c]{0.38\textwidth}{\centering
% TODO: Insert the upper-right microscopy image visible on this page.
}} \\[1em]
\fbox{\parbox[c][4.2cm][c]{0.38\textwidth}{\centering
% TODO: Insert the lower-left microscopy image visible on this page.
}} &
\fbox{\parbox[c][4.2cm][c]{0.38\textwidth}{\centering
% TODO: Insert the lower-right microscopy image visible on this page.
}}
\end{tabular}
\end{center}

\vspace{1em}

操作人/审核人：
% #VALUE_ID: LEX-P0048-V0016
% #FIELD_VALUE: 操作人/审核人
% #TODO #HANDWRITTEN: illegible signature; handwritten name is not clearly discernible
\fieldvalue{\handwritten{illegible signature}}

\hfill 日期：
% #VALUE_ID: LEX-P0048-V0017
% #FIELD_VALUE: 日期
% #HANDWRITTEN: 2024.8.27
\fieldvalue{\handwritten{2024.8.27}}
% LEXOID_PAGE_COMPLETED: 48/60

\newpage

\noindent
Assay:
% #VALUE_ID: LEX-P0049-V0001
% #FIELD_VALUE: Assay
\fieldvalue{MSC-1}

\vspace{0.6em}

\noindent
Cell Type F1:
% #VALUE_ID: LEX-P0049-V0002
% #FIELD_VALUE: Cell Type F1
\fieldvalue{MSC (AO)}

\noindent
Cell Type F2:
% #VALUE_ID: LEX-P0049-V0003
% #FIELD_VALUE: Cell Type F2
\fieldvalue{MSC (PI)}

\vspace{0.6em}

\noindent
Sample ID:
% #VALUE_ID: LEX-P0049-V0004
% #FIELD_VALUE: Sample ID
\fieldvalue{QC-JS-QY0401-A31Z2202508008-20250827-27.4H-2}

\noindent
Dilution:
% #VALUE_ID: LEX-P0049-V0005
% #FIELD_VALUE: Dilution
\fieldvalue{2.00}

\vspace{0.5em}
\noindent\rule{\textwidth}{0.4pt}

\noindent Results:

\vspace{0.4em}

\begin{tabularx}{\textwidth}{@{}X X X@{}}
\textbf{Count} & \textbf{Concentration} & \textbf{Mean Diameter} \\
\texttt{===========} & \texttt{===========} & \texttt{=============} \\[0.2em]
Total:
% #VALUE_ID: LEX-P0049-V0006
% #FIELD_VALUE: Total count
\fieldvalue{126}
&
% #VALUE_ID: LEX-P0049-V0007
% #FIELD_VALUE: Total concentration
\fieldvalue{$4.49\times10^{5}$ cells/mL}
&
% #VALUE_ID: LEX-P0049-V0008
% #FIELD_VALUE: Total mean diameter
\fieldvalue{17.1 micron}
\\
Live:
% #VALUE_ID: LEX-P0049-V0009
% #FIELD_VALUE: Live count
\fieldvalue{114}
&
% #VALUE_ID: LEX-P0049-V0010
% #FIELD_VALUE: Live concentration
\fieldvalue{$4.07\times10^{5}$ cells/mL}
&
% #VALUE_ID: LEX-P0049-V0011
% #FIELD_VALUE: Live mean diameter
\fieldvalue{17.7 microns}
\\
Dead:
% #VALUE_ID: LEX-P0049-V0012
% #FIELD_VALUE: Dead count
\fieldvalue{12}
&
% #VALUE_ID: LEX-P0049-V0013
% #FIELD_VALUE: Dead concentration
\fieldvalue{$4.20\times10^{5}$ cells/mL}
&
% #VALUE_ID: LEX-P0049-V0014
% #FIELD_VALUE: Dead mean diameter
\fieldvalue{11.5 microns}
\\
Viability:
% #VALUE_ID: LEX-P0049-V0015
% #FIELD_VALUE: Viability
\fieldvalue{90.6\%}
&
&
\end{tabularx}

\vspace{1.2em}

\begin{center}
\begin{tabular}{cc}
\fbox{\parbox[c][0.22\textheight][c]{0.40\textwidth}{\centering
% TODO: Insert the upper-left microscopy image visible on this page.
}} &
\fbox{\parbox[c][0.22\textheight][c]{0.40\textwidth}{\centering
% TODO: Insert the upper-right microscopy image visible on this page.
}} \\[1em]
\fbox{\parbox[c][0.22\textheight][c]{0.40\textwidth}{\centering
% TODO: Insert the lower-left microscopy image visible on this page.
}} &
\fbox{\parbox[c][0.22\textheight][c]{0.40\textwidth}{\centering
% TODO: Insert the lower-right microscopy image visible on this page.
}}
\end{tabular}
\end{center}

\vfill

\noindent
操作人/审核人：
% #VALUE_ID: LEX-P0049-V0016
% #FIELD_VALUE: 操作人/审核人
% #TODO #HANDWRITTEN: illegible signature; handwritten blue signature is not clearly readable
\fieldvalue{\handwritten{[illegible signature]}}

\hfill
日期：
% #VALUE_ID: LEX-P0049-V0017
% #FIELD_VALUE: 日期
% #HANDWRITTEN: 2024.8.7
\fieldvalue{\handwritten{2024.8.7}}
% LEXOID_PAGE_COMPLETED: 49/60

\newpage

Assay: MSC-1

Cell Type F1: MSC (AO)\\
Cell Type F2: MSC (PI)

Sample ID: QC-JS-QY0401-A31Z2202508008-20250827-27.4H-3\\
Dilution: 2.00

\noindent\rule{\textwidth}{0.4pt}

Results:

\begin{tabularx}{\textwidth}{XXX}
\textbf{Count} & \textbf{Concentration} & \textbf{Mean Diameter} \\
\(\underline{\hspace{1.5cm}}\) & \(\underline{\hspace{2.0cm}}\) & \(\underline{\hspace{2.0cm}}\) \\
Total: 157 & \(5.55\times10^5\) cells/mL & 15.0 micron \\
Live: 139 & \(4.92\times10^5\) cells/mL & 15.4 microns \\
Dead: 18 & \(6.30\times10^4\) cells/mL & 11.9 microns \\
Viability: 88.6\% & & 
\end{tabularx}

\vspace{1em}

\begin{center}
\begin{tabular}{cc}
% TODO: Image visible on the page; source filename is unavailable.
\fbox{\rule{0pt}{4.0cm}\rule{0.42\textwidth}{0pt}} &
% TODO: Image visible on the page; source filename is unavailable.
\fbox{\rule{0pt}{4.0cm}\rule{0.42\textwidth}{0pt}} \\[1em]
% TODO: Image visible on the page; source filename is unavailable.
\fbox{\rule{0pt}{4.0cm}\rule{0.42\textwidth}{0pt}} &
% TODO: Image visible on the page; source filename is unavailable.
\fbox{\rule{0pt}{4.0cm}\rule{0.42\textwidth}{0pt}}
\end{tabular}
\end{center}

\vfill

操作人/审核人：
% #VALUE_ID: LEX-P0050-V0001
% #FIELD_VALUE: 操作人/审核人
% #TODO #HANDWRITTEN: [illegible signature]; handwritten signature is difficult to decipher
\fieldvalue{\handwritten{[illegible signature]}}

\hfill 日期：
% #VALUE_ID: LEX-P0050-V0002
% #FIELD_VALUE: 日期
% #HANDWRITTEN: 2025.08.27
\fieldvalue{\handwritten{2025.08.27}}
% LEXOID_PAGE_COMPLETED: 50/60

\newpage

\begin{center}
\small
华生生物科技（武汉）有限公司\\
\textit{PLATINUM LIFE EXCELLENCE BIOTECH}\\[2pt]
文件编码：REC-YZ-SMP-QC-01-001-d \hfill 版本号：1.4\\[2pt]
\textbf{\Large 原件请验单}
\end{center}

\begin{tabularx}{\textwidth}{|p{2.5cm}|X|p{2.8cm}|X|}
\hline
\centering 样品名称 &
% #VALUE_ID: LEX-P0051-V0001
% #FIELD_VALUE: 样品名称
% #TODO #HANDWRITTEN: 绝缘胶板及绝缘片; 字迹较模糊
\fieldvalue{\handwritten{绝缘胶板及绝缘片}} &
\centering 样品代码 &
% #VALUE_ID: LEX-P0051-V0002
% #FIELD_VALUE: 样品代码
% #TODO #HANDWRITTEN: QY0401; 首字符及末位字迹略模糊
\fieldvalue{\handwritten{QY0401}}
\\
\hline
\centering 样品批号 &
% #VALUE_ID: LEX-P0051-V0003
% #FIELD_VALUE: 样品批号
% #TODO #HANDWRITTEN: 4712220508023; 数字部分字迹不清
\fieldvalue{\handwritten{4712220508023}} &
\centering 样品规格/数量 &
% #VALUE_ID: LEX-P0051-V0004
% #FIELD_VALUE: 样品规格/数量
% #TODO #HANDWRITTEN: 49?20mL/袋; 首部数字及单位字迹不清
\fieldvalue{\handwritten{49?20mL/袋}}
\\
\hline
\centering 样品类别 &
% #VALUE_ID: LEX-P0051-V0005
% #FIELD_VALUE: 样品类别
\fieldvalue{\ensuremath{\checkmark}\ 生产过程中检验品}
&
\multicolumn{2}{l|}{\quad □ 待包装产品 \qquad □ 成品}\\
\cline{2-4}
& \multicolumn{3}{l|}{\quad □ 其他（\hspace{3cm}）}\\
\hline
\centering 检验类型 &
% #VALUE_ID: LEX-P0051-V0006
% #FIELD_VALUE: 检验类型
\fieldvalue{\ensuremath{\checkmark}\ 初验}
\quad □ 复验 \quad □ 退货 \quad □ 稳定性 \quad □ 留样
& &\\
\hline
\centering 请验目的 &
% #VALUE_ID: LEX-P0051-V0007
% #FIELD_VALUE: 请验目的
\fieldvalue{\ensuremath{\checkmark}\ 全检}
\quad □ 留样 \quad □ 其他（\hspace{3cm}）
& &\\
\hline
\centering 请验部门 &
% #VALUE_ID: LEX-P0051-V0008
% #FIELD_VALUE: 请验部门
\fieldvalue{\ensuremath{\checkmark}\ 生产部}
\quad □ 物料管理部 \quad □ 其他（\hspace{3cm}）
& &\\
\hline
\centering 样品情况 &
% #VALUE_ID: LEX-P0051-V0009
% #FIELD_VALUE: 样品情况
\fieldvalue{\ensuremath{\checkmark}\ 随请验单一起提供}
\quad □ 未随请验单一起提供
& &\\
\hline
\centering 请验人 &
% #VALUE_ID: LEX-P0051-V0010
% #FIELD_VALUE: 请验人
% #TODO #HANDWRITTEN: 袁婷; 姓名笔画较模糊
\fieldvalue{\handwritten{袁婷}}
&
\centering 请验日期 &
% #VALUE_ID: LEX-P0051-V0011
% #FIELD_VALUE: 请验日期
% #HANDWRITTEN: 2025.08.26
\fieldvalue{\handwritten{2025.08.26}}
\\
\hline
\centering 收验人 &
% #VALUE_ID: LEX-P0051-V0012
% #FIELD_VALUE: 收验人
% #TODO #HANDWRITTEN: 杨婷; 姓名笔画较模糊
\fieldvalue{\handwritten{杨婷}}
&
\centering 收验日期 &
% #VALUE_ID: LEX-P0051-V0013
% #FIELD_VALUE: 收验日期
% #HANDWRITTEN: 2025.08.26
\fieldvalue{\handwritten{2025.08.26}}
\\
\hline
\centering 备注 &
% #VALUE_ID: LEX-P0051-V0014
% #FIELD_VALUE: 备注
% #TODO #HANDWRITTEN: -4\%  OK  2L; 首项及字母部分字迹不清
\fieldvalue{\handwritten{-4\% \quad OK \quad 2L}}
& \multicolumn{2}{l|}{}\\
\hline
\end{tabularx}
% LEXOID_PAGE_COMPLETED: 51/60

\newpage

\begin{center}
\small
\begin{tabularx}{\textwidth}{@{}Xr@{}}
\textcolor{cyan}{铂润生物科技（北京）有限公司} & 文件编码：REC-YZ-SMP-QC-01-004-b\quad 版本号：1.4
\end{tabularx}

\vspace{0.4em}
{\Large\textbf{取样记录}}
\end{center}

\renewcommand{\arraystretch}{1.6}
\small
\begin{tabularx}{\textwidth}{|>{\centering\arraybackslash}p{2.2cm}|X|>{\centering\arraybackslash}p{2.2cm}|X|}
\hline
名称 &
% #VALUE_ID: LEX-P0052-V0001
% #FIELD_VALUE: 名称
\fieldvalue{一类生物反应器细胞液} &
代码 &
% #VALUE_ID: LEX-P0052-V0002
% #FIELD_VALUE: 代码
\fieldvalue{QY0401}
\\ \hline

规格 &
% #VALUE_ID: LEX-P0052-V0003
% #FIELD_VALUE: 规格
\fieldvalue{约20ml/袋} &
总数量 &
% #VALUE_ID: LEX-P0052-V0004
% #FIELD_VALUE: 总数量
\fieldvalue{1袋}
\\ \hline

批号 &
% #VALUE_ID: LEX-P0052-V0005
% #FIELD_VALUE: 批号
\fieldvalue{A31Z2202508008} &
进厂批号 &
% #VALUE_ID: LEX-P0052-V0006
% #FIELD_VALUE: 进厂批号
\fieldvalue{/}
\\ \hline

供货单位/厂家 &
\multicolumn{3}{l|}{
% #VALUE_ID: LEX-P0052-V0007
% #FIELD_VALUE: 供货单位/厂家
\fieldvalue{/}}
\\ \hline

取样目的 &
\multicolumn{3}{l|}{
% #VALUE_ID: LEX-P0052-V0008
% #FIELD_VALUE: 取样目的
% #HANDWRITTEN: √ 正常检验
\fieldvalue{\handwritten{√\,正常检验}}\qquad
$\square$ 复检\qquad $\square$ 稳定性
\newline
\hspace{0.2em}$\square$ 退货检验\qquad $\square$ 留样\qquad $\square$ 其他：\underline{\hspace{2.2cm}}}
\\ \hline

取样地点 &
\multicolumn{3}{l|}{
% #VALUE_ID: LEX-P0052-V0009
% #FIELD_VALUE: 取样地点
% #HANDWRITTEN: √ 产品出厂
\fieldvalue{\handwritten{√\,产品出厂}}\qquad
$\square$ 物料库\qquad $\square$ 细胞库\qquad $\square$ 包装间}
\\ \hline

贮存条件 &
\multicolumn{3}{l|}{
% #VALUE_ID: LEX-P0052-V0010
% #FIELD_VALUE: 贮存条件
% #HANDWRITTEN: √ 2--8℃
\fieldvalue{\handwritten{√\,2--8℃}}\qquad
$\square$ -20℃\qquad $\square$ -80℃\qquad $\square$ 常温\qquad
$\square$ 其他：\underline{\hspace{1.5cm}}}
\\ \hline

盛装容器 &
\multicolumn{3}{l|}{
% #VALUE_ID: LEX-P0052-V0011
% #FIELD_VALUE: 盛装容器
% #HANDWRITTEN: √ 保温箱
\fieldvalue{\handwritten{√\,保温箱}}\qquad
$\square$ 泡沫箱\qquad $\square$ 手提袋\qquad
$\square$ 其他：\underline{\hspace{1.8cm}}}
\\ \hline

外观复核 &
\multicolumn{3}{l|}{
% #VALUE_ID: LEX-P0052-V0012
% #FIELD_VALUE: 外观复核
% #HANDWRITTEN: √ 合格
\fieldvalue{\handwritten{√\,合格}}\qquad
$\square$ 不合格}
\\ \hline

取样量 &
% #VALUE_ID: LEX-P0052-V0013
% #FIELD_VALUE: 取样量
\fieldvalue{1袋} &
检验量 &
% #VALUE_ID: LEX-P0052-V0014
% #FIELD_VALUE: 检验量
\fieldvalue{1袋} \\
\cline{3-4}
\multicolumn{2}{c|}{} &
留样量 &
% #VALUE_ID: LEX-P0052-V0015
% #FIELD_VALUE: 留样量
\fieldvalue{/}
\\ \hline

取样人/日期 &
% #VALUE_ID: LEX-P0052-V0016
% #FIELD_VALUE: 取样人
% #TODO #HANDWRITTEN: 签名难以辨认; 蓝色手写签名笔迹不清
\fieldvalue{\handwritten{签名难以辨认}} &
复核人/日期 &
% #VALUE_ID: LEX-P0052-V0018
% #FIELD_VALUE: 复核人
% #TODO #HANDWRITTEN: 签名难以辨认; 蓝色手写签名笔迹不清
\fieldvalue{\handwritten{签名难以辨认}}
\\ \cline{2-2}\cline{4-4}
&
% #VALUE_ID: LEX-P0052-V0017
% #FIELD_VALUE: 取样日期
% #TODO #HANDWRITTEN: 2015.03.26; 手写日期部分笔迹不清
\fieldvalue{\handwritten{2015.03.26}} &
&
% #VALUE_ID: LEX-P0052-V0019
% #FIELD_VALUE: 复核日期
% #HANDWRITTEN: 2015.03.26
\fieldvalue{\handwritten{2015.03.26}}
\\ \hline

备注 &
\multicolumn{3}{l|}{
% #VALUE_ID: LEX-P0052-V0020
% #FIELD_VALUE: 备注
% #TODO #HANDWRITTEN: OK 2/16; 手写内容部分难以辨认
\fieldvalue{\handwritten{OK 2/16}}}
\\ \hline
\end{tabularx}

\vfill
\begin{center}
\scriptsize 第 1 页 共 1 页
\end{center}
% LEXOID_PAGE_COMPLETED: 52/60

\newpage

\begin{center}
\textbf{生物反应器细胞液检验操作记录}
\end{center}

\begin{tabularx}{\textwidth}{|l|X|l|X|}
\hline
\multicolumn{4}{|c|}{葡萄糖浓度检验} \\ \hline
样品名称 &
% #VALUE_ID: LEX-P0053-V0001
% #FIELD_VALUE: 样品名称
\fieldvalue{二级生物反应器细胞液} &
样品规格 &
% #VALUE_ID: LEX-P0053-V0002
% #FIELD_VALUE: 样品规格
\fieldvalue{1 mL/袋} \\ \hline
样品批号 &
% #VALUE_ID: LEX-P0053-V0003
% #FIELD_VALUE: 样品批号
% #TODO #HANDWRITTEN: A?202?; handwriting unclear
\fieldvalue{\handwritten{A?202?}} &
取样日期 &
% #VALUE_ID: LEX-P0053-V0004
% #FIELD_VALUE: 取样日期
% #HANDWRITTEN: 2024年08月26日
\fieldvalue{\handwritten{2024年08月26日}} \\ \hline
\end{tabularx}

\section{仪器设备、试剂耗材}

\subsection{仪器设备}

\begin{tabularx}{\textwidth}{|l|X|c|X|c|}
\hline
名称 & 厂家 & 型号 & 设备编号 & 检查确认 \\ \hline
血糖仪 &
% #VALUE_ID: LEX-P0053-V0005
% #FIELD_VALUE: 血糖仪厂家
\fieldvalue{三诺生物传感股份有限公司} &
% #VALUE_ID: LEX-P0053-V0006
% #FIELD_VALUE: 血糖仪型号
\fieldvalue{安捷} &
% #VALUE_ID: LEX-P0053-V0007
% #FIELD_VALUE: 血糖仪设备编号
% #TODO #HANDWRITTEN: /; handwritten mark is unclear
\fieldvalue{\handwritten{/}} &
% #VALUE_ID: LEX-P0053-V0008
% #FIELD_VALUE: 血糖仪检查确认
\fieldvalue{\(\checkmark\)正常} \\ \hline
移液器 &
% #VALUE_ID: LEX-P0053-V0009
% #FIELD_VALUE: 移液器厂家
\fieldvalue{Eppendorf Corporate} &
% #VALUE_ID: LEX-P0053-V0010
% #FIELD_VALUE: 移液器型号
\fieldvalue{0.1--2.5\,\textmu l} &
% #VALUE_ID: LEX-P0053-V0011
% #FIELD_VALUE: 移液器设备编号
% #TODO #HANDWRITTEN: 137?486; handwriting partially unclear
\fieldvalue{\handwritten{137?486}} &
% #VALUE_ID: LEX-P0053-V0012
% #FIELD_VALUE: 移液器检查确认
\fieldvalue{\(\checkmark\)正常} \\ \hline
\end{tabularx}

\subsection{试剂耗材}

\begin{tabularx}{\textwidth}{|l|X|X|X|}
\hline
名称 & 厂家 & 批号 & 有效期至 \\ \hline
血糖测试条 &
% #VALUE_ID: LEX-P0053-V0013
% #FIELD_VALUE: 血糖测试条厂家
\fieldvalue{三诺生物传感股份有限公司} &
% #VALUE_ID: LEX-P0053-V0014
% #FIELD_VALUE: 血糖测试条批号
% #TODO #HANDWRITTEN: 444-44-01; handwriting unclear
\fieldvalue{\handwritten{444-44-01}} &
% #VALUE_ID: LEX-P0053-V0015
% #FIELD_VALUE: 血糖测试条有效期至
% #TODO #HANDWRITTEN: 2027-1-30; final digit is unclear
\fieldvalue{\handwritten{2027-1-30}} \\ \hline
\end{tabularx}

\section{操作步骤：}

\begin{tabularx}{\textwidth}{|X|X|X|X|}
\hline
\multicolumn{4}{|c|}{操作过程} \\ \hline
取出测试条，取样品自然沉降后的上清液0.6\,\textmu l加入测试条电极反应处，血糖仪自动发出“嘀”提示音，吸入样品完成，血糖仪自动进入倒计时，开始测试； &
\multicolumn{3}{l|}{
取上清液：
% #VALUE_ID: LEX-P0053-V0016
% #FIELD_VALUE: 取上清液体积
% #HANDWRITTEN: 0.6
\fieldvalue{\handwritten{0.6}}\,\textmu l}
\\ \hline
重复以上步骤2次，共计数3次。 &
\multicolumn{3}{l|}{
每个样品分别共计数
% #VALUE_ID: LEX-P0053-V0017
% #FIELD_VALUE: 样品计数次数
% #HANDWRITTEN: 3
\fieldvalue{\handwritten{3}} 次。} \\ \hline
\multicolumn{4}{|c|}{检测结果} \\ \hline
\multicolumn{4}{|c|}{葡萄糖浓度（mmol/L）} \\ \hline
第一次 & 第二次 & 第三次 & 平均值（保留一位小数） \\ \hline
% #VALUE_ID: LEX-P0053-V0018
% #FIELD_VALUE: 第一次检测结果
% #HANDWRITTEN: 1.1
\fieldvalue{\handwritten{1.1}} &
% #VALUE_ID: LEX-P0053-V0019
% #FIELD_VALUE: 第二次检测结果
% #HANDWRITTEN: 2.3
\fieldvalue{\handwritten{2.3}} &
% #VALUE_ID: LEX-P0053-V0020
% #FIELD_VALUE: 第三次检测结果
% #HANDWRITTEN: 2.19
\fieldvalue{\handwritten{2.19}} &
% #VALUE_ID: LEX-P0053-V0021
% #FIELD_VALUE: 平均值
% #HANDWRITTEN: 2.0
\fieldvalue{\handwritten{2.0}} \\ \hline
\end{tabularx}

\begin{tabularx}{\textwidth}{|l|X|}
\hline
备注 &
% #VALUE_ID: LEX-P0053-V0022
% #FIELD_VALUE: 备注
% #TODO #HANDWRITTEN: /; handwritten mark is unclear
\fieldvalue{\handwritten{/}} \\ \hline
实验人/日期 &
% #VALUE_ID: LEX-P0053-V0023
% #FIELD_VALUE: 实验人
% #TODO #HANDWRITTEN: 郭?; handwriting unclear
\fieldvalue{\handwritten{郭?}}\quad
% #VALUE_ID: LEX-P0053-V0024
% #FIELD_VALUE: 实验日期
% #TODO #HANDWRITTEN: 2025.8.6; handwriting partially unclear
\fieldvalue{\handwritten{2025.8.6}} \\ \hline
复核人/日期 &
% #VALUE_ID: LEX-P0053-V0025
% #FIELD_VALUE: 复核人
% #TODO #HANDWRITTEN: 李?; handwriting unclear
\fieldvalue{\handwritten{李?}}\quad
% #VALUE_ID: LEX-P0053-V0026
% #FIELD_VALUE: 复核日期
% #TODO #HANDWRITTEN: 2025.5.26; handwriting partially unclear
\fieldvalue{\handwritten{2025.5.26}} \\ \hline
\end{tabularx}

\begin{center}
第1页 共5页
\end{center}
% LEXOID_PAGE_COMPLETED: 53/60

\newpage

\begin{center}
铂生物\\
\textbf{生物反应器细胞液检验操作记录}
\end{center}

\begin{tabularx}{\textwidth}{|l|X|l|X|}
\hline
密度与活率检验 & \multicolumn{3}{c|}{}\\
\hline
样品名称 &
% #VALUE_ID: LEX-P0054-V0001
% #FIELD_VALUE: 样品名称
\fieldvalue{一级生物反应器细胞泵液} &
样品规格 &
% #VALUE_ID: LEX-P0054-V0002
% #FIELD_VALUE: 样品规格
% #HANDWRITTEN: 10？2
\fieldvalue{\handwritten{10？2}}\ ml/袋\\
\hline
样品批号 &
% #VALUE_ID: LEX-P0054-V0003
% #FIELD_VALUE: 样品批号
% #HANDWRITTEN: 201？？
\fieldvalue{\handwritten{201？？}} &
取样日期 &
% #VALUE_ID: LEX-P0054-V0004
% #FIELD_VALUE: 取样日期
% #HANDWRITTEN: 2026年03月26日
\fieldvalue{\handwritten{2026年03月26日}}\\
\hline
\end{tabularx}

\section{仪器设备、试剂耗材}

\subsection{仪器设备}

\begin{tabularx}{\textwidth}{|l|l|l|l|X|l|}
\hline
名称 & 厂家 & 型号 & 设备编号 & 校准有效期至 & 检查确认\\
\hline
细胞计数仪 & Nexcelom & Cellometer K2 & 1410905 &
% #VALUE_ID: LEX-P0054-V0005
% #FIELD_VALUE: 校准有效期至
% #HANDWRITTEN: 2026.01.20
\fieldvalue{\handwritten{2026.01.20}} &
% #VALUE_ID: LEX-P0054-V0006
% #FIELD_VALUE: 检查确认（正常）
% #HANDWRITTEN: ✓
\fieldvalue{\handwritten{✓}}正常\quad $\square$异常\\
\hline
高速冷冻离心机 & Thermo & ST16R & 1410810 &
% #VALUE_ID: LEX-P0054-V0007
% #FIELD_VALUE: 校准有效期至
% #HANDWRITTEN: 2026.07.07
\fieldvalue{\handwritten{2026.07.07}} &
% #VALUE_ID: LEX-P0054-V0008
% #FIELD_VALUE: 检查确认（正常）
% #HANDWRITTEN: ✓
\fieldvalue{\handwritten{✓}}正常\quad $\square$异常\\
\hline
电热恒温水槽 & 上海一恒 & DK-8AXX & 1411001 &
% #VALUE_ID: LEX-P0054-V0009
% #FIELD_VALUE: 校准有效期至
% #HANDWRITTEN: 2026.07.07
\fieldvalue{\handwritten{2026.07.07}} &
% #VALUE_ID: LEX-P0054-V0010
% #FIELD_VALUE: 检查确认（正常）
% #HANDWRITTEN: ✓
\fieldvalue{\handwritten{✓}}正常\quad $\square$异常\\
\hline
\end{tabularx}

\subsection{试剂耗材}

\begin{tabularx}{\textwidth}{|l|l|X|X|}
\hline
名称 & 厂家 & 批号 & 有效期至\\
\hline
AO/PI 染液 & Nexcelom &
% #VALUE_ID: LEX-P0054-V0011
% #FIELD_VALUE: AO/PI染液批号
% #HANDWRITTEN: 3？4？7
\fieldvalue{\handwritten{3？4？7}} &
% #VALUE_ID: LEX-P0054-V0012
% #FIELD_VALUE: AO/PI染液有效期至
% #HANDWRITTEN: 2026.5.22
\fieldvalue{\handwritten{2026.5.22}}\\
\hline
0.9\%氯化钠注射液 & 石家庄四药 &
% #VALUE_ID: LEX-P0054-V0013
% #FIELD_VALUE: 0.9\%氯化钠注射液批号
% #HANDWRITTEN: 4？510？3？30
\fieldvalue{\handwritten{4？510？3？30}} &
% #VALUE_ID: LEX-P0054-V0014
% #FIELD_VALUE: 0.9\%氯化钠注射液有效期至
% #HANDWRITTEN: 2026.12.06
\fieldvalue{\handwritten{2026.12.06}}\\
\hline
20\%人血白蛋白 & Baxter AG &
% #VALUE_ID: LEX-P0054-V0015
% #FIELD_VALUE: 20\%人血白蛋白批号
% #HANDWRITTEN: 4？5？/？4
\fieldvalue{\handwritten{4？5？/？4}} &
% #VALUE_ID: LEX-P0054-V0016
% #FIELD_VALUE: 20\%人血白蛋白有效期至
% #HANDWRITTEN: 2027.07.18
\fieldvalue{\handwritten{2027.07.18}}\\
\hline
\end{tabularx}

\section{操作步骤}

\begin{tabularx}{\textwidth}{|X|X|}
\hline
\multicolumn{2}{|c|}{溶液准备}\\
\hline
1\%人白蛋白溶液配制：取5 ml 20\%人血白蛋白加入95 ml的0.9\%氯化钠注射液中，混匀即得。 &
\begin{tabular}{l}
20\%人血白蛋白\quad
% #VALUE_ID: LEX-P0054-V0017
% #FIELD_VALUE: 20\%人血白蛋白用量
% #HANDWRITTEN: 1
\fieldvalue{\handwritten{1}}\ ml\\
0.9\%氯化钠注射液\quad
% #VALUE_ID: LEX-P0054-V0018
% #FIELD_VALUE: 0.9\%氯化钠注射液用量
% #HANDWRITTEN: ？9
\fieldvalue{\handwritten{？9}}\ ml
\end{tabular}\\
\hline
名称 & 配制批号 \hfill 有效期至\\
\hline
5$\times$裂解液 &
% #VALUE_ID: LEX-P0054-V0019
% #FIELD_VALUE: 5倍裂解液配制批号
% #HANDWRITTEN: 5025070801
\fieldvalue{\handwritten{5025070801}}
\hfill
% #VALUE_ID: LEX-P0054-V0020
% #FIELD_VALUE: 5倍裂解液有效期至
% #HANDWRITTEN: 2025.10.07
\fieldvalue{\handwritten{2025.10.07}}\\
\hline
\multicolumn{2}{|c|}{操作过程}\\
\hline
按照50 ml离心管的刻度读取样品体积，往取样管中加入对应体积的5$\times$裂解液（每4 ml样品加入1 ml 5$\times$裂解液，按照“只进不舍”原则，如26 ml样品，加入7 ml 5$\times$裂解液），充分混匀后倒入至2.3的样品中，置于37℃电热恒温水槽中裂解20 min后，取出离心管摇匀，使样品充分裂解；再置于37℃电热恒温水槽中裂解10min，得到细胞悬液。 &
样品体积：
% #VALUE_ID: LEX-P0054-V0021
% #FIELD_VALUE: 样品体积
% #HANDWRITTEN: 70
\fieldvalue{\handwritten{70}}\ ml\\
裂解液：
% #VALUE_ID: LEX-P0054-V0022
% #FIELD_VALUE: 裂解液体积
% #HANDWRITTEN: ？0
\fieldvalue{\handwritten{？0}}\ ml\\
电热恒温水槽
% #VALUE_ID: LEX-P0054-V0023
% #FIELD_VALUE: 电热恒温水槽温度
% #HANDWRITTEN: 37
\fieldvalue{\handwritten{37}} ℃，\\
裂解
% #VALUE_ID: LEX-P0054-V0024
% #FIELD_VALUE: 裂解时间
% #HANDWRITTEN: 0.5h--0.5h
\fieldvalue{\handwritten{0.5h--0.5h}}\\
\hline
取出裂解好的细胞悬液300$\times$g离心5分钟，按照读取样品体积加入对应体积的1\%人白蛋白溶液重悬（每1.2 ml样品体积加入0.1 ml 1\%人白蛋白溶液，按照“只进不舍”原则，如26 ml &
% #VALUE_ID: LEX-P0054-V0025
% #FIELD_VALUE: 离心转速
% #HANDWRITTEN: 300
\fieldvalue{\handwritten{300}}\ $\times g$离心
% #VALUE_ID: LEX-P0054-V0026
% #FIELD_VALUE: 离心时间
% #HANDWRITTEN: 5
\fieldvalue{\handwritten{5}}分钟\\
加入1\%人白蛋白溶液总体积：
% #VALUE_ID: LEX-P0054-V0027
% #FIELD_VALUE: 1\%人白蛋白溶液总体积
% #HANDWRITTEN: 3？7
\fieldvalue{\handwritten{3？7}} ml\\
吸取
% #VALUE_ID: LEX-P0054-V0028
% #FIELD_VALUE: 吸取体积
% #HANDWRITTEN: 20
\fieldvalue{\handwritten{20}}\ $\mu$l\\
\hline
\end{tabularx}

\begin{center}
第2页\quad 共5页
\end{center}
% LEXOID_PAGE_COMPLETED: 54/60

\newpage

\noindent
\begin{minipage}{\textwidth}
\small
\textbf{检测记录}

加入 2.2 ml 1\% 人白蛋白消化液，反复吹打至单细胞悬液。与 2\,$\mu$l AO/PI 染料混合均匀。

\begin{flushright}
吸取
% #VALUE_ID: LEX-P0055-V0001
% #FIELD_VALUE: 吸取体积
\fieldvalue{\handwritten{2\,$\mu$l}}
计数。\\
每个样品共计数
% #VALUE_ID: LEX-P0055-V0002
% #FIELD_VALUE: 计数次数
\fieldvalue{\handwritten{2}}
次。
\end{flushright}

振击计数板上下两层保护膜，将细胞悬液混匀，从中吸取 20\,$\mu$l 与 20\,$\mu$l AO/PI 染料混合均匀。

吸取 20\,$\mu$l 细胞悬液与 AO/PI 染料的混合液，注入到计数板的一个检测室中。

将加入混合液的一端插入进样口。

选择 Assay 类型，使用 AO/PI 染料进行双荧光活细胞检测；

在主界面左侧 ID 栏中输入样本批号。

点击主界面左下角 ``Preview'' 预览细胞图片。

调整仪器焦距直至观察到荧光通道下的细胞中心透亮，轮廓清晰。

点击 ``count'' 按钮进行计数并自动拍摄图片。

重复以上步骤 2 次，共计数 3 次。打印数据结果。

数据保存路径：
% #VALUE_ID: LEX-P0055-V0003
% #FIELD_VALUE: 数据保存路径
% #TODO #HANDWRITTEN: 路径内容难以辨认；蓝色手写字迹模糊
\fieldvalue{\handwritten{\text{[路径难以辨认]}}}

\begin{center}
\begin{tabularx}{\textwidth}{|c|X|X|X|X|X|X|}
\hline
\multicolumn{7}{|c|}{检测结果}\\
\hline
检测次数
&
活细胞密度（个细胞/ml）
&
样品密度（个细胞/ml）
&
活细胞密度 CV 值（\%）
&
细胞活率（\%）
&
样品活率（\%）
&
细胞活率 CV 值（\%）
\\
\hline
第一次
&
% #VALUE_ID: LEX-P0055-V0004
% #FIELD_VALUE: 第一次活细胞密度
% #TODO #HANDWRITTEN: 2.71\times10^{?}；指数部分不清晰
\fieldvalue{\handwritten{$2.71\times10^{?}$}}
&
&
&
% #VALUE_ID: LEX-P0055-V0005
% #FIELD_VALUE: 第一次细胞活率
\fieldvalue{\handwritten{91.2}}
&
&
\\
\hline
第二次
&
% #VALUE_ID: LEX-P0055-V0006
% #FIELD_VALUE: 第二次活细胞密度
% #TODO #HANDWRITTEN: 2.88\times10^{?}；指数部分不清晰
\fieldvalue{\handwritten{$2.88\times10^{?}$}}
&
% #VALUE_ID: LEX-P0055-V0007
% #FIELD_VALUE: 样品密度
\fieldvalue{\handwritten{$2.72\times10^{6}$}}
&
% #VALUE_ID: LEX-P0055-V0008
% #FIELD_VALUE: 活细胞密度 CV 值
\fieldvalue{\handwritten{8}}
&
% #VALUE_ID: LEX-P0055-V0009
% #FIELD_VALUE: 第二次细胞活率
\fieldvalue{\handwritten{93.1}}
&
% #VALUE_ID: LEX-P0055-V0010
% #FIELD_VALUE: 样品活率
\fieldvalue{\handwritten{92}}
&
% #VALUE_ID: LEX-P0055-V0011
% #FIELD_VALUE: 细胞活率 CV 值
\fieldvalue{\handwritten{2}}
\\
\hline
第三次
&
% #VALUE_ID: LEX-P0055-V0012
% #FIELD_VALUE: 第三次活细胞密度
% #TODO #HANDWRITTEN: 2.73\times10^{?}；指数部分不清晰
\fieldvalue{\handwritten{$2.73\times10^{?}$}}
&
&
&
% #VALUE_ID: LEX-P0055-V0013
% #FIELD_VALUE: 第三次细胞活率
\fieldvalue{\handwritten{92.11}}
&
&
\\
\hline
\end{tabularx}
\end{center}

\begin{tabularx}{\textwidth}{|l|X|l|X|}
\hline
发酵罐细胞体积（ml）
&
% #VALUE_ID: LEX-P0055-V0014
% #FIELD_VALUE: 发酵罐细胞体积
% #TODO #HANDWRITTEN: 2.1?；末位数字不清晰
\fieldvalue{\handwritten{2.1?}}
&
细胞数量（个细胞/瓶）
&
% #VALUE_ID: LEX-P0055-V0015
% #FIELD_VALUE: 细胞数量
\fieldvalue{\handwritten{$2.7\times10^{7}$}}
\\
\hline
\end{tabularx}

计量要求：样品密度为 3 次活细胞密度结果平均值乘以 1\% 人白蛋白消化液总体积，再除以生物反应器细胞液体积；细胞数量为样品密度乘以发酵罐体积。

\begin{tabularx}{\textwidth}{|l|X|}
\hline
备注
&
% #VALUE_ID: LEX-P0055-V0016
% #FIELD_VALUE: 备注
% #HANDWRITTEN: /
\fieldvalue{\handwritten{/}}
\\
\hline
实验人/日期
&
% #VALUE_ID: LEX-P0055-V0017
% #FIELD_VALUE: 实验人
% #TODO #HANDWRITTEN: 姓名难以辨认；蓝色手写签名模糊
\fieldvalue{\handwritten{\text{[姓名难以辨认]}}}
\quad
% #VALUE_ID: LEX-P0055-V0018
% #FIELD_VALUE: 实验日期
% #HANDWRITTEN: 2025.8.26
\fieldvalue{\handwritten{2025.8.26}}
\\
\hline
复核人/日期
&
% #VALUE_ID: LEX-P0055-V0019
% #FIELD_VALUE: 复核人
% #TODO #HANDWRITTEN: 李宏霞；姓名部分略有不清晰
\fieldvalue{\handwritten{李宏霞}}
\quad
% #VALUE_ID: LEX-P0055-V0020
% #FIELD_VALUE: 复核日期
% #HANDWRITTEN: 2025.8.26
\fieldvalue{\handwritten{2025.8.26}}
\\
\hline
\end{tabularx}

\end{minipage}
% LEXOID_PAGE_COMPLETED: 55/60

\newpage

\begin{center}
\textbf{生物反应器细胞液检测记录}
\end{center}

\noindent
文档编号：REC-YZ-SOP-QC-99-006-a

\begin{center}
\begin{tabularx}{\textwidth}{|c|X|c|X|}
\hline
\multicolumn{4}{|c|}{细胞生长状态检验}\\
\hline
样品名称 &
% #VALUE_ID: LEX-P0056-V0001
% #FIELD_VALUE: 样品名称
\fieldvalue{一级生物反应器细胞液} &
样品规格 & ml/袋\\
\hline
样品批号 & & 取样日期 & 年\quad 月\quad 日\\
\hline
\end{tabularx}
\end{center}

\textbf{1.\quad 仪器设备及试剂}

\textbf{1.1\quad 仪器设备}

\begin{center}
\begin{tabularx}{\textwidth}{|X|X|X|X|X|}
\hline
仪器名称 & 厂家 & 型号 & 设备编号 & 检查确认\\
\hline
% #VALUE_ID: LEX-P0056-V0002
% #FIELD_VALUE: 仪器名称
\fieldvalue{荧光显微镜} &
% #VALUE_ID: LEX-P0056-V0003
% #FIELD_VALUE: 厂家
\fieldvalue{徕卡} &
% #VALUE_ID: LEX-P0056-V0004
% #FIELD_VALUE: 型号
\fieldvalue{DmiL LED} &
% #VALUE_ID: LEX-P0056-V0005
% #FIELD_VALUE: 设备编号
\fieldvalue{1431404} &
$\square$ 正常\quad $\square$ 异常\\
\hline
\end{tabularx}
\end{center}

\textbf{1.2\quad 试剂耗材}

\begin{center}
\begin{tabularx}{\textwidth}{|X|X|X|X|X|}
\hline
名称 & 厂家 & 货号 & 批号 & 有效期至\\
\hline
% #VALUE_ID: LEX-P0056-V0006
% #FIELD_VALUE: 名称
\fieldvalue{细胞活力检测试剂} &
% #VALUE_ID: LEX-P0056-V0007
% #FIELD_VALUE: 厂家
\fieldvalue{江苏凯基生物技术股份有限公司} &
% #VALUE_ID: LEX-P0056-V0008
% #FIELD_VALUE: 货号
\fieldvalue{KGA9501-1000} & &\\
\hline
\end{tabularx}
\end{center}

\textbf{2.\quad 操作步骤}

\begin{center}
\small
\begin{tabularx}{\textwidth}{|X|X|}
\hline
\multicolumn{2}{|c|}{操作流程}\\
\hline
\multicolumn{2}{|c|}{\textbf{操作流程}}\\
\hline
20$\times$的荧光染液配制：取\underline{\hspace{1.2cm}}\,\textmu l 组分A 和\underline{\hspace{1.2cm}}\,\textmu l 组分B混合均匀配成20$\times$的荧光染液。 &
\\
\hline
取混匀后样品200\,\textmu l/孔至96孔板中，取2孔。 &
样品\underline{\hspace{1.2cm}}\,\textmu l/孔\\
& 每个样品取\underline{\hspace{1.2cm}}孔\\
\hline
自然沉降后，吸弃上清，注意不要吸到沉淀。 &
$\square$ 吸弃上清\\
\hline
取20$\times$的荧光染液加入到96孔板中，200\,\textmu l/孔，避光孵育30--60\,min。 &
20$\times$的荧光染液\underline{\hspace{1.2cm}}\,\textmu l/孔\\
& 避光孵育\underline{\hspace{1.2cm}}\\
\hline
染色结束后，吸弃上清，取PBS加入到以上96孔板中，200\,\textmu l/孔，吸弃上清，注意不要吸到沉淀，清洗3次，之后每孔加入100\,\textmu l PBS进行拍照。 &
加入PBS\underline{\hspace{1.2cm}}\,\textmu l/孔\\
& 清洗\underline{\hspace{1.2cm}}次\\
\hline
将荧光显微镜调至4$\times$倍镜，软件中调节荧光相应参数，调节到合适的光源和像距，拍摄4张清晰图像。 &
\\
\hline
进一步调整图像，打印并粘贴在实验记录上。 &
\\
\hline
数据保存路径 & \\
\hline
\multicolumn{2}{|c|}{检测结果}\\
\hline
\end{tabularx}
\end{center}

% TODO: 页面上可见一条蓝色斜线及其附近的蓝色手写批注，文字难以辨认。
% #VALUE_ID: LEX-P0056-V0009
% #HANDWRITTEN: [难以辨认的蓝色手写批注]
\handwritten{[难以辨认的蓝色手写批注]}
% LEXOID_PAGE_COMPLETED: 56/60

\newpage

\begin{center}
\begin{minipage}{0.92\textwidth}
\small
\begin{tabular}{@{}p{0.28\textwidth}p{0.34\textwidth}p{0.30\textwidth}@{}}
\textbf{铂金催化剂科技（北京）有限公司}
&
\begin{center}
\setlength{\fboxsep}{3pt}
\framebox{\begin{minipage}{0.20\textwidth}\centering
{\scriptsize 原件}
\end{minipage}}
\end{center}
&
\begin{flushright}
文件编码：REC-YZ-SOP-QC-99-006-\ldots\\[-2pt]
\setlength{\fboxsep}{2pt}
\framebox{\begin{minipage}{0.19\textwidth}\centering
{\scriptsize 批准生效}\\
{\scriptsize 版本号：%
% #VALUE_ID: LEX-P0057-V0001
% #FIELD_VALUE: 版本号
\fieldvalue{1.4}}
\end{minipage}}
\end{flushright}
\end{tabular}

\vspace{4pt}

\setlength{\fboxsep}{0pt}
\fbox{\begin{minipage}[t][0.82\textheight][t]{\dimexpr\textwidth-2\fboxsep-2\fboxrule\relax}
\vfill

\begin{center}
\rotatebox{51}{\textcolor{blue}{\rule{0.92\textheight}{0.4pt}}}\\[-0.72\textheight]
\rotatebox{51}{%
% #VALUE_ID: LEX-P0057-V0002
% #HANDWRITTEN: 终版-2014.3.10
\handwritten{\textcolor{blue}{终版-2014.3.10}}}
\end{center}

\vfill

\begin{tabular}{|p{0.16\textwidth}|p{0.76\textwidth}|}
\hline
备注 & \\ \hline
检验人/日期 & \\ \hline
复核人/日期 & \\ \hline
\end{tabular}
\end{minipage}}

\vspace{8pt}
\centerline{\scriptsize 第 5 页 共 5 页}
\end{minipage}
\end{center}

% TODO: The diagonal handwritten annotation is partially illegible; transcription rendered as “终版-2014.3.10”.
% LEXOID_PAGE_COMPLETED: 57/60

\newpage

Assay: MSC-1

\vspace{0.8em}

Cell Type F1: MSC (AO)\\
Cell Type F2: MSC (PI)

\vspace{0.8em}

Sample ID: QC-JS-QY0401-A31Z2202508008-20250826-0H-1\\
Dilution: 2.00

\vspace{0.8em}
\hrule
\vspace{0.3em}

Results:

\vspace{0.5em}

\begin{tabular}{lll}
Count & Concentration & Mean Diameter \\
\texttt{===========} & \texttt{===========} & \texttt{=============} \\
Total: 91 & $3.19\times 10^{5}$ cells/mL & 12.4 micron \\
Live: 83 & $2.91\times 10^{5}$ cells/mL & 12.6 microns \\
Dead: 8 & $2.80\times 10^{4}$ cells/mL & 11.0 micron \\
Viability: 91.2\% & & \\
\end{tabular}

\vspace{2em}

\begin{center}
\begin{tabular}{cc}
% TODO: Image visible in the upper-left position; no source filename is available.
\fbox{\rule{0pt}{4.2cm}\rule{7.2cm}{0pt}} &
% TODO: Image visible in the upper-right position; no source filename is available.
\fbox{\rule{0pt}{4.2cm}\rule{7.2cm}{0pt}} \\[0.8em]
% TODO: Image visible in the lower-left position; no source filename is available.
\fbox{\rule{0pt}{4.2cm}\rule{7.2cm}{0pt}} &
% TODO: Image visible in the lower-right position; no source filename is available.
\fbox{\rule{0pt}{4.2cm}\rule{7.2cm}{0pt}}
\end{tabular}
\end{center}

\vfill

\noindent
\begin{tabular}{p{0.48\textwidth}p{0.48\textwidth}}
操作人/审核人：%
% #VALUE_ID: LEX-P0058-V0001
% #FIELD_VALUE: 操作人/审核人
% #TODO #HANDWRITTEN: illegible signature; stylized handwritten signature cannot be reliably transcribed
\fieldvalue{\handwritten{[illegible signature]}}
&
日期：%
% #VALUE_ID: LEX-P0058-V0002
% #FIELD_VALUE: 日期
% #HANDWRITTEN: 2025.8.26
\fieldvalue{\handwritten{2025.8.26}}
\end{tabular}
% LEXOID_PAGE_COMPLETED: 58/60

\newpage

\noindent
Assay:
% #VALUE_ID: LEX-P0059-V0001
% #FIELD_VALUE: Assay
\fieldvalue{MSC-1}

\vspace{0.5em}
\noindent
Cell Type F1:
% #VALUE_ID: LEX-P0059-V0002
% #FIELD_VALUE: Cell Type F1
\fieldvalue{MSC (AO)}

\noindent
Cell Type F2:
% #VALUE_ID: LEX-P0059-V0003
% #FIELD_VALUE: Cell Type F2
\fieldvalue{MSC (PI)}

\vspace{0.8em}
\noindent
Sample ID:
% #VALUE_ID: LEX-P0059-V0004
% #FIELD_VALUE: Sample ID
\fieldvalue{QC-JS-QY0401-A31Z2202508008-20250826-0H-2}

\noindent
Dilution:
% #VALUE_ID: LEX-P0059-V0005
% #FIELD_VALUE: Dilution
\fieldvalue{2.00}

\vspace{0.8em}
\noindent\rule{\textwidth}{0.4pt}

\noindent
Results:

\vspace{0.5em}
\renewcommand{\arraystretch}{1.15}
\begin{tabularx}{\textwidth}{@{}X X X@{}}
\textbf{Count} & \textbf{Concentration} & \textbf{Mean Diameter} \\
\hline
\texttt{=========} & \texttt{=============} & \texttt{=============} \\[0.2em]
Total:
% #VALUE_ID: LEX-P0059-V0006
% #FIELD_VALUE: Total
\fieldvalue{88}
&
% #VALUE_ID: LEX-P0059-V0010
% #FIELD_VALUE: Total concentration
\fieldvalue{3.09x10$^{5}$ cells/mL}
&
% #VALUE_ID: LEX-P0059-V0013
% #FIELD_VALUE: Total mean diameter
\fieldvalue{13.2 micron}
\\
Live:
% #VALUE_ID: LEX-P0059-V0007
% #FIELD_VALUE: Live
\fieldvalue{82}
&
% #VALUE_ID: LEX-P0059-V0011
% #FIELD_VALUE: Live concentration
\fieldvalue{2.88x10$^{5}$ cells/mL}
&
% #VALUE_ID: LEX-P0059-V0014
% #FIELD_VALUE: Live mean diameter
\fieldvalue{13.3 microns}
\\
Dead:
% #VALUE_ID: LEX-P0059-V0008
% #FIELD_VALUE: Dead
\fieldvalue{6}
&
% #VALUE_ID: LEX-P0059-V0012
% #FIELD_VALUE: Dead concentration
\fieldvalue{2.10x10$^{4}$ cells/mL}
&
% #VALUE_ID: LEX-P0059-V0015
% #FIELD_VALUE: Dead mean diameter
\fieldvalue{12.1 microns}
\\
Viability:
% #VALUE_ID: LEX-P0059-V0009
% #FIELD_VALUE: Viability
\fieldvalue{93.2\%}
& &
\end{tabularx}

\vspace{1.2em}

\noindent
\begin{minipage}[t]{0.48\textwidth}
\centering
% TODO: Insert the visible microscopy image at the upper right of this page; no source filename is available.
\fbox{\rule{0pt}{0.25\textwidth}\rule{0.9\linewidth}{0pt}}
\end{minipage}
\hfill
\begin{minipage}[t]{0.48\textwidth}
\centering
% TODO: The upper-left region contains a very faint/overexposed microscopy image; no source filename is available.
\fbox{\rule{0pt}{0.25\textwidth}\rule{0.9\linewidth}{0pt}}
\end{minipage}

\vspace{1em}

\noindent
\begin{minipage}[t]{0.48\textwidth}
\centering
% TODO: Insert the visible microscopy image at the lower left of this page; no source filename is available.
\fbox{\rule{0pt}{0.28\textwidth}\rule{0.9\linewidth}{0pt}}
\end{minipage}
\hfill
\begin{minipage}[t]{0.48\textwidth}
\centering
% TODO: Insert the visible microscopy image at the lower right of this page; no source filename is available.
\fbox{\rule{0pt}{0.28\textwidth}\rule{0.9\linewidth}{0pt}}
\end{minipage}

\vfill

\noindent
操作人/审核人：
% #VALUE_ID: LEX-P0059-V0016
% #FIELD_VALUE: 操作人/审核人
% #TODO #HANDWRITTEN: illegible signature; handwritten signature
\fieldvalue{\handwritten{[illegible signature]}}

\hfill
日期：
% #VALUE_ID: LEX-P0059-V0017
% #FIELD_VALUE: 日期
% #HANDWRITTEN: 24.8.26
\fieldvalue{\handwritten{24.8.26}}
% LEXOID_PAGE_COMPLETED: 59/60

\newpage

\noindent\textbf{Assay: }
% #VALUE_ID: LEX-P0060-V0001
% #FIELD_VALUE: Assay
\fieldvalue{MSC-1}

\vspace{0.8em}

\noindent\textbf{Cell Type F1: }
% #VALUE_ID: LEX-P0060-V0002
% #FIELD_VALUE: Cell Type F1
\fieldvalue{MSC (AO)}

\noindent\textbf{Cell Type F2: }
% #VALUE_ID: LEX-P0060-V0003
% #FIELD_VALUE: Cell Type F2
\fieldvalue{MSC (PI)}

\vspace{0.8em}

\noindent\textbf{Sample ID: }
% #VALUE_ID: LEX-P0060-V0004
% #FIELD_VALUE: Sample ID
\fieldvalue{QC-JS-QY0401-A31Z2202508008-20250826-OH-3}

\noindent\textbf{Dilution: }
% #VALUE_ID: LEX-P0060-V0005
% #FIELD_VALUE: Dilution
\fieldvalue{2.00}

\vspace{0.6em}
\noindent\rule{\textwidth}{0.4pt}

\noindent Results:

\vspace{0.5em}

\renewcommand{\arraystretch}{1.15}
\begin{tabularx}{\textwidth}{@{}X X X@{}}
\textbf{Count} & \textbf{Concentration} & \textbf{Mean Diameter} \\
\texttt{=======} & \texttt{================} & \texttt{================} \\

Total:
% #VALUE_ID: LEX-P0060-V0006
% #FIELD_VALUE: Count -- Total
\fieldvalue{101}
&
% #VALUE_ID: LEX-P0060-V0007
% #FIELD_VALUE: Concentration -- Total
\fieldvalue{$3.55{\times}10^{5}$ cells/mL}
&
% #VALUE_ID: LEX-P0060-V0008
% #FIELD_VALUE: Mean Diameter -- Total
\fieldvalue{13.0 micron}
\\

Live:
% #VALUE_ID: LEX-P0060-V0009
% #FIELD_VALUE: Count -- Live
\fieldvalue{93}
&
% #VALUE_ID: LEX-P0060-V0010
% #FIELD_VALUE: Concentration -- Live
\fieldvalue{$3.27{\times}10^{5}$ cells/mL}
&
% #VALUE_ID: LEX-P0060-V0011
% #FIELD_VALUE: Mean Diameter -- Live
\fieldvalue{13.0 microns}
\\

Dead:
% #VALUE_ID: LEX-P0060-V0012
% #FIELD_VALUE: Count -- Dead
\fieldvalue{8}
&
% #VALUE_ID: LEX-P0060-V0013
% #FIELD_VALUE: Concentration -- Dead
\fieldvalue{$2.80{\times}10^{4}$ cells/mL}
&
% #VALUE_ID: LEX-P0060-V0014
% #FIELD_VALUE: Mean Diameter -- Dead
\fieldvalue{13.4 microns}
\\

Viability:
% #VALUE_ID: LEX-P0060-V0015
% #FIELD_VALUE: Viability
\fieldvalue{92.1\%}
&
&
\end{tabularx}

\vspace{1.5em}

% TODO: Three assay images are visible on this page; source filenames are unavailable.
\begin{center}
\setlength{\fboxsep}{0pt}
\begin{minipage}[t]{0.42\textwidth}
\centering
\fbox{\rule{0pt}{3.0cm}\rule{0.95\linewidth}{0pt}}\\[-0.2em]
\small Assay image
\end{minipage}
\hfill
\begin{minipage}[t]{0.42\textwidth}
\centering
\fbox{\rule{0pt}{3.0cm}\rule{0.95\linewidth}{0pt}}\\[-0.2em]
\small Assay image
\end{minipage}

\vspace{1em}

\begin{minipage}[t]{0.42\textwidth}
\centering
\fbox{\rule{0pt}{3.0cm}\rule{0.95\linewidth}{0pt}}\\[-0.2em]
\small Assay image
\end{minipage}
\end{center}

\vfill

\noindent 操作人/审核人：

% #VALUE_ID: LEX-P0060-V0016
% #HANDWRITTEN: [签名，字迹难以辨认]
\handwritten{[签名]}

\hfill 日期：

% #VALUE_ID: LEX-P0060-V0017
% #FIELD_VALUE: 日期
% #HANDWRITTEN: 2025.8.26
\fieldvalue{\handwritten{2025.8.26}}

\end{document}
% LEXOID_PAGE_COMPLETED: 60/60
